\documentclass[11pt,a4paper]{article}
\usepackage[utf8]{inputenc}
\usepackage[T1]{fontenc}
\usepackage[czech]{babel}
\usepackage[margin=2.2cm]{geometry}
\usepackage{amsmath,amssymb,amsthm}
\usepackage{parskip}
\usepackage{enumitem}
\setlist{nosep,leftmargin=1.4em}
\usepackage{booktabs}
\usepackage{array}
\usepackage{xcolor}
\definecolor{linkblue}{RGB}{20,60,150}
\definecolor{defbg}{RGB}{240,244,252}
\definecolor{defframe}{RGB}{100,130,190}
\definecolor{markcol}{RGB}{176,84,0}
\definecolor{markbg}{RGB}{253,246,236}
\usepackage[most]{tcolorbox}
\newtcolorbox{defbox}[1][]{breakable,colback=defbg,colframe=defframe,boxrule=0.6pt,arc=1.5mm,left=2mm,right=2mm,top=1mm,bottom=1mm,title={#1},fonttitle=\bfseries}
\newtcolorbox{poznbox}[1][]{breakable,colback=markbg,colframe=markcol,boxrule=0.6pt,arc=1.5mm,left=2mm,right=2mm,top=1mm,bottom=1mm,title={#1},fonttitle=\bfseries}
\newtheorem{veta}{Věta}[section]
\theoremstyle{definition}
\newtheorem{definice}{Definice}[section]
\usepackage[colorlinks=true,linkcolor=linkblue,urlcolor=linkblue]{hyperref}
\usepackage{algorithm}
\usepackage{algpseudocode}

\floatname{algorithm}{Algoritmus}
\renewcommand{\algorithmicrequire}{\textbf{Vstup:}}
\renewcommand{\algorithmicensure}{\textbf{Výstup:}}

% Značka pro látku, která není ve slidech EVA I / EVA II.
\newcommand{\mimo}{\textcolor{markcol}{\textsf{\footnotesize[$\blacklozenge$~nad rámec slidů]}}}
\newcommand{\mimoq}[1]{\textcolor{markcol}{\textsf{\footnotesize[$\blacklozenge$~nad rámec slidů: #1]}}}

\title{\textbf{Přírodou inspirované počítání}\\[2mm]
\large Učební text k~okruhu státní závěrečné zkoušky}
\author{SZZ Umělá inteligence, MFF UK}
\date{srpen 2026}

\begin{document}
\maketitle

\begin{center}
\begin{minipage}{0.92\textwidth}
\small
Tento text pokrývá okruh \emph{Přírodou inspirované počítání} magisterské SZZ oboru Umělá
inteligence. Jeho členění se drží oficiálního znění okruhu: každé jeho větě odpovídá
\textbf{jedna ku jedné} právě jedna kapitola, takže je vždycky vidět, kterou část zadání
se čtenář zrovna učí. Přehled té vazby dává mapovací tabulka níže.

Výkladovým základem jsou slidy přednášek NAIL025 (Evoluční algoritmy~I) a~NAIL086
(Evoluční algoritmy~II). Tam, kde slidy látku jen naznačují, ji text doplňuje o~standardní
podání z~učebnic Eiben \& Smith: \emph{Introduction to Evolutionary Computing}, Mitchell:
\emph{An Introduction to Genetic Algorithms}, Michalewicz: \emph{Genetic Algorithms $+$
Data Structures $=$ Evolution Programs} a~Poli, Langdon \& McPhee: \emph{A Field Guide to
Genetic Programming}. Vedle definic a~pseudokódů obsahuje také odvození, číselné příklady
a~srovnávací tabulky. Každá kapitola pak končí shrnutím, které je psané rovnou pro ústní
zkoušku.
\end{minipage}
\end{center}

\begin{center}
\begin{minipage}{0.92\textwidth}
\begin{poznbox}[Rozsah textu a~značení]
\small
Kapitoly nejsou stejně dlouhé, a~je to záměr: jejich rozsah odpovídá váze látky. Nejvíce
prostoru proto dostávají genetické algoritmy spolu s~genetickým programováním, teorie
schémat a~aplikace, tedy právě to, čemu se přednášky věnují nejpodrobněji.

Značka \mimo{} upozorňuje na látku, která \textbf{není v~žádných dostupných slidech},
tj.\ ani v~EVA~I, ani v~EVA~II, ani ve volitelné \emph{Evoluční robotice} (Mráz).
V~textu je taková látka přesto ponechána, a~to ze dvou důvodů: buď ji \emph{jmenuje
oficiální znění okruhu}, nebo by bez ní okolní výklad nedával smysl. Značka tak slouží jako
nabídka ke~škrtnutí: co označeno není, je ve slidech doloženo. Je-li látka ve slidech
zmíněna jen letmo, uvádí značka rovnou i~rozsah té zmínky.

Nejnápadnější případ je \textbf{otevřená evoluce}. Okruh ji výslovně jmenuje, slidy o~ní
mají jedinou větu. Podobně jsou na tom mravenčí algoritmy (ACO): okruh mluví o~rojových
algoritmech v~množném čísle, slidy ale probírají jen PSO. (Rojová robotika, kterou probírá
Evoluční robotika, je něco jiného než rojová optimalizace.) Opačným případem jsou
\emph{gramatická evoluce} a~\emph{interaktivní evoluce}. Ty označeny nejsou: v~EVA sice
chybějí, ale Evoluční robotika je probírá.
\end{poznbox}
\end{minipage}
\end{center}

\vspace{2mm}

\section*{Mapování okruhu na kapitoly}
\addcontentsline{toc}{section}{Mapování okruhu na kapitoly}

\begin{center}
\small
\begin{tabular}{@{}>{\raggedright\arraybackslash}p{8.6cm}>{\raggedright\arraybackslash}p{7.4cm}@{}}
\toprule
\textbf{Věta oficiálního znění okruhu} & \textbf{Kapitola v~tomto textu} \\
\midrule
Genetické algoritmy, genetické a~evoluční programování.
  & \ref{sec:ga}~Genetické algoritmy, genetické a~evoluční programování \\
Teorie schémat, pravděpodobnostní modely jednoduchého genetického algoritmu.
  & \ref{sec:schemata}~Teorie schémat a~pravděpodobnostní modely SGA \\
Evoluční strategie, diferenciální evoluce, koevoluce, otevřená evoluce.
  & \ref{sec:es}~Evoluční strategie, diferenciální evoluce, koevoluce, otevřená evoluce \\
Rojové optimalizační algoritmy.
  & \ref{sec:roje}~Rojové optimalizační algoritmy \\
Memetické algoritmy, hill climbing, simulované žíhání.
  & \ref{sec:lokalni}~Lokální prohledávání a~memetické algoritmy \\
Aplikace evolučních algoritmů (evoluce expertních systémů, neuroevoluce, řešení
kombinatorických úloh, vícekriteriální optimalizace).
  & \ref{sec:aplikace}~Aplikace evolučních algoritmů \\
\bottomrule
\end{tabular}
\end{center}

\vspace{2mm}
\tableofcontents
\newpage

%=====================================================================
\section{Genetické algoritmy, genetické a~evoluční programování}
\label{sec:ga}
%=====================================================================

Tohle je nejobsáhlejší věta celého okruhu a~tomu odpovídá i~stavba kapitoly. Nejdřív přijde
obecné schéma evolučního algoritmu, protože bez něj nedává zbytek textu smysl. Na ně navazuje
Hollandův jednoduchý genetický algoritmus spolu s~reprezentacemi jedince. Kapitola končí
genetickým a~evolučním programováním, tedy evolucí spustitelných struktur.

\subsection{Obecné schéma evolučního algoritmu}
\label{ssec:obecne}
\subsubsection{Zařazení oboru}

Pod hlavičkou \emph{přírodou inspirovaného počítání} (\emph{nature-inspired computing}) se
schází metaheuristické algoritmy, které principy svého fungování přebírají z~toho, co se
pozoruje v~živé, a~někdy i~v~neživé přírodě. Obor se tradičně člení na tři postupně se
zužující okruhy:

\begin{itemize}
\item \textbf{Výpočetní inteligence} (\emph{computational intelligence}, dříve
  \emph{soft computing}) je z~nich nejširší a~patří do ní neuronové sítě, fuzzy
  systémy a~evoluční výpočty. Vymezuje se proti \uv{tvrdé}, logicky a~symbolicky
  orientované AI.
\item \textbf{Evoluční výpočty} (\emph{evolutionary computation}, EC) jsou
  nadmnožinou evolučních algoritmů: vedle nich sem spadají i~metody, které
  evoluční v~úzkém smyslu nejsou, tedy rojové algoritmy, memetické algoritmy nebo
  tabu prohledávání.
\item \textbf{Evoluční algoritmy} (\emph{evolutionary algorithms}, EA) jsou
  podmnožinou EC, která se drží základních principů přírodní evoluce. Pracují
  s~populací kandidátních řešení, mění ji stochastickými operátory (křížením
  a~mutací) a~o~přežití rozhoduje selekce řízená účelovou funkcí.
\end{itemize}

Z~hlediska klasifikace optimalizačních metod jsou EA \textbf{populační stochastické
prohledávací algoritmy}. Sáhne se po nich tam, kde klasické numerické postupy selhávají.
Typicky nemáme k~dispozici gradient ani analytický tvar účelové funkce a~ta se chová jako
černá skříňka (tzv.\ \emph{black-box optimalizace}). Jindy je prostor řešení diskrétní nebo
strukturovaný a~skládá se ze stromů, grafů či permutací. A~konečně se EA hodí tam, kde
nehledáme jedno řešení, ale rovnou několik vzájemně nedominovaných.

\subsubsection{Biologická inspirace}

Terminologie i~mechanika evolučních algoritmů jsou převzaté z~biologie a~vyplatí se vědět,
odkud která myšlenka pochází.

\begin{itemize}
\item \textbf{Darwin (1859, O~původu druhů).} Prostředí nabízí jen omezené zdroje,
  takže o~ně jedinci soupeří a~klíčem k~přežití je reprodukce. Kdo je lépe
  přizpůsobený, má větší šanci se rozmnožit, a~úspěšné znaky fenotypu se proto
  reprodukují, modifikují a~rekombinují.
\item \textbf{Mendel (1866).} Zavádí gen jako základní jednotku dědičnosti
  a~předpokládá, že se alely předávají nezávisle. Realita je složitější:
  \emph{polygenie} znamená, že jeden znak ovlivňuje více genů, \emph{pleiotropie}
  naopak to, že jeden gen ovlivňuje více znaků, a~mimo Mendelovo schéma stojí
  i~mitochondriální DNA a~epigenetika.
\item \textbf{Watson a~Crick (1953).} Popis struktury DNA přidává
  molekulárně-biologický pohled na to, jak se informace ukládá a~dědí. Kodon,
  tedy trojice nukleotidů, kóduje jednu z~aminokyselin a~kód je redundantní.
\item \textbf{Genotyp $\to$ fenotyp.} Cesta od vnitřního zápisu k~vnějším vlastnostem
  je jednosměrná a~složitá, protože vede přes transkripci DNA${\to}$RNA a~translaci
  RNA${\to}$protein. \emph{Lamarckismus} oproti tomu předpokládá existenci
  inverzního zobrazení, tedy dědičnost získaných vlastností. Biologie ho odmítla,
  v~memetických algoritmech se však jako technika používá (kap.~\ref{sec:lokalni}).
\end{itemize}

Překlad mezi obojím slovníkem je přímočarý a~drží se ho celý zbytek textu:

\begin{defbox}[Slovník: příroda vs.\ algoritmus]
\begin{tabular}{@{}ll@{}}
prostředí & optimalizační úloha \\
jedinec / chromozom & kandidátní řešení, bod prohledávaného prostoru \\
gen, alela & složka řešení a~její hodnota \\
genotyp & vnitřní reprezentace (kódování) jedince \\
fenotyp & \uv{vnějšek} jedince, na němž se hodnotí kvalita \\
fitness & míra kvality řešení (účelová funkce) \\
populace & multimnožina kandidátních řešení \\
generace & jedna iterace algoritmu \\
\end{tabular}
\end{defbox}

\subsubsection{Obecné schéma evolučního algoritmu}

Všechny evoluční algoritmy sdílejí jeden a~týž cyklus a~liší se až tím, čím jeho jednotlivé
kroky naplní.

\begin{defbox}[Evoluční algoritmus]
Evoluční algoritmus je iterativní stochastický proces, který udržuje populaci
$P(t)$ kandidátních řešení a~z~ní opakovaně vytváří populaci $P(t+1)$ pomocí
\emph{rodičovské selekce}, \emph{rekombinace}, \emph{mutace} a~\emph{environmentální
selekce}, dokud není splněna ukončovací podmínka.
\end{defbox}

\begin{algorithm}[!ht]
\caption{Obecné schéma evolučního algoritmu}
\begin{algorithmic}[1]
\Require reprezentace, fitness $f$, operátory, parametry
\State $t \gets 0$; inicializuj náhodně populaci $P(0)$
\State ohodnoť všechny jedince v~$P(0)$ funkcí $f$
\While{není splněna ukončovací podmínka}
  \State \textbf{rodičovská selekce:} vyber z~$P(t)$ množinu rodičů (pravděpodobnostně, dle fitness)
  \State \textbf{rekombinace:} z~dvojic (n-tic) rodičů vytvoř potomky s~pravděpodobností $p_C$
  \State \textbf{mutace:} každého potomka náhodně poruš s~pravděpodobností $p_M$
  \State ohodnoť potomky $P'(t+1)$
  \State \textbf{environmentální selekce:} z~$P(t) \cup P'(t+1)$ (nebo jen z~$P'(t+1)$) vyber novou populaci $P(t+1)$
  \State $t \gets t+1$
\EndWhile
\Ensure nejlepší nalezený jedinec
\end{algorithmic}
\end{algorithm}

Uvedený cyklus je společný všem EA. Jednotlivé varianty, tedy GA, ES, EP, GP nebo DE, se liší
až volbou reprezentace, důrazem na jednotlivé operátory a~typem selekce. Uvnitř cyklu přitom
proti sobě stojí dvě síly:

\begin{itemize}
\item \textbf{Variace}, tedy křížení a~mutace, vytváří novou rozmanitost, a~stará
  se tak o~\emph{exploraci} (\emph{exploration}) prostoru.
\item \textbf{Selekce} naopak rozmanitost odebírá a~soustředí hledání do slibných
  oblastí, čímž zajišťuje \emph{exploataci} (\emph{exploitation}).
\end{itemize}

Najít mezi nimi rovnováhu je ústřední problém návrhu každého EA. Když převáží explorace,
algoritmus degeneruje na náhodnou procházku. Když převáží exploatace, populace uvízne
v~lokálním extrému a~dojde k~\emph{předčasné konvergenci} (\emph{premature convergence}).

\subsubsection{Složky evolučního algoritmu}

Návrh konkrétního algoritmu se skládá z~několika rozhodnutí, která na sebe navazují.

\paragraph{Reprezentace.} Volba kódování je vůbec nejdůležitější návrhové rozhodnutí.
Nabízí se binární řetězec, vektor celých čísel, vektor reálných čísel, permutace,
strom, graf, matice, konečný automat i~neuronová síť. Volba není kosmetická, protože
reprezentace určuje, jaké operátory vůbec dávají smysl.

\paragraph{Fitness funkce.} Kvalitu jedince měří zobrazení $f:$ prostor fenotypů
$\to \mathbb{R}$. Obvykle se maximalizuje, takže minimalizační úlohu je potřeba převrátit,
například pomocí $f' = C_{\max} - f$ nebo $f' = 1/(1+f)$. Sama fitness nemusí být pevná
hodnota: bývá \emph{deterministická}, \emph{zašuměná}, počítá-li se simulací,
\emph{dynamická}, mění-li se v~čase (odst.~\ref{sec:koevoluce}), nebo \emph{kontextová},
závisí-li na ostatních jedincích, jak je tomu u~koevoluce.

\paragraph{Populace.} Populace je multimnožina jedinců, obvykle pevné velikosti $n$.
Právě ona je nositelem \emph{diverzity} a~tu lze měřit průměrnou Hammingovou vzdáleností,
entropií alel na jednotlivých pozicích nebo rozptylem fitness.

\paragraph{Inicializace.} Standardem je náhodná uniformní inicializace, existují ale
i~cílenější postupy. \emph{Seeding} vkládá do počáteční populace heuristická nebo
expertní řešení, čímž ovšem riskuje předčasnou konvergenci. Dále se používají konstrukční
heuristiky, například nejbližší soused pro TSP, a~pro reálné vektory latinské hyperkrychle.

\paragraph{Ukončovací podmínka.} Běh se zastaví po vyčerpání rozpočtu, tedy maximálního
počtu generací či vyhodnocení fitness, nebo po dosažení cílové kvality. Dalšími rozumnými
kritérii jsou stagnace nejlepší fitness po $k$ generací a~ztráta diverzity populace.

\subsubsection{Selekční mechanismy}\label{ssec:selekce}

V~selekci se rozhoduje o~velikosti selekčního tlaku, a~proto se jí věnuje celá řada
mechanismů. Liší se tím, kolik informace o~fitness využívají a~jak citlivé jsou na její
měřítko.

\paragraph{Ruletová (fitness-proporcionální) selekce}

\begin{defbox}[Ruletová selekce (\emph{roulette wheel selection})]
Jedinec $x_i$ je vybrán s~pravděpodobností
\[
p_i \;=\; \frac{f(x_i)}{\sum_{j=1}^{n} f(x_j)} \;=\; \frac{f(x_i)}{n\,\bar{f}},
\]
kde $\bar f$ je průměrná fitness populace. Očekávaný počet výběrů jedince při
$n$ tazích je $n p_i = f(x_i)/\bar f$. Vybírá se \emph{s~vracením}, takže jeden
jedinec může uspět i~vícekrát.
\end{defbox}

\textbf{Příklad (klasický Goldbergův).}\label{ex:goldberg} Uvažme populaci čtyř
5-bitových řetězců, které se dekódují jako celá čísla a~jejichž fitness je $f(x) = x^2$:

\begin{center}
\begin{tabular}{@{}llrrrr@{}}
\toprule
$i$ & řetězec & $x$ & $f=x^2$ & $p_i = f/\sum f$ & oček.\ počet $f/\bar f$ \\
\midrule
1 & 01101 & 13 & 169 & 0{,}144 & 0{,}58 \\
2 & 11000 & 24 & 576 & 0{,}492 & 1{,}97 \\
3 & 01000 &  8 &  64 & 0{,}055 & 0{,}22 \\
4 & 10011 & 19 & 361 & 0{,}309 & 1{,}23 \\
\midrule
  & \multicolumn{2}{l}{součet} & 1170 & 1{,}000 & 4{,}00 \\
  & \multicolumn{2}{l}{průměr $\bar f$} & 292{,}5 & & \\
\bottomrule
\end{tabular}
\end{center}

Poslední sloupec je poměr fitness jedince k~průměru populace, tedy očekávaný počet kopií
v~mezipopulaci. Ruleta rozdělí obvod kola na výseče o~velikostech $p_i$ a~pak se $n$-krát
roztočí. Jedinec 2 tak obsadí v~mezipopulaci zhruba dvě místa, kdežto jedinec 3 z~ní nejspíš
vypadne.

\textbf{Nevýhody ruletové selekce:}
\begin{itemize}
\item \emph{Předčasná dominance} \uv{supermana}: jedinec s~extrémní fitness
  ovládne populaci a~diverzita zmizí.
\item \emph{Ztráta selekčního tlaku} na konci běhu: jakmile jsou všechny fitness
  podobné, jsou podobné i~pravděpodobnosti výběru a~selekce se blíží náhodnému výběru.
\item Ruleta vyžaduje kladné fitness a~je citlivá na afinní transformace, protože
  $f$ a~$f + c$ dávají jiné pravděpodobnosti. Obojí se řeší \emph{škálováním}
  (odst.~\ref{ssec:skalovani}).
\item Výběr má vysoký rozptyl, takže se skutečné počty kopií mohou od očekávaných
  značně lišit.
\end{itemize}

\paragraph{Stochastické univerzální vzorkování (SUS)}

Poslední zmíněnou nevýhodu odstraňuje \emph{stochastic universal sampling}. Kolo se roztočí
jen jednou, zato je po jeho obvodu rozmístěno $n$ ukazatelů vzdálených od sebe o~$1/n$.
Očekávané počty výběrů zůstávají stejné jako u~rulety, jenže rozptyl klesne na minimum:
jedinec s~očekávaným počtem $e_i$ je vybrán $\lfloor e_i \rfloor$ nebo
$\lceil e_i \rceil$-krát. Je to tedy \uv{spravedlivější ruleta}.

\paragraph{Turnajová selekce}

\begin{defbox}[$k$-turnajová selekce (\emph{tournament selection})]
Vyber náhodně (uniformně, s~vracením) $k$ jedinců a~vrať toho nejlepšího.
Parametr $k$ (typicky 2, 3, 5) přímo řídí selekční tlak: pravděpodobnost, že
bude vybrán $i$-tý nejlepší jedinec z~$n$ (\emph{zde} $i=1$ značí
\textbf{nejlepšího}), je
$\big((n-i+1)^k - (n-i)^k\big)/n^k$ --- je to pravděpodobnost, že všech $k$
losovaných padne mezi $i$-tého a~horší, mínus pravděpodobnost, že všech $k$
padne ostře pod $i$-tého. Deterministický turnaj lze změkčit:
lepší vyhraje s~pravděpodobností $p < 1$.
\end{defbox}

Turnaj má hned několik výhod. Nepotřebuje globální informaci o~populaci, protože nic
nesčítá, a~snadno se proto paralelizuje. Je invariantní vůči monotónním transformacím
fitness. A~funguje i~tam, kde fitness není explicitně dána a~porovnání dvou jedinců se
získává až \emph{souhrou}, tedy odehranou partií, simulací nebo koevolucí.

\paragraph{Pořadová selekce}

Ještě dál jde \emph{rank-based selection}, která hodnotu fitness zahodí zcela a~nahradí ji
pořadím jedince. Lineární pořadová selekce přiřadí jedinci s~pořadím $i$ (\emph{pozor, zde
opačně než u~turnaje:} nejhorší $i=1$, nejlepší $i=n$)
pravděpodobnost
\[
p_i \;=\; \frac{1}{n}\Big(s^{-} + (s^{+}-s^{-})\,\frac{i-1}{n-1}\Big),
\qquad 1 \le s^{+} \le 2,\quad s^{-} = 2 - s^{+},
\]
kde $s^{+}$ je očekávaný počet kopií nejlepšího jedince. Exponenciální varianta místo toho
volí $p_i \propto (1-c)c^{\,n-i}$. Protože pořadí na absolutních hodnotách nezávisí, udržuje
pořadová selekce konstantní selekční tlak během celého běhu a~neublíží jí ani \uv{superman},
ani změna měřítka fitness.

\paragraph{Boltzmannova selekce}

Boltzmannova selekce vybírá jedince s~pravděpodobností $p_i \propto e^{f(x_i)/T}$, kde $T$ je
\emph{teplota} klesající v~čase. Na začátku je $T$ vysoká a~selekce se blíží uniformnímu
výběru, takže převládá explorace; na konci je $T$ nízká a~selekce se stává téměř
deterministickou, tedy exploatační. Stejnou myšlenku používá \emph{Boltzmannův turnaj},
v~němž horší jedinec vyhraje s~pravděpodobností závislou exponenciálně na poměru rozdílu
fitness a~teploty. Souvislost se simulovaným žíháním je přímá (kap.~\ref{sec:lokalni}).

\paragraph{Environmentální selekce, elitismus, modely populace}

Rodičovská selekce rozhoduje o~tom, kdo se bude rozmnožovat; environmentální selekce
o~tom, kdo se dožije další generace. Podle toho, jak velkou část populace vymění, se
rozlišují tyto modely:

\begin{itemize}
\item \textbf{Generační model} (\emph{generational}) nahradí celou populaci
  potomky naráz, takže $|P'| = |P| = n$.
\item \textbf{Model s~ustáleným stavem} (\emph{steady-state}) mění populaci po
  kouscích: v~každém kroku vznikne jen $\lambda$ potomků, často 1--2, a~ti nahradí
  vybrané jedince. Nahradit lze nejhoršího, náhodně zvoleného, nebo nejpodobnějšího,
  čemuž se říká \emph{crowding}. Podíl nahrazované populace udává parametr
  \emph{generation gap} $G = \lambda/n$.
\item \textbf{Elitismus} (\emph{elitism}) bezpodmínečně kopíruje $k$ nejlepších
  jedinců do další generace. Zaručuje tím monotonii nejlepší nalezené fitness a~je
  nutnou podmínkou pro řadu konvergenčních vět; příliš silný elitismus však snižuje
  diverzitu.
\item \textbf{Deterministická selekce ES} nenechává nic náhodě: $(\mu+\lambda)$
  vybírá $\mu$ nejlepších z~rodičů i~potomků dohromady, a~má tedy implicitní
  elitismus, kdežto $(\mu,\lambda)$ vybírá jen z~potomků (kap.~\ref{sec:es}).
\end{itemize}

\begin{defbox}[Selekční tlak a~doba převzetí]
\emph{Selekční tlak} (\emph{selection pressure}) je míra zvýhodnění lepších
jedinců. Kvantifikuje se \emph{dobou převzetí} (\emph{takeover time}) $\tau^*$, tedy
počtem generací, za který jediná kopie nejlepšího jedince zaplní bez variace
celou populaci. Pro turnaj velikosti $k$ platí zhruba
$\tau^* \approx \ln n/\ln k$. Pro fitness-proporcionální selekci je doba převzetí
delší a~závisí na rozložení fitness.
\end{defbox}

\subsubsection{Teoretické meze: No Free Lunch}

Nabízí se otázka, jestli nejde navrhnout algoritmus, který bude lepší než ostatní na
čemkoli. Odpověď je záporná a~má podobu věty.

\begin{veta}[No Free Lunch, Wolpert \& Macready 1995--1997]
Zprůměrováno přes \emph{všechny} možné účelové funkce na konečné doméně mají
libovolné dva deterministické neopakující se prohledávací algoritmy stejný
očekávaný výkon. Neexistuje tedy univerzálně nejlepší optimalizační algoritmus.
\end{veta}

Praktický důsledek je jasný: vyplácí se stavět \textbf{doménově specifické varianty} EA,
tedy volit reprezentaci a~operátory tak, aby zohledňovaly strukturu úlohy. Co věta naopak
neříká, je to, že by si na \emph{reálných} úlohách byly všechny algoritmy rovny. Reálné funkce
tvoří velmi malou a~strukturovanou podmnožinu všech funkcí a~právě té struktury se dá využít.

\subsubsection{Historické větve evolučních algoritmů}

Evoluční algoritmy nevyrostly z~jediného kořene. Nezávisle na sobě vznikly čtyři školy,
které se lišily reprezentací i~tím, který operátor považovaly za hlavní:

\begin{center}
\small
\begin{tabular}{@{}>{\raggedright\arraybackslash}p{2.3cm}
                  >{\raggedright\arraybackslash}p{3.1cm}
                  >{\raggedright\arraybackslash}p{3.1cm}
                  >{\raggedright\arraybackslash}p{3.1cm}
                  >{\raggedright\arraybackslash}p{3.1cm}@{}}
\toprule
 & \textbf{GA} & \textbf{ES} & \textbf{EP} & \textbf{GP} \\
\midrule
autor, rok & Holland, 1975 & Rechenberg, Schwefel, 1964 &
  L.~Fogel, Owens, Walsh, 1965 & Koza, 1989--92 \\[2pt]
reprezentace & binární řetězec & vektor $\mathbb{R}^n$ + strategické parametry &
  konečný automat (genotyp = fenotyp) & syntaktický strom (LISP) \\[2pt]
hlavní operátor & křížení & mutace & mutace & křížení podstromů \\[2pt]
křížení & 1-bodové & rekombinace více rodičů (lokální/globální) &
  \emph{není} & výměna podstromů \\[2pt]
mutace & bit-flip $p_M$ & gaussovská, samoadaptivní $\sigma$ &
  \uv{chytré} mutace automatu & více typů \\[2pt]
rodičovská selekce & ruleta & náhodná & každý jednou & turnaj \\[2pt]
env.\ selekce & generační náhrada &
  deterministická $(\mu,\lambda)$, $(\mu+\lambda)$ &
  turnaj (rodiče $+$ potomci) & generační \\[2pt]
teorie & schémata & konvergence, 1/5 pravidlo & --- & bloat, schémata \\
\bottomrule
\end{tabular}
\end{center}

Dnes se hranice mezi větvemi stírají; EP a~ES třeba prakticky splynuly a~mluví se jednotně
o~evolučních algoritmech.

\subsection{Genetické algoritmy}
\label{ssec:sga}

\subsubsection{Jednoduchý genetický algoritmus (SGA)}

Původnímu Hollandovu návrhu z~roku 1975 se dnes říká \emph{simple genetic algorithm} (SGA).
Vystačí si s~minimální sadou operátorů a~s~nejjednodušším možným kódováním, a~to záměrně:
jednoduchost byla podmínkou toho, aby šel algoritmus teoreticky analyzovat (teorie schémat,
kap.~\ref{sec:schemata}).

\begin{defbox}[SGA]
SGA udržuje populaci $n$ binárních řetězců délky $m$. Rodiče vybírá
fitness-proporcionální (ruletovou) selekcí, kříží je jednobodovým křížením
s~pravděpodobností $p_C$ a~na každý bit potomka aplikuje bitovou mutaci
s~pravděpodobností $p_M$. Populace se nahrazuje generačně: $P(t+1)$ zcela
nahradí $P(t)$.
\end{defbox}

\begin{algorithm}[!ht]
\caption{Jednoduchý genetický algoritmus (SGA)}
\begin{algorithmic}[1]
\State $t \gets 0$; vygeneruj náhodně $P(0)$ = $n$ řetězců délky $m$
\While{není splněna ukončovací podmínka}
  \State spočítej $f(x)$ pro každé $x \in P(t)$
  \State $P(t+1) \gets \emptyset$
  \For{$n/2$ krát}
    \State vyber ruletou dvojici rodičů $x, y$ z~$P(t)$
    \State s~pravděpodobností $p_C$ zkřiž $x, y$ jednobodovým křížením
    \State každý bit $x$ i~$y$ převrať s~pravděpodobností $p_M$
    \State vlož $x, y$ do $P(t+1)$
  \EndFor
  \State $t \gets t+1$
\EndWhile
\end{algorithmic}
\end{algorithm}

Parametry se v~praxi volí takto: $n \in [50, 500]$, $p_C \in [0{,}6; 0{,}9]$
a~$p_M \approx 1/m$, což znamená, že se v~průměru změní jeden bit na jedince.

\subsubsection{Binární reprezentace}

SGA pracuje s~binárními řetězci, takže cokoli jiného než bity je nutné do bitů zakódovat.

\paragraph{Kódování celých a~reálných čísel.} Reálné číslo $x \in [a,b]$
zakódované na $k$ bitů jako celé číslo $z \in \{0,\dots,2^k-1\}$ dekódujeme vztahem
\[
x \;=\; a \;+\; z \cdot \frac{b-a}{2^k - 1}.
\]
Dosažitelná přesnost je $(b-a)/(2^k-1)$. Výhodou je explicitní kontrola nad přesností
a~komprese prohledávaného prostoru. Nevýhodou jsou \emph{Hammingovy útesy}
(\emph{Hamming cliffs}): dvě sousední hodnoty se mohou lišit ve všech bitech naráz,
třeba 0111 $\to$ 1000. Malá změna fenotypu pak vyžaduje velkou změnu genotypu, což
mutaci po jednotlivých bitech znesnadňuje.

\paragraph{Grayův kód.} Hammingovy útesy odstraňuje Grayův kód, v~němž se sousední celá
čísla liší právě v~jednom bitu. Z~binárního zápisu $b$ na Grayův $g$ se přechází předpisem
$g_1 = b_1$, $g_i = b_{i-1} \oplus b_i$; zpět pak $b_1 = g_1$,
$b_i = b_{i-1} \oplus g_i$.

\paragraph{Hollandův argument pro binární abecedu.} Holland zdůvodňoval volbu dvouprvkové
abecedy počtem schémat. Binární řetězce délky 100 kódují zhruba stejnou informaci jako
desítkové délky 30, obsahují ale více schémat ($3^{100}$ oproti $11^{30}$), takže GA
\uv{zpracovává} více informací naráz. Dnes víme, že schémata nejsou tak zásadní, jak se
Holland domníval, a~že rozhodující je \textbf{přirozenost kódování pro danou úlohu}. Pro
reálnou optimalizaci se proto používají reálné vektory, pro TSP permutace atd.

\subsubsection{Operátory křížení}

Genetický algoritmus stojí a~padá s~křížením, takže se vyplatí projít jeho základní varianty
a~všímat si přitom, jak moc která rozbíjí to, co už populace našla.

\begin{defbox}[Křížení (\emph{crossover, recombination})]
Operátor, který ze dvou (nebo více) rodičů vytvoří potomka kombinací jejich
genetického materiálu. V~GA je hlavním operátorem; vyjadřuje naději, že
rekombinací dobrých částí (\emph{stavebních bloků}) vzniknou ještě lepší
jedinci. Aplikuje se s~pravděpodobností $p_C$, jinak potomci = kopie rodičů.
\end{defbox}

\paragraph{Jednobodové křížení.} Náhodně se zvolí bod řezu $k \in \{1,\dots,m-1\}$
a~rodiče si vymění koncové části za ním.

\textbf{Příklad.} Rodiče $P_1 = 1011|0110$, $P_2 = 0010|1101$, bod řezu $k=4$:
\[
O_1 = \mathbf{1011}\,1101, \qquad O_2 = \mathbf{0010}\,0110.
\]

\paragraph{Dvoubodové křížení.} Zvolí se dva body $k_1 < k_2$ a~vymění se úsek mezi nimi.
Proti jednobodové variantě se řetězec chová jako kruh, takže geny na okrajích nejsou
znevýhodněné: u~jednobodového křížení je totiž pravděpodobnost rozbití úseku úměrná jeho
\emph{definující délce} (viz kap.~\ref{sec:schemata}).

\textbf{Příklad.} $P_1 = 10|1101|10$, $P_2 = 00|1011|01$ $\Rightarrow$
$O_1 = 10\,1011\,10$, $O_2 = 00\,1101\,01$.

\paragraph{Uniformní křížení.} Pro každou pozici $j$ se zvlášť hodí mince:
s~pravděpodobností $0{,}5$ pochází gen od prvního rodiče, jinak od druhého, přičemž druhý
potomek se sestaví komplementárně. Sousedící pozice tedy nejsou nijak zvýhodněné, zato
uniformní křížení silně narušuje schémata s~velkou definující délkou, takže má vysokou
\emph{disruptivitu}.

\textbf{Příklad.} $P_1 = 11111111$, $P_2 = 00000000$, maska $10110010$:
$O_1 = 10110010$, $O_2 = 01001101$.

Jak moc který operátor rozbíjí schémata, shrnuje následující tabulka:

\begin{center}
\begin{tabular}{@{}lll@{}}
\toprule
\textbf{operátor} & \textbf{p-st rozbití schématu délky $d$} & \textbf{poznámka} \\
\midrule
jednobodové & $d/(m-1)$ & poziční bias, chrání krátká schémata \\
dvoubodové & $\dfrac{2d(m-d)}{m(m-1)}$ & řetězec jako kruh, odstraňuje bias konců \\
$k$-bodové & roste s~$k$ & \\
uniformní & $1 - (1/2)^{o(S)-1}$ & distribuční bias, největší disruptivita \\
\bottomrule
\end{tabular}
\end{center}

\textbf{Poznámka k~dvoubodovému křížení.} Vzorec ve druhém řádku stojí za odvození.
Chápeme-li řetězec jako kruh, vybíráme dva ze $m$ řezů. Schéma je rozbito, padne-li dovnitř
jeho definujícího úseku o~$d$ mezerách \emph{právě jeden} z~nich, tj. s~pravděpodobností
$d(m-d)/\binom{m}{2} = 2d(m-d)/(m(m-1))$. Pro schéma sahající od prvního
k~poslednímu bitu, tedy pro $d = m-1$, klesne rozbití z~jistoty (jednobodové:
$d/(m-1) = 1$) na pouhé $2/m$, a~to je přesně ono \uv{odstranění okrajového
biasu}. U~krátkých schémat je tomu naopak: dvoubodové křížení je asi dvakrát
destruktivnější než jednobodové.

\subsubsection{Mutace a~inverze}

\paragraph{Bitová mutace.} Každý bit se nezávisle na ostatních překlopí
s~pravděpodobností $p_M$. V~SGA má mutace roli \emph{sekundárního} operátoru a~slouží
jako pojistka: brání uvíznutí v~lokálním extrému a~trvalé ztrátě alely z~populace.
V~EP a~v~raných ES je tomu jinak, tam je mutace jediným zdrojem variability.
Doporučenou hodnotou je $p_M \approx 1/m$.

\paragraph{Inverze.} Původní Hollandův návrh počítal ještě s~operátorem \emph{inverze}.
Ten obrátí pořadí části řetězce, ale \emph{zachová význam bitů}, protože každý gen si
s~sebou nese identifikátor své pozice. Cílem bylo nechat evolvovat samotné uspořádání genů
tak, aby spolupracující geny ležely blízko sebe a~nebyly rozbíjeny křížením. Technicky je
inverze náročná a~její přínos se neprokázal. U~permutační reprezentace se však používá jako
velmi účinná mutace (odst.~\ref{sec:reprez}).

\subsubsection{Škálování fitness}\label{ssec:skalovani}

Ruletová selekce reaguje na absolutní hodnoty fitness, a~to jí dělá potíže na obou koncích
běhu. Zpočátku bývá rozptyl fitness velký a~populaci ovládne několik \uv{supermanů}. Ke konci
je rozptyl naopak malý a~selekce ztrácí tlak. Oboje se dá napravit tím, že se fitness před
selekcí transformuje předpisem $f \mapsto f'$:

\begin{itemize}
\item \textbf{Lineární škálování} počítá $f' = a f + b$, kde se $a, b$ volí tak,
  aby $\bar{f'} = \bar{f}$ a~$f'_{\max} = C_{\mathrm{mult}} \cdot \bar{f}$,
  typicky s~$C_{\mathrm{mult}} \in [1{,}2; 2]$. Nutná je pojistka proti záporným
  hodnotám u~nejhorších jedinců.
\item \textbf{Sigma truncation} počítá $f' = \max\big(0,\; f - (\bar f - c\,\sigma_f)\big)$,
  kde $\sigma_f$ je směrodatná odchylka fitness a~$c \in [1;3]$. Odstraní
  odlehlé nízké hodnoty a~automaticky se přizpůsobuje rozptylu.
\item \textbf{Mocninné škálování} umocňuje, $f' = f^{\,k}$, přičemž $k$ se
  v~průběhu běhu zvyšuje, takže tlak roste.
\item \textbf{Pořadové škálování} nahradí fitness pořadím (viz
  odst.~\ref{ssec:selekce}). Je nejrobustnější a~dnes nejobvyklejší.
\item \textbf{Boltzmannovo škálování} používá $f' = e^{f/T}$ s~klesající teplotou.
\end{itemize}

\textbf{Příklad.} Vezměme populaci s~fitness $(169, 576, 64, 361)$, kde $\bar f = 292{,}5$
a~$f_{\max} = 576$, a~použijme lineární škálování s~$C_{\mathrm{mult}} = 1{,}5$.
Chceme $f'_{\max} = 438{,}75$ a~$\bar{f'} = 292{,}5$. Z~podmínek
$a f_{\max} + b = 438{,}75$ a~$a \bar f + b = 292{,}5$ dostáváme
$a = 146{,}25/283{,}5 = 0{,}516$ a~$b = 292{,}5 - 0{,}516 \cdot 292{,}5 = 141{,}6$.
Nové fitness jsou $(228{,}8;\, 438{,}8;\, 174{,}6;\, 327{,}9)$ a~podíl nejlepšího jedince
klesl z~$49{,}2\,\%$ na $37{,}5\,\%$. Selekční tlak se tím zmírnil, a~přitom ani nejhorší
jedinec neztrácí šanci.

\subsection{Reprezentace jedince a~genetické operátory}
\label{sec:reprez}

Volba reprezentace rozhoduje o~tom, které operátory mají vůbec smysl, a~tím zároveň
o~tom, jak vypadá \emph{fitness krajina} (\emph{fitness landscape}), tedy jak snadné
bude prohledávání. Při návrhu evolučního algoritmu je to nejdůležitější praktické
rozhodnutí. Následující oddíly proto procházejí obvyklé reprezentace jednu po druhé
a~u~každé ukazují, jaké mutace a~jaké křížení k~ní patří.

\subsubsection{Celočíselná reprezentace}

Jedinec je vektor $(z_1,\dots,z_m) \in \prod_j D_j$, kde $D_j$ jsou konečné domény.
Záleží ale na tom, jakou povahu hodnoty v~doméně mají, protože od ní se odvíjí volba
mutace. Rozlišujeme dva případy:
\begin{itemize}
\item \textbf{Kardinální (ordinální) atributy} mají smysluplné uspořádání a~metriku,
  takže se dá říci, které dvě hodnoty jsou si blízké. Patří sem třeba počet strojů
  nebo věk.
\item \textbf{Nominální atributy} jsou pouhé kódy bez vnitřního uspořádání, například
  barva nebo typ komponenty. O~\uv{blízkých} hodnotách tu nemá smysl mluvit.
\end{itemize}

\paragraph{Mutace.} Podle povahy atributu se používá jeden ze dvou operátorů:
\begin{itemize}
\item \emph{Nezaujatá} mutace (\emph{unbiased}, \emph{random resetting}) losuje novou
  hodnotu náhodně z~celé domény. Pro nominální atributy je to jediná korektní volba.
\item \emph{Zaujatá} mutace (\emph{biased}, \emph{creep mutation}) ponechá původní
  hodnotu a~posune ji o~malý náhodný krok, například z~diskretizovaného normálního
  rozdělení. Vhodná je jen pro ordinální atributy, protože jinde by \uv{malý krok}
  nic neznamenal.
\end{itemize}

\paragraph{Křížení.} S~křížením je to jednodušší, protože se na rozdíl od mutace nemusí
ohlížet na povahu hodnot. Používají se stejné operátory jako u~binární reprezentace, tedy
jednobodové, vícebodové i~uniformní. Jediný rozdíl je v~tom, že se místo jednotlivých
bitů kopírují celá čísla.

\subsubsection{Reálná reprezentace}

Jedinec je $\vec{x} = (x_1,\dots,x_n) \in \mathbb{R}^n$. Historicky se reálná čísla
kódovala do bitových řetězců a~pracovalo se s~nimi binárními operátory. Dnes se pracuje
přímo s~reálnými hodnotami a~operátory se tomu přizpůsobují.

\paragraph{Mutace.}
\begin{itemize}
\item \textbf{Gaussovská (normální) mutace} přičte ke~složce normálně rozdělený šum:
  $x_i' = x_i + \sigma_i \cdot N(0,1)$. Klíčovým parametrem je šířka kroku $\sigma$.
  Může být pevná, deterministicky klesající podle předem daného předpisu
  (např.\ $\sigma(t) = 1 - 0{,}9\,t/T$), adaptivní (pravidlo 1/5), nebo samoadaptivní,
  kdy se stane součástí genomu a~evolvuje spolu s~jedincem (viz kap.~\ref{sec:es}).
\item \textbf{Cauchyho mutace} má těžší chvosty, takže velké skoky nastávají častěji
  a~jedinec snáze unikne z~lokálního extrému. Používá se ve \emph{fast EP}.
\item \textbf{Uniformní mutace} losuje $x_i'$ náhodně z~intervalu $[a_i, b_i]$, a~je
  tedy nezaujatá.
\item Zvlášť je potřeba ošetřit meze prohledávaného prostoru. Nabízí se oříznutí
  (\emph{clipping}), odraz (\emph{reflection}), periodické zabalení (\emph{wrapping})
  nebo převzorkování, dokud výsledek nepadne dovnitř.
\end{itemize}

\paragraph{Křížení.} Rozlišujeme dva typy:
\begin{itemize}
\item \textbf{Strukturální (diskrétní)} křížení, tedy 1-bodové a~uniformní, hodnoty jen
  přeskupuje mezi rodiči. Žádné nové číslo nevznikne a~potomek proto leží ve vrcholu
  hyperkvádru daného rodiči.
\item \textbf{Aritmetické} křížení hodnoty kombinuje, takže potomek může nést čísla,
  která ani jeden rodič neměl.
\end{itemize}

\begin{defbox}[Aritmetické křížení (\emph{arithmetic crossover})]
Pro rodiče $\vec x, \vec y$ a~parametr $\alpha \in (0,1)$ vzniká potomek jako
\[
\vec z \;=\; \alpha \vec x + (1-\alpha)\vec y \qquad\text{(konvexní kombinace)}.
\]
Varianty se liší tím, kolika složek se kombinace týká. \emph{Jednoduché} křížení
kombinuje všechny složky a~je nejběžnější, \emph{jednoduché aritmetické} vybere náhodně
jedinou složku a~\emph{celé/single} se kombinuje s~jednobodovým křížením, takže
aritmeticky se míchá teprve za bodem řezu. Pro $\alpha = 1/2$ vyjde prostý průměr rodičů.
\end{defbox}

\textbf{Příklad.} $\vec x = (2{,}0;\, -1{,}0;\, 4{,}0)$,
$\vec y = (0{,}0;\, 3{,}0;\, 1{,}0)$, $\alpha = 0{,}25$:
\[
\vec z = 0{,}25\,\vec x + 0{,}75\,\vec y = (0{,}5;\; 2{,}0;\; 1{,}75).
\]

Čistě aritmetické křížení má ovšem podstatnou nevýhodu. Potomci leží uvnitř konvexního
obalu rodičů, takže rozptyl populace nutně klesá a~populace se \uv{smršťuje}. Proto se
používají rozšířené varianty, které dovolí potomkovi vyjít i~ven:
\begin{itemize}
\item \textbf{BLX-$\alpha$} (\emph{blend crossover}) losuje $z_i$ uniformně z~intervalu
  $[\min - \alpha I,\; \max + \alpha I]$, kde $I = |x_i - y_i|$; typicky se volí
  $\alpha = 0{,}5$. Interval je tedy oproti rodičům rozšířený na obě strany a~potomek
  smí ležet i~vně nich.
\item \textbf{SBX} (\emph{simulated binary crossover}) napodobuje chování jednobodového
  binárního křížení v~reálném prostoru: potomci jsou rozděleni symetricky kolem rodičů
  s~hustotou, kterou řídí parametr $\eta$. Je to standardní operátor v~NSGA-II.
\end{itemize}

\subsubsection{Permutační reprezentace}

Jedinec je permutace $\pi$ na $\{1,\dots,n\}$, což je přirozený zápis pro obchodního
cestujícího (TSP), rozvrhování, přiřazovací úlohy nebo uspořádání operací. Klasické
operátory zde ale selhávají, protože výsledek jednobodového křížení dvou permutací
obecně permutace není. Potřebujeme operátory, které \emph{zachovávají přípustnost}
a~zároveň dědí něco smysluplného. Co přesně se má dědit, závisí na úloze:

\begin{center}
\begin{tabular}{@{}lll@{}}
\toprule
\textbf{operátor} & \textbf{co zachovává} & \textbf{typické použití} \\
\midrule
PMX & absolutní pozice měst & úlohy s~významnou pozicí \\
OX & relativní pořadí & TSP (cyklické cesty) \\
CX & absolutní pozice (cykly) & úlohy s~významnou pozicí \\
ER & hrany (sousednosti) & TSP --- nejlepší kvalita \\
\bottomrule
\end{tabular}
\end{center}

\paragraph{PMX --- částečně mapované křížení}

\begin{defbox}[PMX (\emph{partially mapped crossover}, Goldberg)]
Dvoubodový operátor. (1)~Vyber úsek mezi dvěma body řezu a~vyměň jej mezi rodiči.
(2)~Z~vyměněných úseků odvoď \emph{mapování} odpovídajících si hodnot, tedy dvojic, které
se v~úsecích ocitly nad sebou. (3)~Zbylé pozice zkopíruj z~původního rodiče; kde by vznikl
konflikt, protože hodnota už v~potomkovi je, použij mapování, a~to případně i~opakovaně.
\end{defbox}

\textbf{Příklad.} $P_1 = (123|4567|89)$, $P_2 = (452|1876|93)$.
\begin{enumerate}
\item Výměna úseků: $O_1 = (\ldots|1876|\ldots)$, $O_2 = (\ldots|4567|\ldots)$.
\item Mapování z~úseků: $1{\leftrightarrow}4$, $8{\leftrightarrow}5$,
  $7{\leftrightarrow}6$, $6{\leftrightarrow}7$.
\item Doplnění bezkonfliktních pozic z~$P_1$ do $O_1$:
  $O_1 = (.23|1876|.9)$ (hodnoty 1 a~8 by kolidovaly).
\item Konflikty vyřešíme mapováním: $1 \to 4$, $8 \to 5$, tedy
  $O_1 = (423|1876|59)$. Analogicky $O_2 = (182|4567|93)$.
\end{enumerate}

\paragraph{OX --- křížení zachovávající pořadí}

\begin{defbox}[OX (\emph{order crossover}, Davis)]
Dvoubodový operátor zachovávající \emph{relativní pořadí}. (1)~Zkopíruj úsek mezi body
řezu z~prvního rodiče. (2)~Ze druhého rodiče čti hodnoty cyklicky počínaje pozicí za
druhým bodem řezu a~vynech ty, které už v~potomkovi jsou. (3)~Zbylé hodnoty vkládej do
volných pozic potomka, opět cyklicky od pozice za druhým bodem řezu.
\end{defbox}

\textbf{Příklad.} $P_1 = (123|4567|89)$, $P_2 = (452|1876|93)$.
\begin{enumerate}
\item $O_1$ převezme úsek z~$P_1$: $O_1 = ({\cdot}{\cdot}{\cdot}|4567|{\cdot}{\cdot})$.
\item Z~$P_2$ čteme cyklicky od pozice 8: $9,3,4,5,2,1,8,7,6$.
\item Vyškrtneme hodnoty už obsažené ($4,5,6,7$): zbývá $9,3,2,1,8$.
\item Doplňujeme od pozice 8 cyklicky (pozice $8,9,1,2,3$):
  $O_1 = (218|4567|93)$.
\item Symetricky (úsek z~$P_2$, pořadí z~$P_1$): $O_2 = (345|1876|92)$.
\end{enumerate}
Princip je tedy jednoduchý: vnitřní úsek se převezme beze změny a~zbytek se doplní
v~relativním pořadí druhého rodiče, jen se vynechají duplicity. Pro TSP je OX vhodnější
než PMX, protože u~okružní jízdy nezáleží na absolutní pozici města, ale na pořadí,
v~jakém se města projíždějí.

\paragraph{CX --- cyklické křížení}

\begin{defbox}[CX (\emph{cyclic crossover}, Oliver)]
Zachovává \emph{absolutní pozici}: každý gen potomka pochází ze stejné pozice u~některého
z~rodičů. Nejprve se sestrojí cykly. Začni na pozici 1 a~vezmi hodnotu z~$P_1$; podívej
se, jaká hodnota je na téže pozici v~$P_2$, najdi ji v~$P_1$ a~pokračuj stejným způsobem,
dokud se cyklus neuzavře. Sudé cykly potom ber z~jednoho rodiče, liché z~druhého.
\end{defbox}

\textbf{Příklad.} $P_1 = (123456789)$, $P_2 = (412876935)$.
Začneme hodnotou 1 z~$P_1$ na pozici 1. Pod ní je v~$P_2$ hodnota 4, tu najdeme
v~$P_1$ na pozici 4 $\Rightarrow$ přidáme; pod ní je 8, ta je v~$P_1$ na pozici
8 $\Rightarrow$ přidáme; pod ní je 3 (pozice 3), pod ní 2 (pozice 2) a~pod 2 je
1, čímž se cyklus uzavřel. Máme
$O_1 = (1\,2\,3\,4\,\_\,\_\,\_\,8\,\_)$; zbytek doplníme z~$P_2$:
$O_1 = (123476985)$, analogicky $O_2 = (412856739)$.

\paragraph{ER --- rekombinace hran}

Předchozí tři operátory spojuje jedno pozorování: PMX, OX i~CX zachovávají jen asi
$60\,\%$ hran rodičů. Pro TSP je ale právě \emph{hrana} nositelem kvality, protože délka
cesty je součtem délek hran. \emph{Edge recombination} (Whitley) proto pracuje rovnou
s~nimi:
\begin{enumerate}
\item Pro každé město sestaví \emph{seznam sousedů} z~\textbf{obou} rodičů.
\item Začne v~náhodném (nebo prvním) městě.
\item Dalším městem je ten dosud nenavštívený soused, který má \emph{nejméně}
  zbývajících sousedů; heuristika za tím je, aby žádné město nezůstalo \uv{osamocené}.
  Při shodě rozhodne náhoda.
\item Pokud seznam sousedů dojde prázdný, pokračuje se náhodným nenavštíveným
  městem. Nastává to jen v~$1$--$1{,}5\,\%$ případů.
\end{enumerate}
Vylepšená varianta \textbf{ER2} navíc značí symbolem \# hrany, které se vyskytují
v~\emph{obou} rodičích, a~při výběru je upřednostňuje.

\textbf{Příklad.} $P_1 = (123456789)$, $P_2 = (412876935)$ (cyklické cesty).
Seznamy sousedů: $1{:}\,9,2,4$; $2{:}\,1,3,8$; $3{:}\,2,4,9,5$; $4{:}\,3,5,1$;
$5{:}\,4,6,3$; $6{:}\,5,7,9$; $7{:}\,6,8$; $8{:}\,7,9,2$; $9{:}\,8,1,6,3$.
Startujeme v~1. Kandidáty jsou 9 (4~sousedi), 2 (3) a~4 (3), takže 9 vypadává a~mezi
2 a~4 losujeme; vyjde 4. Sousedé 4 jsou 3 a~5, přičemž 3 má ještě 3 volné sousedy
a~5 jen 2, bereme tedy 5, a~tak dále. Výsledkem je např.\ $(145678239)$.

\emph{Technická poznámka:} korektně se navštívené město nejprve odstraní ze
\emph{všech} seznamů a~teprve pak se počítají zbývající sousedé. Výše jsou u~startu
uvedeny počty ještě \emph{před} odebráním jedničky (po odebrání by bylo 9$:$3, 2$:$2,
4$:$2). Na pořadí kandidátů to zde nic nemění, ale při vlastním počítání na papíře je
třeba postupovat takto, jinak vycházejí jiné remízy.

\paragraph{Mutace permutací}

Mutace musí rovněž zachovat přípustnost, takže se omezuje na přeuspořádání již
přítomných hodnot:
\begin{itemize}
\item \textbf{Swap} prohodí dvě náhodně vybraná města.
\item \textbf{Insert} vyjme jedno město a~vloží je jinam, čímž se zbytek posune.
\item \textbf{Inverze (2-opt krok)} obrátí souvislý úsek cesty. Pro TSP je nejúčinnější,
  protože mění jen dvě hrany.
\item \textbf{Scramble} náhodně přeháže souvislý úsek.
\item Používá se také \textbf{výměna podcest} a~heuristické mutace, tedy 2-opt, 3-opt
  a~Lin-Kernighan.
\end{itemize}

\subsubsection{Další reprezentace}

Výčtem výše repertoár nekončí. Jedincem může být strom (odst.~\ref{sec:gp}), graf nebo
DAG, jak je tomu v~Cartesian GP a~v~neuroevoluci, případně matice, kterou se popisují
rozvrhy i~TSP zapsaný jako matice předchůdců. Dále se používají konečné automaty
v~evolučním programování, seznamy pravidel v~systémech LCS (odst.~\ref{sec:lcs})
a~celí agenti umělého života.

\subsection{Genetické a~evoluční programování}
\label{sec:gp}

\subsubsection{Historie}

Evoluci lze pustit i~na struktury, které samy něco počítají, a~ten nápad je podstatně
starší než většina evolučních algoritmů.
Myšlenku evoluce programů zmínil už Alan Turing v~50.\ letech. Forsyth ji v~roce 1980
aplikoval v~systému BEAGLE na rozpoznávání vzorů, Nichael Cramer (1985) poprvé popsal
stromové jedince a~John Koza formuloval v~letech 1989--1992 stromové genetické
programování v~podobě, v~jaké je známe dnes; nechal si je také patentovat.

\begin{defbox}[Genetické programování (\emph{genetic programming}, GP)]
Evoluční algoritmus, jehož jedinci jsou \textbf{spustitelné struktury}, tedy programy,
výrazy, obvody nebo modely. Fitness se určuje \emph{spuštěním} jedince na testovacích
datech. Velikost ani tvar jedince nejsou pevné a~mění se v~průběhu evoluce.
\end{defbox}

\begin{algorithm}[!ht]
\caption{Obecné schéma GP}
\begin{algorithmic}[1]
\State vygeneruj počáteční populaci náhodných programů
\While{není nalezeno dost dobré řešení}
  \State ohodnoť programy jejich \textbf{spuštěním} na trénovacích datech
  \State selekce podle fitness (typicky turnaj)
  \State křížení dvou programů, mutace programů
  \State vytvoř novou populaci
\EndWhile
\end{algorithmic}
\end{algorithm}

\subsubsection{Stromové GP}

\paragraph{Reprezentace.} Program je \emph{syntaktický strom}. Abychom takové stromy
mohli generovat, definujeme dvě množiny:
\begin{itemize}
\item \textbf{Množina terminálů} $T$ obsazuje listy stromu. Patří do ní proměnné,
  konstanty a~funkce bez argumentů, například čtení senzoru.
\item \textbf{Množina neterminálů (funkcí)} $F$ obsazuje vnitřní uzly a~u~každého prvku
  je určena jeho \emph{arita}. Bývají to aritmetické operace, logické spojky, podmínky
  nebo cykly.
\end{itemize}
Na obě množiny se kladou dva požadavky. \emph{Uzavřenost} (\emph{closure}) žádá, aby
každá funkce uměla zpracovat libovolný výstup libovolné jiné; odtud pochází
\uv{chráněné dělení} $\%$, které pro nulového dělitele vrací 1. \emph{Dostatečnost}
(\emph{sufficiency}) pak žádá, aby z~$T \cup F$ bylo řešení vůbec vyjádřitelné.

Příklad: strom \texttt{(+ (* x x) (sin y))} reprezentuje výraz $x^2 + \sin y$.

\paragraph{Inicializace.} Počáteční populaci lze vygenerovat několika způsoby, které se
liší tvarem výsledných stromů:
\begin{itemize}
\item \textbf{Grow} losuje uzly z~$T \cup F$, dokud se nedosáhne limitu hloubky nebo
  počtu uzlů. Vznikají tak nepravidelné stromy různé velikosti.
\item \textbf{Full} losuje do dané hloubky jen z~$F$ a~na poslední úrovni jen z~$T$,
  takže vznikají plné, pravidelné stromy.
\item \textbf{Ramped half-and-half} (Koza) vytvoří polovinu populace metodou grow
  a~polovinu metodou full, přičemž hloubky rozloží rovnoměrně (např.\ 2--6). Je to
  standardní volba. Vytváří spíše \uv{košaté} pravidelné tvary, což vyhovuje úlohám se
  symetriemi, kde jsou všechny vstupy podobně důležité.
\item \textbf{Uniformní vzorkování} všech možných stromů dává nepravidelnější tvary,
  protože se některé terminály ocitnou blíže ke kořeni.
\item \textbf{Seeding} vkládá do počáteční populace částečná či expertní řešení. Evoluci
  to zrychlí, hrozí ale předčasná konvergence.
\end{itemize}

\paragraph{Křížení.}
\begin{itemize}
\item \textbf{Výměna podstromů} je standardní Kozovo křížení: v~každém rodiči se zvolí
  náhodný uzel a~podstromy pod zvolenými uzly se vymění. Neterminály mívají vyšší pravděpodobnost
  výběru, typicky 90 \% neterminál a~10 \% terminál, protože jinak by se často
  vyměňovaly jen listy.
\item \textbf{Homologní křížení} se snaží zachovat pozici genetického materiálu.
  Nejprve se najde \emph{společná oblast} $C$, a~to procházením od kořene, dokud se
  shoduje arita uzlů:
  \begin{itemize}
  \item \emph{Jednobodové} křížení volí bod křížení ze společné oblasti.
  \item \emph{Uniformní} (GPUX, Poli a~Langdon) vymění každý uzel společné oblasti
    s~konstantní pravděpodobností. Uzly \emph{uvnitř} $C$ se vymění bez dopadu na jejich
    podstromy, zatímco uzly na \emph{hranici} $C$ se vymění i~s~podstromy. V~raných
    fázích se tak vyměňují velké podstromy a~s~konvergencí populace se operátor stává
    stále lokálnějším.
  \end{itemize}
\item \textbf{Křížení zachovávající kontext} (\emph{context preserving}) identifikuje
  uzel jeho \emph{souřadnicemi}, tedy posloupností odboček na cestě od kořene
  (např.\ $(2,1,3,1)$). \emph{Silná} varianta povoluje křížení jen mezi uzly se
  \emph{stejnými} souřadnicemi, \emph{slabá} je volnější.
\end{itemize}

\paragraph{Mutace.} V~GP je nutné (nikoli jen užitečné) používat \emph{více
typů} mutací, protože každý typ zasahuje strom jinak:
\begin{itemize}
\item \textbf{Podstromová mutace} nahradí náhodný podstrom novým náhodným.
\item \textbf{Size-fair} podstromová mutace navíc losuje velikost nového podstromu tak,
  aby odpovídala odstraněnému.
\item \textbf{Mutace uzlu} (\emph{node replacement}) zamění uzel za jiný \emph{stejné
  arity}.
\item \textbf{Hoist} nahradí jedince jeho vlastním podstromem, a~tím ho zmenší.
\item \textbf{Shrink} nahradí podstrom terminálem, což jedince zmenší také.
\item \textbf{Permutační mutace} prohodí podstromy jednoho neterminálu.
\item \textbf{Mutace konstant} řeší to, že GP tradičně špatně dolaďuje číselné hodnoty.
  Pomáhá aritmetická mutace konstant nebo rovnou lokální optimalizace \emph{všech}
  konstant ve stromu, ať už hill climbingem, nebo gradientní metodou; jde tedy
  o~memetický přístup.
\end{itemize}

\paragraph{Fitness a~selekce.} Fitness se odvozuje od toho, co má program dělat:
u~symbolické regrese je to chyba na trénovacích datech, jinde počet správných odpovědí
nebo kvalita chování v~simulaci. Selekce bývá turnajová. Protože se program učí z~dat,
je riziko přeučení stejné jako u~strojového učení a~řeší se stejně, tedy validační
množinou a~penalizací složitosti.

\subsubsection{Bloat}

Se stromovou reprezentací je spojený jeden charakteristický problém, který u~pevně
dlouhého genomu vůbec nemůže nastat.

\begin{defbox}[Bloat]
\emph{Bloat} je tendence programů v~GP růst do velikosti bez odpovídajícího zlepšení
fitness. Typicky se v~programech hromadí \emph{introny}, tedy části kódu, které nemají
vliv na výstup; příkladem je \texttt{(+ x 0)} nebo \texttt{(if false ...)}.
\end{defbox}

Důvody jsou dva a~oba jsou docela pochopitelné. V~prostoru programů je mnohem víc velkých
programů se stejným chováním než malých, takže už samotný neutrální drift tlačí
k~růstu. Navíc introny chrání jedince před destruktivním křížením, a~tak se
\uv{svezou} s~dobrým kódem.

V~přírodě růst naráží na fyzikální a~energetické limity, které ho samy zastaví, v~GP
ale žádná taková brzda není a~s~bloatem se musí bojovat explicitně:
\begin{itemize}
\item \textbf{Tvrdé limity} na hloubku či počet uzlů. Potomek, který limit překročí, se
  zahodí nebo nahradí rodičem.
\item \textbf{Penalizace ve fitness} (\emph{parsimony pressure}) odečte od fitness člen
  úměrný velikosti: $f' = f - \alpha \cdot \text{velikost}$. \emph{Lexikografická}
  varianta velikost nezapočítává, jen při shodné fitness nechá vyhrát menšího jedince.
\item \textbf{Anti-bloat operátory} hlídají, aby křížení a~mutace jedince příliš
  nezvětšovaly (to zajišťují size-fair varianty), a~doplňují je speciální zmenšující
  mutace \emph{hoist} a~\emph{shrink}.
\item \textbf{Vícekriteriální přístup} bere velikost jako druhé optimalizované kritérium
  (odst.~\ref{sec:moea}).
\end{itemize}

\subsubsection{ADF a~typované GP}

Stromové GP se dá výrazně posílit tím, že se programu dovolí mít vlastní podprogramy.

\begin{defbox}[ADF (\emph{automatically defined functions})]
Podprogramy evolvované společně s~hlavním programem. Jsou nutnou podmínkou
\textbf{modularity}, vyšší složitosti a~schopnosti zachytit symetrie úlohy. Jedinec se
skládá z~větve \texttt{main} a~jedné či více větví ADF; každá ADF je charakterizována
svou \emph{aritou} a~může mít omezenou (jinou) množinu terminálů a~neterminálů. Volání
ADF se v~hlavním programu chová jako \textbf{nový neterminál} dané arity. Genetické
operátory pracují na větvích \textbf{odděleně}, takže se main kříží s~mainem a~ADF$_1$
s~ADF$_1$.

Myšlenku lze rozšířit na automaticky definované smyčky, iterace a~rekurzi. Alternativou
je \emph{koevoluce} dvou populací, hlavních programů a~podprogramů
(odst.~\ref{sec:koevoluce}).
\end{defbox}

\textbf{Vynucení struktury.} Prohledávaný prostor se dá také omezit tím, že se zakážou
nesmyslné stromy. \emph{Silně typované GP} dává každému uzlu typ vstupů a~výstupu
a~operátory smějí vytvářet jen typově korektní stromy. \emph{Gramatické GP} místo toho
popíše přípustné tvary stromů bezkontextovou gramatikou. Obojí redukuje prohledávaný
prostor na smysluplné programy.

\subsubsection{Lineární a~grafové GP}

Strom není jediná struktura, do které se program dá zapsat. Následující varianty GP se
liší právě tím, jak jedince reprezentují, a~s~tím se mění i~celá sada operátorů.

\paragraph{Lineární GP.} Program je posloupnost instrukcí, často ve strojovém nebo byte
kódu, které pracují nad registry. Výhod je několik: operátory jsou jednoduché, protože
se kříží a~mutuje jako u~vektorů, interpretace je rychlá a~celá reprezentace je
přirozeně blízká skutečnému hardwaru. Nevýhodou je vysoké riziko vzniku nesmyslných
programů. Lineární GP je oblíbené v~umělém životě a~při evoluci herních botů. Slavným
příkladem je \textbf{Tierra} (T.~Ray), simulace umělého ekosystému, ve které jsou
jedinci sebekopírující programy ve strojovém kódu soutěžící o~paměť a~procesorový čas;
emergentně v~ní vznikli paraziti, hyperparaziti a~obrana proti nim. Následovníky jsou
Avida a~Avida-ED.

\paragraph{Grafové GP.} Program je obecný, často acyklický graf. Původně šlo
o~rozšíření stromového GP na paralelní programy, ale ukázalo se, že grafy dobře popisují
také elektrické obvody, konečné automaty, neuronové sítě a~plány. Problémem jsou složité
operátory, protože není zřejmé, jak křížit obecné grafy.

\paragraph{Kartézské GP (CGP).} \mimo{} Elegantní obchvat: DAG se reprezentuje svým
rovinným obrazem v~2D mřížce se souřadnicemi uzlů (fenotyp), zatímco genotyp je
\textbf{lineární} posloupnost celých čísel, která pro každý uzel popisuje jeho funkci
(index do tabulky funkcí) a~jeho vstupy a~nakonec vyjmenuje výstupní uzly grafu. Z~toho
plyne několik vlastností:
\begin{itemize}
\item Mutace je bodová mutace lineárního genomu. Musí se přitom hlídat, aby vznikaly jen
  platné funkce a~platná propojení, což znamená omezení hloubky spojů.
\item Malá změna genotypu může mít velký dopad na fenotyp. Naopak řada mutací je
  \emph{neutrální}, protože uzly nepřipojené k~výstupu tvoří přirozené introny, což
  CGP prospívá.
\item Křížení je málo důležité. Naivní jednobodové je příliš destruktivní, a~tak se
  používá aritmetická varianta a~řezy na hranicích uzlů.
\item Selekce je obvykle $(1+\lambda)$-ES s~$\lambda = 4$.
\item Genom má konstantní délku, takže bloat je omezen už samotnou mřížkou.
\end{itemize}

\paragraph{Vývojové (developmental) GP.} Místo přímé evoluce grafu se evolvuje
\emph{strom-program, který graf postupně staví}. Začíná se \uv{embryem} a~aplikují se
operace pro strukturu (sériové či paralelní zapojení, třícestné dělení) a~pro elementy
(vlož rezistor, vlož kondenzátor). Takto postupoval Koza při evoluci elektrických
obvodů, kde se podařilo znovuobjevit patentovaná zapojení, Gruau ve svém
\emph{cellular encoding} pro neuronové sítě a~AutoML při evoluci modelů scikit-learn.

\subsubsection{Gramatická evoluce}

Gramatiku lze posunout ještě dál než u~gramatického GP: nemusí jen filtrovat přípustné
stromy, může být přímo tím, co z~genomu program odvozuje.

\begin{defbox}[Gramatická evoluce (\emph{grammatical evolution}, GE)]
Genotyp je \textbf{lineární seznam celých čísel} proměnné délky, fenotyp je program
odvozený z~bezkontextové gramatiky. Číslo z~genomu určuje, které pravidlo se použije pro
expanzi nejlevějšího neterminálu:
\[
\text{pravidlo} \;=\; (\text{gen} \bmod \text{počet pravidel pro tento neterminál}).
\]
\end{defbox}

Použití modula činí kódování \textbf{robustním}, protože každý genom je
interpretovatelný. A~protože je genom lineární, lze používat \textbf{standardní operátory
GA}, což je hlavní praktická výhoda celého přístupu. Gramatika navíc zaručuje
syntaktickou správnost výsledku.

Nevýhodou je nelokálnost: malá změna genotypu, zejména na jeho začátku, může zcela
změnit fenotyp (\emph{ripple effect}). Nedokončí-li se derivace, čte se genom znovu od
začátku (\emph{wrapping}), a~to s~limitem na počet obalení.

\subsubsection{Evoluční programování}

K~evoluci programů se dá dojít i~úplně jinou cestou, než jakou šel Koza, a~historicky
se tak stalo dřív.

\begin{defbox}[Evoluční programování (\emph{evolutionary programming}, EP)]
Větev EA (L.~Fogel, 1965, tedy starší než GA), jejímž cílem bylo evolvovat \uv{umělou
inteligenci}: agenta, který predikuje stav prostředí a~volí vhodnou akci. Agent je
reprezentován \textbf{konečným automatem} a~platí \textbf{genotyp = fenotyp}, žádné
zvláštní kódování se tedy nepoužívá. \textbf{Křížení se nepoužívá} vůbec a~veškerá
variabilita pochází z~mutací.
\end{defbox}

\paragraph{EP nad automaty.} Prostředí je posloupnost symbolů konečné abecedy. Automat
dostává symboly na vstup a~jeho výstup se porovnává s~následujícím symbolem posloupnosti;
fitness je úspěšnost predikce, měřená absolutní nebo střední kvadratickou chybou, často
s~penalizací za velikost automatu. Mutovat se dá pěti způsoby: změnou výstupního symbolu,
změnou cílového stavu přechodu, přidáním stavu, odebráním stavu a~změnou počátečního
stavu. Novými automaty se nahrazuje polovina populace.

\paragraph{Slavný příklad --- predikce prvočísel.} Cílem je predikovat posloupnost
\[
111010100010100010100010000010\ldots,
\]
kde 1 znamená \uv{toto číslo je prvočíslo}; automat se tedy má naučit Eratosthenovo
síto. Postupuje se iterativně: nejprve se učí krátká posloupnost, pak se prodlužuje
o~jeden symbol. Fitness dává bod za každý uhodnutý symbol a~odečítá
$0{,}01 \cdot N$ za počet stavů. Výsledek je poučný. Automaty se postupně
zjednodušovaly, až nakonec vznikl jednostavový automat generující stále 0. Prvočísla
jsou totiž řídká, takže konstantní odpověď \uv{není prvočíslo} má vysokou úspěšnost.
Je to klasická ukázka toho, jak evoluce zneužije špatně navrženou fitness.

\paragraph{EP dnes.} Z~původní podoby zůstalo několik zásad:
\begin{itemize}
\item Kódování má být \emph{přirozené} pro úlohu, obvykle tedy genotyp = fenotyp.
\item Používá se sada \uv{chytrých}, doménově specifických mutací a~žádné křížení.
\item Rodičovská selekce je nastavena tak, že každý jedinec je právě jednou rodičem.
\item Environmentální selekce probíhá turnajem mezi rodiči a~potomky.
\item EP nad reálnými vektory obsahuje, stejně jako ES, rozptyly mutací přímo v~genomu
  ($\sigma_j' = \sigma_j(1 + \alpha N(0,1))$, poté $x_j' = x_j + \sigma_j' N(0,1)$).
  EP a~ES tak prakticky splynuly.
\end{itemize}

\subsubsection{Je křížení k~něčemu? Experiment s~bezhlavým kuřetem}

Tradiční obhajoba křížení stojí na tom, že si rodiče vyměňují stavební bloky. Kritika,
typická pro EP, oponuje, že křížení je jen podivná velká náhodná mutace, která navíc
nerespektuje kódování. Spor se dá rozhodnout experimentem.

\begin{defbox}[Headless chicken test (Jones, 1995)]
Porovnáme standardní křížení s~\emph{náhodným} křížením, kde se jedinec kříží
s~\textbf{náhodně vygenerovaným} řetězcem. Je-li výsledek stejný nebo lepší, pak zkoumaný
operátor ve skutečnosti nefunguje jako křížení, ale jako makromutace:
\uv{nenazývejme bezhlavé kuře kuřetem, i~když má mnoho kuřecích rysů}.
\end{defbox}

Interpretace je přímočará. Existují-li v~úloze jasné stavební bloky, je pravé křížení
lepší. Vyhraje-li náhodné křížení, znamená to, že v~dané úloze nebo v~daném kódování
stavební bloky nejsou; na vině je pak špatné kódování, špatný operátor, nebo prostě
úloha, pro kterou se křížení nehodí.

\subsection{Shrnutí ke~zkoušce}

\begin{itemize}
\item Evoluční algoritmus je populační stochastické prohledávání a~běží v~cyklu
  \emph{inicializace $\to$ ohodnocení $\to$ rodičovská selekce $\to$ rekombinace $\to$
  mutace $\to$ ohodnocení $\to$ environmentální selekce}.
\item Variace, tedy křížení a~mutace, vytvářejí v~populaci diverzitu, kdežto selekce ji
  odebírá a~tím směřuje hledání. Navrhnout evoluční algoritmus proto znamená najít rovnováhu
  mezi explorací a~exploatací.
\item Selekční schémata se liší vyvíjeným tlakem a~rozptylem počtu potomků. \textbf{Ruleta}
  vybírá rodiče s~pravděpodobností $p_i = f_i/\sum f_j$, má vysoký rozptyl a~je citlivá na
  škálování fitness. \textbf{SUS} zachovává stejné střední hodnoty, ale rozptyl srazí na
  minimum. \textbf{Turnaj} řídí tlak velikostí $k$, je invariantní k~transformacím fitness
  a~vystačí si bez explicitní fitness. \textbf{Pořadová} selekce drží tlak konstantní,
  zatímco u~\textbf{Boltzmannovy} tlak s~časem roste.
\item Elitismus zaručuje, že nejlepší dosažená fitness už nikdy neklesne, tedy její monotonii.
  Generační a~steady-state model se od sebe liší tím, jak velkou část populace v~jednom kroku
  nahradí.
\item Věta No Free Lunch říká, že univerzálně nejlepší algoritmus neexistuje. Důsledkem je
  specializace reprezentace i~operátorů na konkrétní doménu.
\item Historické větve GA, ES, EP a~GP i~jejich charakteristické volby shrnuje srovnávací
  tabulka a~je dobré umět je vyjmenovat zpaměti.
\item SGA kombinuje binární řetězce, ruletovou selekci, 1-bodové křížení, bit-flip mutaci
  a~generační náhradu. Nastavují se u~něj parametry $n$, $p_C \approx 0{,}7$
  a~$p_M \approx 1/m$.
\item Reálná čísla se do binárního řetězce kódují vztahem $x = a + z(b-a)/(2^k-1)$. Sousední
  hodnoty se přitom mohou lišit v~mnoha bitech naráz a~tyto Hammingovy útesy odstraňuje
  Grayův kód.
\item Operátory křížení se liší tím, jak silně narušují schémata. Jednobodové rozbije schéma
  s~pravděpodobností $d(S)/(m-1)$, dvoubodové nemá okrajový bias a~uniformní je nejvíce
  disruptivní, zato bez pozičního biasu.
\item Mutace má v~GA sekundární roli: brání nevratné ztrátě alel. Inverze je historický
  operátor, kterým se evolvovalo uspořádání genů.
\item Škálování fitness (lineární, sigma truncation, mocninné, pořadové a~Boltzmannovo)
  opravuje dvě slabiny fitness-proporcionální selekce: dominanci supermanů na začátku běhu
  a~ztrátu tlaku na jeho konci.
\item Reprezentace spolu s~operátory určuje fitness krajinu. Pro nominální atributy přicházejí
  v~úvahu jen nezaujaté mutace, pro ordinální i~zaujaté (creep).
\item V~reálné reprezentaci se mutuje gaussovsky, $x' = x + \sigma N(0,1)$. Aritmetické
  křížení $\alpha \vec x + (1-\alpha)\vec y$ zmenšuje rozptyl populace, a~proto se místo něj
  používá BLX-$\alpha$ a~SBX.
\item Permutační křížení se rozlišují podle toho, co po rodičích dědí: PMX přenáší pozice
  a~zbytek doplní mapováním, OX relativní pořadí, CX absolutní pozice pomocí cyklů a~ER hrany,
  což je pro TSP nejlepší volba. PMX a~OX je nutné umět předvést na konkrétním osmiprvkovém
  příkladu!
\item Permutace se mutují operátory swap, insert, inverze (2-opt) a~scramble.
\item Stromové GP podle Kozy staví programy z~terminálů a~neterminálů, které musí splňovat
  uzavřenost a~dostatečnost. Populace se inicializuje metodou grow, full nebo ramped
  half-and-half, křížením je výměna podstromů (neterminály se vybírají s~vyšší
  pravděpodobností), mutací existuje více typů včetně mutace konstant a~fitness se získá
  spuštěním programu.
\item Bloat je růst velikosti jedinců bez zlepšení fitness, na kterém se podílejí introny.
  Vysvětlují ho tři teorie: replication accuracy, removal bias a~crossover bias. Bránit se
  proti němu dá limity velikosti, penalizací (parsimony pressure), anti-bloat operátory
  (hoist, shrink, size-fair) a~vícekriteriálním přístupem.
\item ADF jsou evolvované podprogramy, které do GP vnášejí modularitu; operátory pak pracují
  na jednotlivých větvích odděleně. Typované GP a~gramatické GP omezují prohledávaný prostor
  na smysluplné programy.
\item Z~variant stojí za zapamatování lineární GP (byte kód, Tierra), grafové GP, Cartesian GP
  (lineární genotyp $\to$ 2D DAG, bodová mutace, $(1+4)$-ES) a~vývojové GP, kde strom teprve
  staví výsledný graf, tedy obvod nebo síť.
\item Gramatická evoluce pracuje s~lineárním genomem celých čísel, jimiž se pomocí modula
  vybírají pravidla gramatiky, a~vzniklý derivační strom dává fenotyp. Kódování je robustní
  a~operátory standardní jako v~GA, mapování je ale nelokální.
\item EP evolvuje konečné automaty, genotyp splývá s~fenotypem, používá se jen mutace
  a~turnajová selekce z~rodičů i~potomků. Učebnicovým příkladem je predikce prvočísel, na
  které se ukáže degenerace fitness.
\item Headless chicken experiment ověřuje, zda křížení opravdu kombinuje stavební bloky, nebo
  zda funguje pouze jako makromutace.
\end{itemize}

%=====================================================================
\section{Teorie schémat a~pravděpodobnostní modely SGA}
\label{sec:schemata}
%=====================================================================

Tahle kapitola odpovídá na otázku, \emph{proč} genetické algoritmy vlastně fungují.
Nejprve přijde klasická Hollandova teorie schémat i~s~kritikou, která k~ní neodmyslitelně
patří, a~potom exaktní Voseho model, který teorii schémat nahradil.

\subsection{Teorie schémat}
\label{ssec:teorie-schemat}

Teorie schémat je vůbec první pokus vysvětlit úspěch genetických algoritmů formálně, a~ne
jen odkazem na to, že se v~praxi osvědčily. Zformuloval ji Holland (1975) a~do širokého
povědomí ji dostal Goldberg (1989). Dnešní pohled na ni je smířlivý, ale střízlivý:
historicky je to zásadní krok, jako vysvětlení ale nestačí. U~zkoušky se přesto probírá
standardně a~její kritika váží stejně jako věta sama.

\subsubsection{Schéma, jeho řád a~definující délka}

Základním pojmem je schéma. Je to způsob, jak mluvit najednou o~celé skupině jedinců,
kteří se shodují v~tom podstatném a~liší se jen v~pozicích, na nichž zrovna nezáleží.

\begin{defbox}[Schéma (\emph{schema})]
Nechť jedinec je slovo délky $m$ nad abecedou $\{0,1\}$. \textbf{Schéma} $S$ je
slovo délky $m$ nad abecedou $\{0,1,*\}$, kde $*$ znamená \uv{na hodnotě
nezáleží} (\emph{don't care}). Schéma reprezentuje množinu všech jedinců, kteří
se s~ním shodují na všech pevných pozicích. Obsahuje-li $r$ hvězdiček,
reprezentuje $2^r$ jedinců.
\end{defbox}

Například schéma $1*0*$ zastupuje čtveřici $\{1000, 1001, 1100, 1101\}$, zatímco
schéma $**\ldots*$ ze samých hvězdiček zastupuje rovnou celý prostor.

\begin{defbox}[Řád, definující délka, fitness schématu]
Pro schéma $S$ délky $m$ definujeme:
\begin{itemize}
\item \textbf{Řád} $o(S)$ je počet pevných pozic, tedy počet nul a~jedniček
  ve~schématu.
\item \textbf{Definující délka} $d(S)$ je vzdálenost mezi první a~poslední
  pevnou pozicí; pro schéma s~jedinou pevnou pozicí vychází $d(S) = 0$.
\item \textbf{Fitness schématu} $F(S,t)$ je průměrná fitness těch jedinců
  v~populaci $P(t)$, kteří schéma $S$ reprezentují. \emph{Pozor:} tato hodnota
  závisí na aktuální populaci, nikoli jen na $S$ a~$f$.
\end{itemize}
\end{defbox}

\textbf{Příklad.} Schéma $S = 1{*}{*}0{*}1$ délky $m=6$ má pevné pozice 1, 4 a~6,
takže jeho řád je $o(S) = 3$; první a~poslední z~nich dělí $d(S) = 6 - 1 = 5$ míst.

\paragraph{Kombinatorika schémat.}
Než přijde věta samotná, vyplatí se spočítat, s~kolika schématy vlastně máme co do činění.
\begin{itemize}
\item Schémat délky $m$ je celkem $3^m$, protože na každé pozici může stát nula,
  jednička, nebo hvězdička.
\item Jeden jedinec délky $m$ je reprezentantem $2^m$ schémat: na každé pozici lze
  buď ponechat jeho bit, nebo tam napsat $*$.
\item Populace $n$ jedinců reprezentuje mezi $2^m$ a~$n \cdot 2^m$ schématy. Kde přesně
  v~tomto rozmezí leží, závisí na míře podobnosti jedinců.
\end{itemize}

\subsubsection{Hollandova věta o~schématech}

Věta o~schématech odhaduje zdola, kolik reprezentantů daného schématu bude v~populaci
o~generaci později.

\begin{veta}[Věta o~schématech (TST), Holland 1975]
V~jednoduchém genetickém algoritmu s~fitness-proporcionální selekcí,
jednobodovým křížením s~pravděpodobností $p_C$ a~bitovou mutací
s~pravděpodobností $p_M$ platí pro očekávaný počet reprezentantů schématu $S$
v~následující generaci:
\[
\mathbb{E}\big[C(S,t+1)\big] \;\ge\; C(S,t)\,\frac{F(S,t)}{\bar f(t)}
\left[\,1 - p_C \frac{d(S)}{m-1} - p_M\, o(S)\right].
\]
Slovně řečeno: schémata, která jsou \textbf{krátká} (mají malou definující délku),
\textbf{nízkého řádu} a~\textbf{nadprůměrná}, se v~po sobě jdoucích generacích
množí exponenciálně.
\end{veta}

\paragraph{Odvození}

Označme $C(S,t)$ počet jedinců populace $P(t)$, kteří reprezentují $S$, dále
$n = |P(t)|$ velikost populace a~$\bar f(t)$ její průměrnou fitness. Odvození
postupuje ve třech krocích, na každý operátor jeden.

\paragraph{Krok 1 --- selekce.} Pravděpodobnost, že ruleta vybere jedince $v$, je
$p_s(v) = f(v)/F(t)$, kde $F(t) = \sum_{u \in P(t)} f(u)$. Sečtením přes všechny
reprezentanty schématu dostaneme pravděpodobnost, že vybraný jedinec reprezentuje $S$:
\[
p_s(S) \;=\; \sum_{v \in S \cap P(t)} \frac{f(v)}{F(t)}
\;=\; \frac{C(S,t)\,F(S,t)}{F(t)}.
\]
Nová populace vzniká $n$ nezávislými tahy, takže
\[
\mathbb{E}\big[C(S,t+1)\big] \;=\; n\, p_s(S) \;=\; C(S,t)\,\frac{F(S,t)}{\bar f(t)},
\qquad \bar f(t) = F(t)/n .
\]
Je-li schéma nadprůměrné o~$\varepsilon$, tj.\
$F(S,t) = \bar f(t)\,(1+\varepsilon)$, a~drží-li se tento poměr v~čase, dostáváme
\[
C(S,t+1) = C(S,t)(1+\varepsilon)
\qquad\Longrightarrow\qquad
C(S,t) = C(S,0)\,(1+\varepsilon)^{t},
\]
tedy \textbf{exponenciální růst}, případně exponenciální pokles pro $\varepsilon<0$.
Samotná selekce tedy schéma rozmnožuje, ostatní operátory mu ubližují.

\paragraph{Krok 2 --- křížení.} Jednobodové křížení vybírá bod řezu z~$m-1$
možností. Schéma se \emph{jistě} rozpadne, padne-li řez mezi jeho první
a~poslední pevnou pozici, což nastane s~pravděpodobností nejvýše $d(S)/(m-1)$.
Slovo \uv{nejvýše} je podstatné: i~při řezu uvnitř schématu může být schéma
obnoveno, pokud druhý rodič nese na týchž pozicích tytéž hodnoty. Ke~křížení
navíc dochází jen s~pravděpodobností $p_C$, takže schéma přežije s~pravděpodobností
\[
p_{\mathrm{surv}}^{\mathrm{cross}}(S) \;\ge\; 1 - p_C\,\frac{d(S)}{m-1}.
\]

\paragraph{Krok 3 --- mutace.} Mutace schématu neublíží, dokud se nezmění žádná
z~jeho $o(S)$ pevných pozic; co dělá na hvězdičkových pozicích, je schématu jedno:
\[
p_{\mathrm{surv}}^{\mathrm{mut}}(S) = (1-p_M)^{o(S)} \;\approx\; 1 - p_M\,o(S)
\quad \text{pro } p_M \ll 1 .
\]

\paragraph{Složení.} Zbývá obě přežití složit. Předpokládáme, že jsou oba
destruktivní jevy (přibližně) nezávislé, a~jejich součin zanedbáme:
\[
\mathbb{E}\big[C(S,t+1)\big] \;\ge\; C(S,t)\,\frac{F(S,t)}{\bar f(t)}
\left[1 - p_C \frac{d(S)}{m-1} - p_M o(S)\right]. \qquad\blacksquare
\]

\paragraph{Číselný příklad}

Na číslech je vidět, co věta vlastně tvrdí. Vezměme Goldbergovu populaci ze
str.~\pageref{ex:goldberg}, tedy $m=5$ a~$n=4$: jedinci $01101\,(169)$,
$11000\,(576)$, $01000\,(64)$ a~$10011\,(361)$ dávají $\bar f = 292{,}5$.
Operátory nastavíme na $p_C = 0{,}7$ a~$p_M = 0{,}001$.

\begin{center}
\begin{tabular}{@{}lccccc@{}}
\toprule
schéma $S$ & reprezentanti & $F(S)$ & $o(S)$, $d(S)$ & přežití & $\mathbb{E}[C(S,t{+}1)]$ \\
\midrule
$1{*}{*}{*}{*}$ & 11000, 10011 & $468{,}5$ & $1;\ 0$ &
  $1 - 0 - 0{,}001 = 0{,}999$ &
  $2 \cdot 1{,}602 \cdot 0{,}999 = 3{,}2$ \\
$1{*}{*}{*}1$ & 10011 & $361$ & $2;\ 4$ &
  $1 - 0{,}7 - 0{,}002 = 0{,}298$ &
  $1 \cdot 1{,}234 \cdot 0{,}298 = 0{,}37$ \\
$0{*}{*}{*}{*}$ & 01101, 01000 & $116{,}5$ & $1;\ 0$ &
  $0{,}999$ &
  $2 \cdot 0{,}398 \cdot 0{,}999 = 0{,}80$ \\
\bottomrule
\end{tabular}
\end{center}

Tabulka ukazuje všechny tři možné osudy schématu. Krátké nadprůměrné schéma
$1{*}{*}{*}{*}$ se rozšíří ze dvou zástupců zhruba na 3,2. Schéma $1{*}{*}{*}1$ je
sice také nadprůměrné ($F/\bar f = 1{,}23$), jenže jeho velká definující délka je
fatální: křížení ho rozbije a~očekáváme, že vymizí. A~podprůměrné $0{*}{*}{*}{*}$
prostě ubývá. To je přesně obsah věty.

\subsubsection{Hypotéza stavebních bloků}

Věta o~schématech říká, které kusy řešení v~populaci přežívají. Hypotéza stavebních
bloků z~toho dělá tvrzení o~tom, jak genetický algoritmus hledá.

\begin{defbox}[Hypotéza stavebních bloků (\emph{building blocks hypothesis})]
Genetický algoritmus hledá dobré řešení tak, že \textbf{rekombinuje krátká,
nízkého řádu a~nadprůměrná schémata}, tzv.\ \emph{stavební bloky},
do stále lepších řešení.
Goldberg to popsal metaforou: \uv{tak jako dítě staví hrad z~jednoduchých dřevěných
kostek, tak GA hledá téměř optimální řešení skládáním stavebních bloků}.
\end{defbox}

Z~hypotézy plynou tři praktické důsledky:
\begin{itemize}
\item \textbf{Na kódování záleží.} Geny, které spolu interagují, mají v~řetězci ležet
  blízko sebe (\emph{linkage}), protože jinak je křížení rozbíjí dřív, než se stihnou
  osvědčit. Odtud pocházejí snahy o~inverzi, \emph{messy GA}, \emph{linkage learning}
  a~algoritmy odhadu rozdělení (EDA).
\item \textbf{Na velikosti populace záleží.} Populace musí obsahovat dostatek vzorků
  od každého stavebního bloku, jinak není z~čeho skládat; tomu se věnuje teorie
  \emph{population sizing}.
\item \textbf{Předčasná konvergence škodí.} Se ztrátou diverzity zmizí z~populace
  materiál, z~něhož by šly chybějící bloky objevit.
\end{itemize}

\subsubsection{Implicitní paralelismus a~víceruký bandita}

Populace vypadá zvenčí jako pouhých $n$ řetězců. Kombinatorika schémat ale ukázala, že
každý z~nich je zároveň svědectvím o~$2^m$ schématech, takže algoritmus s~jedním
vyhodnocením fitness posouvá odhad u~obrovského množství schémat najednou. Právě tomu se
říká implicitní paralelismus.

\begin{defbox}[Implicitní paralelismus (\emph{implicit parallelism})]
GA pracuje explicitně s~$n$ jedinci, ale implicitně přitom \uv{zpracovává}
mnohem více schémat: mezi $2^m$ a~$n \cdot 2^m$. Holland odhadl, že těch schémat,
se kterými algoritmus zachází \emph{užitečně} (rostou exponenciálně a~nejsou
příliš rozbíjena), je řádu $n^3$.
\end{defbox}

Bývá to komentováno jako jediný případ, kdy kombinatorická exploze pracuje
v~náš prospěch, a~ne proti nám.

\paragraph{Souvislost s~úlohou o~banditovi.} Hollandova původní motivace mířila na
\uv{adaptivní plán}, který hledá rovnováhu mezi explorací a~exploatací. Formálně se
tahle úvaha dá postavit takto:
\begin{itemize}
\item \textbf{Dvouruký bandita.} Máme $N$ mincí a~automat se dvěma pákami, jejichž
  střední výplaty $\mu_1, \mu_2$ ani rozptyly neznáme. Alokujeme
  $n$ mincí horší a~$N-n$ lepší páce. Analytické řešení říká, že optimální
  strategie přiděluje momentálně vítězné páce \emph{exponenciálně} více pokusů:
  $N - n^* = O(e^{c\, n^*})$.
\item \textbf{GA jako bandita.} Věta o~schématech tvrdí o~genetickém algoritmu totéž:
  úspěšnějším schématům přiděluje exponenciálně více \uv{slotů v~populaci}, a~řeší
  tedy dilema explorace a~exploatace (téměř) optimálně.
\item Původně se soudilo, že SGA hraje jedinou hru s~$3^m$ pákami. Ve
  skutečnosti hraje mnoho her paralelně: schémata spolu soutěží vždy o~konkrétní
  pevné pozice, takže schémata řádu $k$ hrají hru s~$2^k$ pákami.
\end{itemize}

\subsubsection{Kritika teorie schémat}

Věta o~schématech je pravdivá, jenže se z~ní dlouho vyvozovalo víc, než z~ní plyne.
Kritika míří na tři místa: na úlohy, kde stavební bloky vedou špatným směrem, na to,
co vlastně GA o~schématech měří, a~na samotný tvar věty.

\paragraph{Kdy GA selhává --- deceptivní úlohy}

Uvažujme funkci, v~níž jsou nadprůměrná schémata
\[
S_1 = (111{*}\ldots{*}), \qquad S_2 = ({*}\ldots{*}11),
\]
ale zároveň
\[
f(111{*}{*}{*}{*}{*}11) \;\ll\; f(000{*}{*}{*}{*}{*}00).
\]
Selekce vede algoritmus k~jedničkovým blokům, přestože optimum leží úplně jinde,
u~samých nul. Takovým úlohám se říká \textbf{deceptivní} (\emph{deceptive},
Goldberg 1987): dílčí stavební bloky nevedou po složení ke globálnímu optimu.
Deception bývá navrhována jako míra obtížnosti úlohy pro EA, i~když ne všichni
s~tím souhlasí.

Opačným testem jsou \textbf{Royal Road funkce} (Mitchell, Forrest, Holland). Fitness
je tu součtem příspěvků za kompletní \uv{bloky} jedniček, například osm bloků po
osmi bitech. Krajina je tedy ušitá přesně na míru hypotéze stavebních bloků a~GA by na
ní měl excelovat. Experiment dopadl opačně: jednoduchý GA prohrál s~náhodně mutujícím
horolezeckým algoritmem (RMHC). Na vině je \emph{hitchhiking} (\uv{stopování}): spolu
s~nalezeným dobrým blokem se selekcí šíří i~nahodilé \uv{parazitní} bity na ostatních
pozicích a~ničí diverzitu, kterou by algoritmus potřeboval k~nalezení zbývajících bloků.
Royal Road funkce jsou proto standardním empirickým protipříkladem k~naivnímu čtení BBH.

\paragraph{Statické průměrné fitness --- statická BBH}

Druhá výtka pochází od Grefenstetta (1991) a~týká se toho, co vlastně GA o~schématu ví.
Běžný výklad BBH tvrdí, že GA konverguje k~řešením s~dobrou \emph{skutečnou statickou}
průměrnou fitness schématu, tedy s~průměrem přes celý prostor. GA však odhaduje fitness
schématu jen \emph{ze vzorků, které má právě v~populaci}, a~ten odhad je vychýlený,
a~to hned ze dvou důvodů:

\begin{itemize}
\item \textbf{Kolaterální konvergence} (\emph{collateral convergence}).
  Jakmile GA někam konverguje, přestává vzorkovat schémata rovnoměrně.
  Je-li například schéma $111{*}\ldots{*}$ dobré, má po několika generacích
  skoro každý jedinec tento prefix. Pak je ale téměř každý vzorek schématu
  ${*}{*}{*}000\ldots$ zároveň vzorkem $111000\ldots$ a~fitness prvního
  schématu vyjde zcela zkresleně.
\item \textbf{Velký rozptyl fitness schématu.} Uvažujme funkci, pro kterou je
  $f(x) = 2$ pro $x \in 111{*}\ldots{*}$, $f(x)=1$ pro $x \in 0{*}\ldots{*}$
  a~$0$ jinak. Statická průměrná fitness schématu $1{*}\ldots{*}$ vychází $\approx 1/2$,
  zatímco $F(0{*}\ldots{*}) = 1$. GA však odhadne
  $F(1{*}\ldots{*}) \approx 2$, protože v~populaci převáží jedinci
  s~prefixem 111. Vysoký rozptyl tedy vede k~systematickému zkreslení a~k~tomu,
  že GA nevzorkuje schémata nezávisle.
\end{itemize}

\paragraph{Souhrn slabin}

Poslední skupina výhrad se netýká konkrétních úloh, ale samotného tvaru věty a~předpokladů,
za nichž byla odvozena:

\begin{itemize}
\item Věta je pouze \textbf{nerovnost} a~navíc \textbf{jednokrokový} výsledek. Iterovat
  ji přes více generací lze jen za předpokladu, že poměr $F(S,t)/\bar f(t)$ zůstává
  konstantní, což typicky neplatí: mění se dynamicky obě jeho složky.
\item Populace jsou \textbf{konečné a~malé}, takže se uplatní výběrové chyby;
  exponenciální růst je jen střední hodnota, ne záruka.
\item Věta zohledňuje jen \textbf{destruktivní} účinek operátorů a~pomíjí jejich
  \emph{konstruktivní} roli, přestože křížení může schéma i~vytvořit.
\item Soutěž schémat je \textbf{silně závislá}, zatímco páky bandity jsou nezávislé.
\item Analýza je \uv{on-line}: SGA maximalizuje průběžný výkon, hodí se tedy
  spíše pro adaptivní úlohy. Nejběžnější použití je přitom nalezení jednoho
  nejlepšího řešení, což je kritérium off-line.
\end{itemize}

Přes všechny výhrady zůstává TST užitečnou intuicí: \emph{co je krátké, jednoduché
a~dobré, to se šíří}.

\subsection{Pravděpodobnostní modely jednoduchého GA}
\label{sec:vose}


Tam, kde teorie schémat nabízí jen nerovnost pro jediný krok, chtěli mít lidé rovnost.
V~90.\ letech proto vznikly \emph{exaktní} modely SGA (Vose, Liepins, Nix, Whitley),
které přibližnou teorii schémat nahradily. Program byl jasný: přesně popsat, jak vypadá
populace, najít zobrazení, které jednu populaci převádí na následující, prozkoumat jeho
vlastnosti a~z~nich odvodit asymptotické chování SGA. Přístupy jsou dva a~liší se v~tom,
jak velkou populaci připustí. \textbf{Nekonečné populace} dávají deterministický
dynamický systém, \textbf{konečné populace} Markovův řetězec.

\subsubsection{Zjednodušený SGA}

Aby byl SGA analyzovatelný, ještě se zjednoduší. V~každém kroku se vyberou dva rodiče,
s~pravděpodobností $p_C$ se zkříží a~\textbf{jeden z~potomků se zahodí}; ten zbylý
zmutuje a~vloží se do nové populace. Zahození jednoho potomka není libovůle: díky němu
vzniká každý jedinec nové populace nezávisle na ostatních.

\subsubsection{Formalizace populace}
\label{ssec:vose}

Každý binární řetězec délky $l$ ztotožníme s~číslem
$i \in \{0,1,\dots,2^l-1\}$, například $00000111 \equiv 7$. Celou populaci v~čase $t$
pak popíšou dva vektory délky $2^l$:

\begin{defbox}[Vektor populace a~vektor selekce]
\begin{itemize}
\item $p(t) = (p_0(t),\dots,p_{2^l-1}(t))$, kde $p_i(t)$ je \emph{relativní
  četnost} řetězce $i$ v~populaci. Platí $p_i \ge 0$ a~$\sum_i p_i = 1$, takže
  $p(t)$ leží v~simplexu $\Lambda$.
\item $s(t) = (s_0(t),\dots,s_{2^l-1}(t))$, kde $s_i(t)$ je pravděpodobnost,
  že selekce vybere řetězec $i$.
\end{itemize}
Vektor $p$ popisuje \emph{složení} populace, vektor $s$ \emph{selekční
pravděpodobnosti}.
\end{defbox}

\subsubsection{Operátor $\mathcal{G}$}

Cílem je najít operátor $\mathcal{G}: \Lambda \to \Lambda$ takový, že
\[
p(t+1) = \mathcal{G}\big(p(t)\big),
\]
a~rozložit ho na dvě části, $\mathcal{G} = \mathcal{M} \circ \mathcal{F}$,
kde $\mathcal{F}$ odpovídá selekci (fitness) a~$\mathcal{M}$ tzv.\ míchání
(\emph{mixing}), tedy křížení a~mutaci dohromady. Jakmile takový operátor máme,
stačí ho iterovat z~počáteční populace a~známe úplné chování SGA.

\paragraph{Část $\mathcal{F}$ --- selekce.} Definujme diagonální matici
$\mathbf{F}$ s~$\mathbf{F}_{ii} = f(i)$ a~$\mathbf{F}_{ij} = 0$ pro $i \ne j$.
Fitness-proporcionální selekce je pak přesně
\[
s(t) \;=\; \frac{\mathbf{F}\,p(t)}{\sum_j \mathbf{F}_{jj}\,p_j(t)} .
\]
Čitatel $(\mathbf{F}p)_i = f(i)\,p_i$ je \uv{váha} řetězce $i$ v~populaci,
jmenovatel je normalizační konstanta, totiž průměrná fitness populace.

Vypustíme-li nyní křížení a~mutaci, je nová populace prostě $n$ nezávislých
tahů podle rozdělení $s(t)$, takže $\mathbb{E}[p(t+1)] = s(t)$. Dosazením
do definice selekce dostáváme $\mathbb{E}[s(t+1)] \propto \mathbf{F}\,s(t)$,
takže iterování $\mathcal{F}$ není nic jiného než (až na normalizaci) opakované
násobení maticí $\mathbf{F}$, tedy \emph{mocninná metoda}. Ta konverguje k~vlastnímu
vektoru příslušnému největší diagonální hodnotě, čili k~populaci složené
výhradně z~kopií nejlepšího jedince \emph{přítomného v~počáteční populaci}.
U~konečných populací způsobují statistické chyby odchylky od střední hodnoty, a~čím
větší populace, tím menší odchylka. Nekonečné populace jsou exaktní, byť nepraktické.

\textbf{Mini-příklad.} Pro $l = 2$ máme čtyři řetězce $00,01,10,11$, tedy
indexy $0,1,2,3$. Nechť $f = (1, 2, 3, 4)$ a~populace osmi jedinců obsahuje
$2\times 00$, $4\times 01$, $1\times 10$, $1\times 11$, tj.\
$p = (0{,}25;\,0{,}5;\,0{,}125;\,0{,}125)$. Pak
$\mathbf{F}p = (0{,}25;\,1{,}0;\,0{,}375;\,0{,}5)$, součet je $2{,}125$
(= průměrná fitness populace) a
\[
s \;=\; (0{,}118;\; 0{,}471;\; 0{,}176;\; 0{,}235).
\]
Podíl nejlepšího řetězce $11$ vzrostl z~$12{,}5\,\%$ na $23{,}5\,\%$ a~podíl
nejhoršího $00$ klesl z~$25\,\%$ na $11{,}8\,\%$, a~to přesně o~faktor
$f(i)/\bar f$, jak víme z~věty o~schématech.

\paragraph{Část $\mathcal{M}$ --- míchání.} Zavedeme funkci
\[
r(i,j,k) \;=\; \Pr[\text{potomek } k \text{ vznikne z~rodičů } i, j],
\]
která v~sobě zahrnuje jak křížení (přes všechny body řezu), tak mutaci.
Pak
\[
\mathbb{E}\big[p_k(t+1)\big] \;=\; \sum_{i}\sum_{j} s_i(t)\, s_j(t)\, r(i,j,k).
\]
Míchání je tedy \emph{kvadratický} operátor na simplexu. Odvodit $r(i,j,k)$ je
technicky náročné, ale možné. Klíčové zjednodušení plyne z~toho, že jak
jednobodové křížení, tak bitová mutace \uv{nevidí} absolutní hodnoty bitů,
nýbrž jen jejich shody a~neshody. Formálně to znamená, že pro libovolné $k$ platí
\[
r(i,j,k) \;=\; r(i \oplus k,\; j \oplus k,\; 0),
\]
kde $\oplus$ je bitový XOR. Stačí tedy znát jedinou \emph{míchací matici}
$\mathbf{M}$ s~prvky $\mathbf{M}_{ij} = r(i,j,0)$ (rozměru $2^l \times 2^l$)
a~zbytek se dopočte permutacemi $\sigma_k : x \mapsto x \oplus k$:
$\mathcal{M}(s)_k = (\sigma_k s)^{\!\top}\, \mathbf{M}\, (\sigma_k s)$.
V~tom je hlavní technický přínos Voseho formalismu: úloha se smrskne
z~$2^{3l}$ čísel na $2^{2l}$.

\subsubsection{Výsledky pro nekonečné populace}

Co z~modelu plyne, se dá shrnout do čtyř bodů:
\begin{itemize}
\item SGA je \textbf{diskrétní dynamický systém} na simplexu a~posloupnost
  $p(0), p(1), \dots$ je trajektorií tohoto systému.
\item \textbf{Pevné body $\mathcal{F}$} jsou populace tvořené pouze kopiemi
  nejlepšího jedince z~počáteční populace; selekce tedy \uv{zaostřuje}.
\item \textbf{Pevný bod $\mathcal{M}$} je jediný, totiž rovnoměrné rozdělení.
  Míchání naopak \uv{rozostřuje} a~maximalizuje entropii.
\item Pevné body složeného $\mathcal{G}$ obecně známy nejsou a~jejich hledání
  je hlavní otevřený problém. Chování je kombinací zaostřování a~rozostřování
  a~projevuje se jako \textbf{přerušované rovnováhy}
  (\emph{punctuated ekvilibria}), tedy dlouhá období metastability střídaná
  rychlými přechody. Tentýž jev je znám i~z~biologie.
\end{itemize}

\subsubsection{Konečné populace: Markovův řetězec}

Skutečný běh SGA má populaci konečnou, takže se v~něm uplatní náhoda a~deterministická
trajektorie nestačí. Zato je celý proces bezpaměťový: nová populace závisí výhradně na té
současné, a~to je přesně situace pro Markovův řetězec.

\begin{defbox}[Markovův řetězec]
Stochastický proces v~diskrétním čase se stavy, kde pravděpodobnost přechodu
do dalšího stavu závisí \emph{jen} na aktuálním stavu, nikoli na historii.
Úplným popisem takového procesu je matice přechodových pravděpodobností.
\end{defbox}

Modelujeme SGA s~populací velikosti $n$ nad řetězci délky $l$. Potřebujeme tři věci:
\begin{itemize}
\item \textbf{Stavy} jsou jednotlivé možné populace, tedy multimnožiny $n$ řetězců.
  Jejich počet je
  \[
  N \;=\; \binom{n + 2^l - 1}{\,2^l - 1\,},
  \]
  což je počet multimnožin velikosti $n$ z~$2^l$ typů.
\item \textbf{Matice $\mathbf{Z}$}, jejíž sloupce odpovídají populacím a~jejíž prvek
  $\mathbf{Z}(y,i)$ udává počet výskytů jedince $y$ v~populaci $i$. Sloupce
  matice $\mathbf{Z}$ jsou tedy přímo stavy řetězce.
\item \textbf{Matice přechodů $\mathbf{Q}$} rozměru $N \times N$, kterou lze
  explicitně odvodit z~$\mathcal{G}$, protože přechod není nic jiného než
  multinomické losování $n$ jedinců podle rozdělení $\mathcal{G}(p)$:
  \[
  \mathbf{Q}_{i,j} \;=\; n! \prod_{y=0}^{2^l-1}
  \frac{\big(\mathcal{G}(p_i)_y\big)^{\mathbf{Z}(y,j)}}{\mathbf{Z}(y,j)!}.
  \]
\end{itemize}

\textbf{Problém velikosti.} Už pro $n = l = 8$ má $\mathbf{Q}$ řádově
$10^{29}$ prvků. Model je tedy exaktní, ale k~počítání prakticky nepoužitelný;
jeho hodnota je teoretická.

\paragraph{Kvalitativní důsledky.}
\begin{itemize}
\item Je-li $p_M > 0$, je řetězec \textbf{ireducibilní a~aperiodický}, protože
  z~libovolné populace lze mutacemi dosáhnout libovolné jiné. Existuje tedy
  jediné stacionární rozdělení a~SGA navštíví každou populaci, tu s~globálním
  optimem nevyjímaje, nekonečně mnohokrát s~pravděpodobností 1.
\item Bez elitismu ale algoritmus ke~globálnímu optimu \textbf{nekonverguje}:
  optimum se objeví a~zase zmizí. \emph{Elitistický} GA, který si nejlepší nalezené
  řešení uchovává, konverguje k~globálnímu optimu s~pravděpodobností 1 při
  $t \to \infty$. To je standardní konvergenční výsledek, neříká však nic
  o~\emph{rychlosti}.
\item Mezi modelem nekonečné a~konečné populace existuje korespondence:
  pro $n \to \infty$ se chování konečného modelu blíží deterministické
  trajektorii $\mathcal{G}$.
\end{itemize}

\subsubsection{Algoritmy odhadu rozdělení (EDA)}
\label{ssec:eda}

\noindent\mimoq{ve slidech EVA není; je to ale přímé pokračování Voseho pohledu}

\begin{defbox}[EDA --- obecné schéma]
Voseho model ukazuje, že SGA je ve skutečnosti stroj na transformaci pravděpodobnostního
rozdělení nad prostorem řetězců. \emph{Algoritmy odhadu rozdělení} berou toto pozorování
doslova: zahodí křížení i~mutaci a~místo nich udržují \textbf{explicitní model}
$p_\theta(x)$. Jedna iterace vypadá tak, že se z~$p_\theta$ navzorkuje populace, vybere
se z~ní lepší část, na té se model \emph{přeučí}, a~znovu. Model $\theta$ je jediná
informace, která se přenáší mezi generacemi.

Jednotlivé algoritmy se liší tím, jak bohaté závislosti model zachytí: \textbf{žádné}
(UMDA, PBIL), \textbf{párové} (MIMIC), \textbf{obecné} (BOA s~Bayesovskou sítí) nebo
\textbf{spojité} (CMA-ES). PBIL například posouvá vektor pravděpodobností k~nejlepšímu
jedinci: $p_i \gets (1-\alpha)\,p_i + \alpha\, x_i^{\mathrm{best}}$.
\end{defbox}

Univariátní model selže, jakmile mezi geny existuje \emph{epistáze}, a~právě proto
vznikl BOA, který se strukturu závislostí učí. Tím EDA řeší problém, kvůli němuž Holland
navrhoval inverzi: kde mají spolupracující geny ležet, si algoritmus zjistí sám.

\subsection{Shrnutí ke~zkoušce}

\begin{itemize}
\item Schéma je slovo nad abecedou $\{0,1,*\}$. Řád $o(S)$ je počet pevných pozic,
  definující délka $d(S)$ je vzdálenost první a~poslední pevné pozice a~$F(S)$ je
  průměrná fitness reprezentantů \emph{v~dané populaci}. Celkem existuje $3^m$ schémat
  a~jeden jedinec je reprezentantem $2^m$ z~nich.
\item Věta o~schématech zní
  $\mathbb{E}[C(S,t+1)] \ge C(S,t)\frac{F(S)}{\bar f}
  [1 - p_C \frac{d(S)}{m-1} - p_M o(S)]$ a~odvozuje se ve třech krocích:
  selekce dává exponenciální růst, křížení přispívá členem $d(S)/(m-1)$
  a~mutace členem $(1-p_M)^{o(S)}$.
\item Hypotéza stavebních bloků říká, že GA skládá krátká, nízkého řádu
  a~nadprůměrná schémata, tzv.\ stavební bloky. Z~toho plyne, že záleží na kódování
  i~na velikosti populace a~že škodí předčasná konvergence.
\item Implicitní paralelismus znamená, že se užitečně zpracovává $\sim n^3$ schémat.
  Analogie s~banditou je v~tom, že schémata řádu $k$ soupeří o~tytéž $k$ pevných pozic,
  hrají tedy hru s~$2^k$ pákami; analytické řešení přiděluje vítězi exponenciálně
  více pokusů, $N - n^* = O(e^{c n^*})$.
\item Kritika má šest bodů: věta je nerovnost pro jediný krok, populace jsou konečné,
  existují deceptivní úlohy, uplatňuje se kolaterální konvergence a~velký rozptyl
  fitness schématu, a~model ignoruje konstruktivní efekty i~vzájemnou závislost schémat.
\item Voseho model ztotožní řetězce s~čísly $0..2^l-1$. Populaci pak popisuje
  vektor relativních četností $p(t)$ ležící v~simplexu spolu s~vektorem selekčních
  pravděpodobností $s(t)$.
\item Hledá se $\mathcal{G} = \mathcal{M} \circ \mathcal{F}$ s~$p(t+1)=\mathcal{G}(p(t))$.
  Selekční část $\mathcal{F}$ je dána diagonální maticí fitness,
  $s = \mathbf{F}p / \sum_j \mathbf{F}_{jj}p_j$, kdežto míchací část $\mathcal{M}$
  (křížení a~mutace) je kvadratický operátor
  $\mathbb{E}[p_k'] = \sum_{i,j} s_i s_j\, r(i,j,k)$. Díky symetrii
  $r(i,j,k) = r(i\oplus k, j \oplus k, 0)$ stačí jediná míchací matice.
\item Pevné body $\mathcal{F}$ jsou populace kopií nejlepšího jedince, pevným bodem
  $\mathcal{M}$ je rovnoměrné rozdělení. Jejich kombinace vede k~přerušovaným
  rovnováhám a~pevné body $\mathcal{G}$ známy nejsou.
\item Pro konečné populace se staví Markovův řetězec, jehož stavy jsou populace.
  Stavů je $N = \binom{n+2^l-1}{2^l-1}$ a~matice $\mathbf{Q}$ je sice exaktní, ale
  astronomicky velká. Při $p_M>0$ je řetězec ireducibilní a~elitistický GA konverguje
  k~optimu s~pravděpodobností 1.
\item Nekonečné populace dávají deterministický model, který je exaktní, ale
  nepraktický; konečné populace model stochastický, tedy realistický, zato
  nezvládnutelný. Limitní přechod obě varianty spojuje.
\item \textbf{EDA} berou pohled \uv{GA je transformace rozdělení} doslova: místo
  operátorů se učí model $p_\theta$. UMDA a~PBIL pracují s~nezávislými marginálami,
  MIMIC a~BMDA s~párovými závislostmi, BOA s~Bayesovskou sítí a~CMA-ES s~gaussovkou.
  Cyklus je pokaždé stejný: navzorkuj, vyber lepší část, přeuč model. PBIL konkrétně
  počítá $p_i \gets (1-\alpha)p_i + \alpha x_i^{\mathrm{best}}$.
\end{itemize}

%=====================================================================
\section{Evoluční strategie, diferenciální evoluce, koevoluce, otevřená evoluce}
\label{sec:es}
%=====================================================================

Kapitola spojuje čtyři témata, jejichž společným jmenovatelem je, že jdou \emph{za} klasický
genetický algoritmus. První dvě, evoluční strategie a~diferenciální evoluce, jsou metody pro
spojitou optimalizaci: pracují s~vektory reálných čísel a~těžiště hledání přesouvají z~křížení
na mutaci.

Zbylá dvě témata mění samotnou fitness. U~koevoluce přestává být fitness absolutní veličinou
a~závisí na ostatních jedincích, u~otevřené evoluce není fitness vůbec pevně dána.

\subsection{Evoluční strategie}
\label{ssec:es}
\subsubsection{Motivace a~základní rysy}

Evoluční strategie (\emph{evolution strategies}, ES) navrhli Rechenberg a~Schwefel v~Berlíně
v~60.\ letech, a~to kvůli velmi konkrétní potřebě: optimalizovat reálné parametry v~obtížných
technických úlohách, jako je tvar trysky nebo ohyb potrubí. Úloha byla od začátku spojitá
a~té je přizpůsobená celá stavba algoritmu.

Od genetického algoritmu se ES liší především ve čtyřech rysech.

\begin{itemize}
\item Jedinec je vektor reálných čísel a~\textbf{mutace je hlavní operátor} --- právě ona
  nese tíhu prohledávání prostoru.
\item Kromě \emph{genetických parametrů}, tedy vlastního řešení, nese jedinec
  i~\emph{strategické, endogenní parametry} (\emph{endogenous parameters}), které řídí
  samotnou evoluci; typicky to jsou směrodatné odchylky mutací. Vyvíjí se tak nejen řešení,
  ale i~způsob, jakým se hledá, a~odtud pochází slogan \uv{evoluce evoluce}.
\item Rodičovská selekce je \textbf{náhodná} a~na fitness vůbec nebere ohled, zatímco
  environmentální selekce je \textbf{deterministická}: přežívá $\mu$ nejlepších. Veškerý
  selekční tlak je tedy soustředěn do jediného místa cyklu.
\item Nejmodernějším a~nejúspěšnějším zástupcem celé rodiny je CMA-ES, kterému je níže
  věnován samostatný oddíl.
\end{itemize}

\begin{defbox}[Notace ES]
Symbol $\mu$ značí velikost populace, tedy počet rodičů, $\lambda$ počet vytvořených potomků
a~$\rho$ počet rodičů, kteří se podílejí na jednom potomkovi.
\begin{itemize}
\item V~$(\mu + \lambda)$-ES se nová populace vybírá z~\emph{$\mu$ rodičů i~$\lambda$
  potomků}; schéma je tedy elitistické a~dosud nejlepší jedinec z~populace nikdy nezmizí.
\item V~$(\mu, \lambda)$-ES se vybírá \emph{jen z~$\lambda$ potomků} a~rodiče vždy vymřou.
  Schéma je neelitistické a~nutně vyžaduje $\lambda > \mu$, jinak by nebylo z~čeho vybírat.
\item Zápis $(\mu/\rho \,\overset{+}{,}\, \lambda)$-ES je explicitní varianta obojího, která
  navíc uvádí, že potomek vzniká rekombinací $\rho$ rodičů.
\end{itemize}
\end{defbox}

Volba mezi čárkou a~plusem není kosmetická. Obě schémata se liší téměř ve všem podstatném,
od schopnosti uniknout z~lokálního optima až po to, jestli se samoadaptace vůbec chytne:

\begin{center}
\begin{tabular}{@{}p{4cm}p{5.5cm}p{5.5cm}@{}}
\toprule
 & $(\mu,\lambda)$ & $(\mu+\lambda)$ \\
\midrule
pool pro selekci & jen potomci & rodiče + potomci \\
elitismus & ne (fitness může klesnout) & ano \\
únik z~lokálních optim & lepší & horší \\
konvergence & pomalejší & rychlejší \\
doporučení & spojité domény, zašuměná či dynamická fitness & diskrétní domény,
  drahá fitness \\
samoadaptace $\sigma$ & funguje dobře (špatné $\sigma$ vymře) & může uváznout
  (starý rodič s~malým $\sigma$ přežívá) \\
poměr & typicky $\lambda/\mu \approx 7$ & i~$\mu \ge \lambda$; $\lambda = 1$ =
  steady-state ES \\
\bottomrule
\end{tabular}
\end{center}

\subsubsection{$(1+1)$-ES a~pravidlo 1/5}

Nejjednodušší evoluční strategie si vystačí s~jediným jedincem. V~každé iteraci z~něj
gaussovská mutace vytvoří jediného potomka a~ten se přijme, jen pokud je lepší než rodič.
Zůstává tím jediná otevřená otázka: jak velké kroky má mutace dělat. Odpověď na ni dává
pravidlo 1/5.

\begin{algorithm}[!ht]
\caption{$(1+1)$-ES s~pravidlem 1/5 (minimalizace $f:\mathbb{R}^n \to \mathbb{R}$)}
\begin{algorithmic}[1]
\State inicializuj $\vec x(0)$ náhodně, $\sigma > 0$, $t \gets 0$
\Repeat
  \State $\vec y(t) \gets \vec x(t) + \sigma \cdot N(\vec 0, \mathbf{I})$
  \If{$f(\vec y(t)) < f(\vec x(t))$} $\vec x(t+1) \gets \vec y(t)$
  \Else\ $\vec x(t+1) \gets \vec x(t)$
  \EndIf
  \State sleduj $p$ = podíl úspěšných mutací za posledních $k$ iterací
  \State každých $k$ iterací uprav $\sigma$:
  \State \quad \textbf{if} $p > 1/5$ \textbf{then} $\sigma \gets \sigma / c$
    \Comment{daří se: dělej větší kroky, více explorace}
  \State \quad \textbf{if} $p < 1/5$ \textbf{then} $\sigma \gets \sigma \cdot c$
    \Comment{nedaří se: zmenši kroky, více exploatace}
  \State \quad \textbf{if} $p = 1/5$ \textbf{then} $\sigma$ beze změny
    \Comment{typicky $0{,}8 \le c \le 1$}
  \State $t \gets t+1$
\Until{$\vec x(t)$ je dost dobré}
\end{algorithmic}
\end{algorithm}

\begin{defbox}[Pravidlo 1/5 úspěchu (\emph{1/5 success rule})]
Rechenberg odvodil z~analýzy dvou modelových funkcí, kulové a~koridorové, heuristiku, podle
níž je \emph{optimální poměr úspěšných mutací přibližně $1/5$.} Je-li podíl úspěchů vyšší,
postupujeme příliš opatrně: krok je malý a~$\sigma$ se zvětší. Je-li nižší, je krok naopak
příliš velký a~$\sigma$ se zmenší.

Řízení parametru je zde \emph{adaptivní}, nikoli samoadaptivní. Algoritmus sice využívá
zpětnou vazbu z~vlastního běhu, ale $\sigma$ není součástí genomu a~selekce na ně nepůsobí.
\end{defbox}

\subsubsection{Samoadaptace strategických parametrů}

Pravidlo 1/5 řídí jediný parametr zvenčí a~pro celý prostor stejně. Samoadaptace jde ještě o~krok
dál: nastavení mutace svěří evoluci samotné, takže se velikost i~tvar kroku vyladí souběžně
s~řešením.

\begin{defbox}[Samoadaptace (\emph{self-adaptation})]
Strategické parametry, tedy odchylky $\sigma$ a~případně úhly natočení, jsou
\textbf{součástí genomu} a~podléhají mutaci i~selekci společně s~řešením. Selekce na ně
působí jen nepřímo: jedinci s~vhodným $\sigma$ produkují lepší potomky, a~jejich $\sigma$
se tím šíří populací.

Na pořadí obou mutací přitom záleží. \textbf{Nejprve se mutují strategické parametry, teprve
pak se s~nimi mutují genetické parametry}; při opačném pořadí by selekce hodnotila staré
$\sigma$.
\end{defbox}

Jedinec má tvar $C = [\vec x, \vec \sigma]$, kde $\vec x \in \mathbb{R}^n$ je vlastní řešení
a~$\vec\sigma$ má $k$ složek. Počet složek rozhoduje o~tom, jaké tvary dokáže mutace
v~prostoru opsat:

\begin{center}
\begin{tabular}{@{}clp{8.2cm}@{}}
\toprule
$k$ & název & geometrie mutace \\
\midrule
$1$ & jedno globální $\sigma$ & izotropní --- vrstevnice hustoty jsou koule \\
$n$ & nekorelované mutace & elipsoid s~osami rovnoběžnými se souřadnicemi \\
$n(n+1)/2$ & korelované mutace & obecný elipsoid ($n$ odchylek $+$ $n(n-1)/2$
  úhlů natočení $\alpha_{ij}$), odpovídá plné kovarianční matici \\
\bottomrule
\end{tabular}
\end{center}

V~nejobecnějším případě je kovarianční matice $\mathbf{C}$ součinem rotační matice a~její
prvky vycházejí jako $c_{ii} = \sigma_i^2$ a~$c_{ij} = \tfrac12(\sigma_i^2 -
\sigma_j^2)\tan(2\alpha_{ij})$.

\paragraph{Standardní pravidla mutace (Schwefel).} Konkrétní podobu mutace obou skupin
parametrů zavedl Schwefel a~používá se dodnes:
\[
\sigma_i' = \sigma_i \cdot \exp\big(\tau' N(0,1) + \tau N_i(0,1)\big),
\qquad
x_i' = x_i + \sigma_i' \cdot N_i(0,1),
\]
kde $\tau' \propto 1/\sqrt{2n}$ je společná porucha pro celého jedince a~$\tau \propto
1/\sqrt{2\sqrt{n}}$ individuální porucha jednotlivé složky. Lognormální tvar plní dva úkoly
najednou: udržuje $\sigma$ kladné a~činí zvětšení a~zmenšení symetrickými. Aby $\sigma$
nezdegenerovalo na nulu a~evoluce se nezastavila, zavádí se dolní mez $\sigma_{\min}$. Úhly
natočení se na rozdíl od odchylek mutují aditivně, tedy
$\alpha_{ij}' = \alpha_{ij} + \beta \cdot N(0,1)$.

\subsubsection{Rekombinace v~ES}

Rekombinace v~ES má oproti křížení jednu zvláštnost: z~$\rho$ rodičů vzniká \emph{jeden}
potomek. Podle počtu zapojených rodičů se rozlišuje rekombinace \textbf{lokální}, kde
$\rho = 2$, a~\textbf{globální}, kde $\rho = \mu$ a~potomek tedy čerpá z~celé populace.
Druhé dělení se týká toho, jak se rodičovské hodnoty skládají dohromady.

\begin{itemize}
\item \textbf{Diskrétní (dominantní) rekombinace} vybírá hodnotu na každé pozici náhodně od
  jednoho z~rodičů, takže potomek je mozaikou nezměněných rodičovských hodnot.
\item \textbf{Intermediární (aritmetická) rekombinace} bere na každé pozici průměr hodnot
  rodičů; potomek tedy nepřebírá beze změny ani jednu z~nich.
\end{itemize}
V~praxi se obě varianty kombinují. Genetické parametry se rekombinují diskrétně, protože to
zachovává rozmanitost, kdežto strategické parametry intermediárně, protože průměrování
stabilizuje adaptaci $\sigma$.

\begin{algorithm}[!ht]
\caption{Obecný cyklus ES}
\begin{algorithmic}[1]
\State $t \gets 0$; inicializuj náhodně populaci $P_t$ o~$\mu$ jedincích $[\vec x, \vec\sigma]$
\State ohodnoť jedince v~$P_t$
\While{není splněna ukončovací podmínka}
  \State $C_{t+1} \gets \emptyset$
  \For{$\lambda$ krát}
    \State vyber náhodně $\rho$ rodičů z~$P_t$
    \State rekombinuj je do jednoho potomka $[\vec x, \vec\sigma]$
    \State mutuj strategické parametry $\vec\sigma$
    \State mutuj genetické parametry $\vec x$ pomocí nových $\vec\sigma$
    \State ohodnoť potomka a~vlož do $C_{t+1}$
  \EndFor
  \State $(\mu,\lambda)$: $P_{t+1} \gets$ $\mu$ nejlepších z~$C_{t+1}$
  \State $(\mu+\lambda)$: $P_{t+1} \gets$ $\mu$ nejlepších z~$P_t \cup C_{t+1}$
  \State $t \gets t+1$
\EndWhile
\end{algorithmic}
\end{algorithm}

\subsubsection{CMA-ES}

\noindent\mimoq{slidy uvádějí jednou větou jako \uv{dnes nejúspěšnější a~nejsložitější} ES}

\begin{defbox}[CMA-ES --- princip]
\emph{Covariance Matrix Adaptation ES} (Hansen, Ostermeier) je dnes referenční algoritmus
pro spojitou black-box optimalizaci. Samoadaptaci jednotlivých $\sigma$ v~genomu nahrazuje
jiným nápadem: udržuje \textbf{jediné rozdělení pro celou populaci}
$\mathcal{N}(\vec m, \sigma^2 \mathbf{C})$ a~jeho kovarianční matici se učí z~historie
úspěšných kroků.

Jedna iterace vypadá takto. Navzorkuje se $\lambda$ bodů a~vybere se $\mu$ nejlepších
\emph{podle pořadí}. Střed $\vec m$ se posune do jejich váženého průměru a~matice
$\mathbf{C}$ spolu s~délkou kroku $\sigma$ se aktualizují podle \emph{evolučních cest}, tedy
exponenciálně vyhlazených záznamů o~tom, kudy se střed pohyboval. Míří-li po sobě jdoucí
kroky souhlasným směrem, $\sigma$ se zvětší; jsou-li protichůdné, zmenší se.
\end{defbox}

Výsledkem je algoritmus, který se sám přizpůsobí anizotropii a~korelacím úlohy; v~podstatě
se bez gradientu učí aproximaci inverzní Hessovy matice. K~tomu je invariantní vůči
monotónním transformacím fitness a~má velmi málo laditelných parametrů, takže se s~ním dá
pracovat i~na úloze, o~které se nic neví.

\subsection{Diferenciální evoluce}
\label{sec:de}
\subsubsection{Motivace}

Diferenciální evoluce (\emph{differential evolution}, DE; Storn a~Price, 1997) je jednoduchý
a~překvapivě silný algoritmus pro black-box optimalizaci funkcí
$f: \mathbb{R}^D \to \mathbb{R}$. Klíčová myšlenka je \textbf{geometrická}: velikost ani směr
mutačního kroku se neodhadují z~parametru $\sigma$, nýbrž se přebírají
z~\emph{rozdílů dvojic jedinců v~populaci}.

Populace tak sama kóduje měřítko a~orientaci úlohy. Na počátku je rozptýlená a~kroky jsou
proto velké; jakmile začne konvergovat, rozdíly se zmenší a~mutace přejde k~jemnému doladění.
Algoritmus tedy dostává samoadaptaci zdarma, aniž by se cokoli muselo evolvovat nebo ladit.

Aplikace sahají od strojového učení až po plánování optimálních trajektorií kosmických sond
v~Evropské kosmické agentuře.

\subsubsection{Algoritmus DE/rand/1/bin}

Populace je $N$ vektorů $\vec x_{1},\dots,\vec x_{N} \in \mathbb{R}^D$. Rodičovská selekce
je triviální, protože rodičem se postupně stává \textbf{každý} jedinec $\vec x_i$; pro
každého z~nich proběhnou tři kroky.

\begin{enumerate}
\item \textbf{Mutace.} Vyberou se tři navzájem různí jedinci
  $\vec x_{a}, \vec x_{b}, \vec x_{c}$, všichni odlišní od $\vec x_i$, a~vznikne z~nich
  \emph{dárce} (\emph{donor vector})
  \[
  \vec v_i \;=\; \vec x_{a} \;+\; F\,(\vec x_{b} - \vec x_{c}),
  \qquad F \in [0,2],\ \text{typicky } F \approx 0{,}8 .
  \]
  Rozdíl $\vec x_{b} - \vec x_{c}$ dodává směr i~délku kroku, koeficient $F$ jej škáluje.
\item \textbf{Křížení (binomické, uniformní \uv{s~pojistkou}).}
  Z~dárce a~rodiče se poskládá \emph{zkušební vektor} (\emph{trial vector}) $\vec u_i$:
  \[
  u_{j,i} =
  \begin{cases}
  v_{j,i} & \text{pokud } \mathrm{rand}_j \le C \text{ nebo } j = I_{\mathrm{rand}},\\[2pt]
  x_{j,i} & \text{jinak},
  \end{cases}
  \]
  kde $\mathrm{rand}_j \sim U(0,1)$, $C \in [0,1]$ je práh křížení a~$I_{\mathrm{rand}}$
  náhodný index z~$\{1,\dots,D\}$. Podmínka $j = I_{\mathrm{rand}}$ je zmíněná
  \uv{pojistka}: zaručuje, že alespoň jedna složka pochází od dárce, takže potomek nikdy
  není identickou kopií rodiče.
\item \textbf{Environmentální selekce (souboj 1:1).} Zkušební vektor se utká přímo se svým
  rodičem a~lepší z~dvojice postupuje dál:
  \[
  \vec x_i(t+1) =
  \begin{cases}
  \vec u_i & \text{pokud } f(\vec u_i) \le f(\vec x_i(t)),\\
  \vec x_i(t) & \text{jinak.}
  \end{cases}
  \]
  Selekce je tedy elitistická na úrovni jednotlivce a~populace se nikdy nezhorší.
\end{enumerate}

\begin{algorithm}[!ht]
\caption{DE/rand/1/bin}
\begin{algorithmic}[1]
\Require velikost populace $N$, měřítko $F$, práh křížení $C$
\State inicializuj populaci $\{\vec x_i\}_{i=1}^{N}$ náhodně v~mezích úlohy
\While{není splněna ukončovací podmínka}
  \For{$i \gets 1$ \textbf{to} $N$}
    \State vyber náhodně různé indexy $a,b,c \ne i$
    \State $\vec v \gets \vec x_a + F(\vec x_b - \vec x_c)$
    \State $I_{\mathrm{rand}} \gets$ náhodný index z~$\{1,\dots,D\}$
    \For{$j \gets 1$ \textbf{to} $D$}
      \State $u_j \gets v_j$ pokud $\mathrm{rand}() \le C$ nebo $j = I_{\mathrm{rand}}$, jinak $u_j \gets x_{j,i}$
    \EndFor
    \If{$f(\vec u) \le f(\vec x_i)$} $\vec x_i' \gets \vec u$ \Else\ $\vec x_i' \gets \vec x_i$ \EndIf
  \EndFor
  \State $\{\vec x_i\} \gets \{\vec x_i'\}$
\EndWhile
\end{algorithmic}
\end{algorithm}

\subsubsection{Číselný příklad jednoho kroku}

Celý cyklus se dá spočítat na papíře. Nechť $D = 3$ a~minimalizuje se
$f(\vec x) = \sum_j x_j^2$ s~parametry $F = 0{,}8$ a~$C = 0{,}9$. Aktuálním jedincem je
$\vec x_i = (4;\, -3;\, 2)$ s~hodnotou $f(\vec x_i) = 16+9+4 = 29$ a~náhodně vybráni byli
jedinci
\[
\vec x_a = (2;\, -1;\, 0{,}5), \quad
\vec x_b = (1;\, 0;\, -1), \quad
\vec x_c = (-1;\, 2;\, 1).
\]

\textbf{Mutace:} nejprve se spočítá rozdílový vektor a~z~něj dárce.
\[
\vec x_b - \vec x_c = (2;\, -2;\, -2), \qquad
\vec v = (2;\,-1;\,0{,}5) + 0{,}8\,(2;\,-2;\,-2) = (3{,}6;\; -2{,}6;\; -1{,}1).
\]

\textbf{Křížení:} nechť $I_{\mathrm{rand}} = 1$ a~náhodná čísla vyšla
$\mathrm{rand} = (0{,}40;\ 0{,}95;\ 0{,}20)$. Složku po složce potom platí:
\begin{itemize}
\item pro $j=1$ je $0{,}40 \le 0{,}9$, a~navíc jde o~pojistkový index
  $j = I_{\mathrm{rand}}$, takže $u_1 = v_1 = 3{,}6$;
\item pro $j=2$ je $0{,}95 > 0{,}9$ a~zároveň $j \ne I_{\mathrm{rand}}$, hodnota tedy
  zůstává rodičovská, $u_2 = x_{2,i} = -3$;
\item pro $j=3$ je $0{,}20 \le 0{,}9$, a~proto $u_3 = v_3 = -1{,}1$.
\end{itemize}
Zkušební vektor je $\vec u = (3{,}6;\, -3;\, -1{,}1)$.

\textbf{Selekce:} platí $f(\vec u) = 12{,}96 + 9 + 1{,}21 = 23{,}17 < 29 = f(\vec x_i)$,
takže $\vec u$ nahrazuje $\vec x_i$ v~nové populaci.

\subsubsection{Varianty mutace}

Varianty DE se označují zápisem \texttt{DE/$x$/$y$/$z$}, v~němž $x$ říká, jak se volí
základní vektor, $y$ udává počet rozdílových vektorů a~$z$ určuje typ křížení: \texttt{bin}
znamená binomické uniformní křížení, \texttt{exp} exponenciální, které přebírá souvislý
úsek složek.

\begin{center}
\begin{tabular}{@{}ll@{}}
\toprule
\textbf{varianta} & \textbf{dárce $\vec v$} \\
\midrule
DE/rand/1 & $\vec x_a + F(\vec x_b - \vec x_c)$ \\
DE/rand/2 & $\vec x_a + F(\vec x_b - \vec x_c + \vec x_d - \vec x_e)$ \\
DE/best/1 & $\vec x_{\mathrm{best}} + F(\vec x_b - \vec x_c)$ \\
DE/best/2 & $\vec x_{\mathrm{best}} + F(\vec x_b - \vec x_c + \vec x_d - \vec x_e)$ \\
DE/current-to-best/1 & $\vec x_i + F(\vec x_{\mathrm{best}} - \vec x_i) + F(\vec x_b - \vec x_c)$ \\
\bottomrule
\end{tabular}
\end{center}

Volba základního vektoru rozhoduje o~kompromisu mezi rychlostí a~robustností. Varianty
s~\texttt{best} konvergují rychleji, ale snáze uvíznou v~lokálním optimu, kdežto \texttt{rand}
je robustnější. Moderní varianty jDE, SaDE, JADE, SHADE a~L-SHADE adaptují $F$ a~$C$ za běhu
a~patří ke špičce v~soutěžích CEC.

\subsubsection{Srovnání DE s~ES a~GA}

Tři algoritmy je nejnázornější postavit vedle sebe podle toho, odkud berou velikost mutačního
kroku a~kde v~nich působí selekční tlak.

\begin{center}
\begin{tabular}{@{}p{3.2cm}p{3.7cm}p{3.7cm}p{3.7cm}@{}}
\toprule
 & \textbf{GA (reálný)} & \textbf{ES} & \textbf{DE} \\
\midrule
zdroj kroku mutace & pevné $\sigma$ & evolvované $\sigma$ & rozdíly jedinců \\
rodičovská selekce & podle fitness & náhodná & každý jedinec jednou \\
env.\ selekce & generační / elitismus & $\mu$ nejlepších & souboj rodič vs.\ potomek \\
hlavní operátor & křížení & mutace & diferenční mutace \\
adaptace měřítka & žádná & explicitní (v~genomu) & implicitní (z~populace) \\
\bottomrule
\end{tabular}
\end{center}

\subsection{Koevoluce}
\label{sec:koevoluce}
\subsubsection{Kontextová fitness}

\begin{defbox}[Koevoluce (\emph{coevolution})]
Evoluční algoritmus, v~němž fitness jedince \textbf{není absolutní}, ale závisí na ostatních
jedincích, ať už na jeho vlastní populaci, nebo na jiné, souběžně se vyvíjející populaci.
Fitness je tedy \emph{kontextová} a~fitness krajina \textbf{dynamická}: mění se tím, jak se
mění protivníci či partneři.
\end{defbox}

Důvod, proč se ke koevoluci sahá, je zpravidla praktický. U~mnoha úloh neumíme napsat
absolutní účelovou funkci, protože na otázku \uv{jak dobrá je strategie pro dámu?} se přímo
odpovědět nedá. Dva jedince spolu ale \emph{porovnat} umíme: stačí je nechat zahrát partii.

Druhý důvod je ambicióznější. Koevoluce slibuje \emph{evoluční závody ve zbrojení}, v~nichž
se protivníci navzájem tlačí ke stále lepším řešením. Podle toho, zda jsou si jedinci soupeři,
nebo spoluhráči, se rozlišují dva typy koevoluce:

\begin{center}
\begin{tabular}{@{}p{3cm}p{6cm}p{6cm}@{}}
\toprule
 & \textbf{kompetitivní} & \textbf{kooperativní} \\
\midrule
vztah & fitness jednoho roste na úkor druhého & jedinci se doplňují do
  společného řešení \\
příklad & predátor--kořist, host--parazit, hráč--protihráč, řadicí sítě vs.\
  posloupnosti čísel & dekompozice problému na podproblémy, moduly a~plány
  (CoDeepNEAT), hlavní programy a~ADF \\
fitness & výsledek souboje (turnaje) & kvalita složeného řešení, na němž se
  jedinec podílel \\
riziko & cyklení, odpojení, přespecializace & kredit assignment, líní partneři \\
\bottomrule
\end{tabular}
\end{center}

\subsubsection{Kompetitivní koevoluce}

\paragraph{Klasický příklad: Hillisovy řadicí sítě (1990).} Populace řadicích sítí zde
koevolvuje proti populaci testovacích posloupností. Fitness sítě je podíl správně seřazených
testů, fitness testu naopak to, kolik sítí na něm selhalo. Parazitické testy se proto
specializují na slabiny sítí a~sítím nezbývá než se zlepšovat.

Hillis takto nalezl síť pro 16 prvků s~\textbf{61} komparátory. \emph{Pozor na častý omyl:}
nejlepší známý \emph{lidský} návrh (Green, 1969) má 60 komparátorů, koevoluce ho tedy
nepřekonala. Podstatné je srovnání uvnitř experimentu, a~to vypadá jinak: týž genetický
algoritmus \emph{bez} koevolvujících parazitů dosáhl jen 65 komparátorů. Koevoluce zde
prokazatelně zlepšila optimalizační schopnost EA a~přiblížila se lidskému výsledku na jediný
komparátor.

\paragraph{Predátor--kořist} je druhý klasický vzorec a~je doma v~evoluční robotice:
roboti-lovci a~roboti-uprchlíci se navzájem zdokonalují, protože každé zlepšení jedné strany
je pro druhou novým zadáním. Dvojice \textbf{host--parazit} se objevuje v~modelech umělého
života (Tierra) a~při testování software.

\paragraph{Turnajové hodnocení.} Fitness se získává soubojem a~způsobů, jak souboje
uspořádat, je několik. Nejpoctivější je turnaj každý s~každým, ten je ale drahý. Levnější
a~zároveň zašuměnou variantou je souboj s~náhodným vzorkem protivníků. Metoda \emph{hall of
fame} nechává jedince utkat se s~archivem historicky nejlepších, a~brání tak zapomínání.
\emph{Turnaj se sdílenou fitness} (\emph{competitive fitness sharing}) navíc odstupňuje
odměnu: za poražení protivníka, kterého poráží málokdo, je odměna vyšší.

\subsubsection{Patologie koevoluce}

Koevoluce má pověst vrtkavého nástroje a~důvodem je řada dobře popsaných patologií, které
se u~ní objevují znovu a~znovu.

\begin{itemize}
\item \textbf{Cyklení} (\emph{cycling}, intranzitivita) nastává, když populace obíhají
  dokola jako v~kámen--nůžky--papír. Pokrok je pak jen zdánlivý a~dynamika koevoluce může
  být až deterministicky chaotická.
\item \textbf{Odpojení} (\emph{disengagement}) znamená, že jedna populace zcela přeroste
  druhou. Všechny souboje pak končí stejně, fitness ztrácí rozlišovací schopnost a~selekce
  přestane fungovat, protože zmizel gradient.
\item \textbf{Efekt Červené královny} (\emph{Red Queen effect}) nastane, když se všichni
  zlepšují, ale relativní fitness zůstává konstantní. Z~naměřených hodnot pak nepoznáme,
  jestli je pokrok skutečný; řešením je měřit proti pevné externí sadě protivníků
  (\emph{master tournament}).
\item \textbf{Průměrné stabilní stavy} (\emph{mediocre stable states}) vznikají, když se
  populace ustálí ve vzájemné rovnováze na nízké úrovni a~už nic nenutí ani jednu z~nich
  se zlepšovat.
\item \textbf{Přespecializace} (\emph{overspecialization}, \emph{focusing}) postihuje
  jedince, který se vyladí na aktuálního protivníka a~ztratí obecnost.
\item \textbf{Zapomínání} (\emph{forgetting}) je ztráta schopnosti porazit staré protivníky.
  Pomáhají proti němu archivy: hall of fame, Pareto-koevoluční archivy a~\emph{Nash memory}.
\end{itemize}

\subsubsection{Kooperativní koevoluce}

\begin{defbox}[CCGA (\emph{cooperative coevolutionary GA}, Potter \& De Jong)]
Problém se rozdělí na $k$ komponent, jimiž mohou být například souřadnice, moduly nebo
pravidla, a~pro každou komponentu se udržuje \textbf{samostatná populace}. Fitness jedince
z~populace $i$ se určí tak, že se jedinec \emph{spojí} se zástupci ostatních populací
(typicky s~jejich aktuálně nejlepšími jedinci, případně i~s~náhodnými) a~ohodnotí se
výsledné složené řešení.
\end{defbox}

Výhody jsou nasnadě: velký prostor se rozloží na menší, dekompozice se přirozeně
paralelizuje a~pro modulární úlohy je takové rozdělení věcně správné.

Problémy jsou tři. \emph{Přiřazení zásluh} (\emph{credit assignment}) řeší, jak rozdělit
kvalitu složeného řešení mezi jednotlivé komponenty. \emph{Provázanost} brání dobrému
oddělení tam, kde spolu komponenty silně interagují. A~\emph{relativní přespecializace}
znamená, že se populace přizpůsobí konkrétním partnerům místo úloze jako celku.

Praktických příkladů je řada: koevoluce hlavních programů a~ADF v~GP, dále CoDeepNEAT
s~populací modulů a~populací \uv{blueprintů} (odst.~\ref{sec:neuro}) a~konečně
meta-lamarckovské memetické algoritmy, kde vedle sebe žije populace memů a~populace řešení
(kap.~\ref{sec:lokalni}).

\subsubsection{Evoluce kooperace a~vězňovo dilema}

Koevoluce úzce souvisí se základní otázkou evoluční biologie: \emph{jak může vzniknout
altruismus}, když z~definice snižuje fitness jedince? Darwinismus je ve své podstatě
soutěživý, jak napovídá heslo \uv{přežití nejzdatnějšího}, a~přesto je v~přírodě
i~ve společnosti spolupráce všude kolem. Vysvětlení nabízí několik teorií.

\begin{itemize}
\item \textbf{Skupinová selekce} předpokládá, že evoluce působí i~na skupiny (Darwin).
  Naráží ale na otázku, jak vysvětlit podvodníky uvnitř skupiny.
\item \textbf{Příbuzenská selekce} (\emph{kin selection}) vysvětluje altruismus zachováním
  téměř identických genů u~příbuzných. Nevysvětlí ovšem altruismus vůči cizím jedincům
  a~jiným druhům.
\item \textbf{Sobecký gen} (Dawkins) posouvá jednotku evoluce od jedince ke genu.
  (Wilson to shrnuje větou \uv{organismus je jen způsob, jak DNA vyrábí další DNA}.)
\item \textbf{Reciproční altruismus} (Trivers, 1971) staví na vzájemném prospěchu
  a~\emph{stínu budoucnosti}; jeho formálním modelem je iterované vězňovo dilema.
\end{itemize}

\begin{defbox}[Vězňovo dilema (\emph{prisoner's dilemma})]
Hra dvou hráčů s~akcemi $C$ (kooperace) a~$D$ (zrada) a~výplatami
\[
\begin{array}{c|cc}
 & C & D\\\hline
C & R/R & S/T\\
D & T/S & P/P
\end{array}
\qquad
T > R > P > S \quad (\text{a pro iterovanou verzi } 2R > T+S).
\]
Typicky se volí $R=-1$, $T=0$, $P=-2$ a~$S=-3$. Protože $T > R$ a~zároveň $P > S$, je zrada
\textbf{dominantní strategií} obou hráčů a~dvojice $DD$ je \textbf{Nashovým ekvilibriem}.
Jako jediný z~možných výsledků přitom $DD$ \textbf{není Pareto-optimální}: společný zisk
maximalizuje $CC$.
\end{defbox}

Racionální hráči tedy dospějí k~výsledku, který je pro oba horší než dosažitelná spolupráce.
Tentýž mechanismus v~$n$-hráčské podobě se nazývá \emph{tragédie obecní pastviny}
(\emph{tragedy of the commons}): každému se individuálně vyplatí čerpat sdílený zdroj
o~něco víc, a~zdroj proto zanikne. Jádrem celé diskuse kolem vězňova dilematu je právě
otázka, co je vlastně racionální a~zda se tak lidé opravdu chovají.

\begin{defbox}[Dominantní strategie, Nashovo ekvilibrium, Pareto optimalita]
Strategie $s$ je \emph{dominantní} pro hráče $i$, dává-li stejný nebo lepší výsledek než
kterákoli jiná strategie $i$ proti \emph{všem} strategiím protihráče. Dvojice $(s_i, s_j)$
je v~\emph{Nashově ekvilibriu}, jsou-li obě strategie vzájemně nejlepšími odpověďmi, takže
žádný hráč nemá důvod jednostranně změnit svou volbu. Řešení je \emph{Pareto-optimální},
pokud nelze zlepšit výsledek jednoho hráče, aniž by se zhoršil výsledek jiného.

Nashova věta říká, že každá hra s~konečně mnoha strategiemi má ekvilibrium ve \emph{smíšených}
strategiích, tedy v~náhodné volbě mezi čistými; příkladem je kámen--nůžky--papír.
\end{defbox}

\paragraph{Iterovaná verze a~Axelrodovy turnaje.} Hraje-li se opakovaně a~hráči si pamatují
historii, může se spolupráce vyplatit, protože nad hrou visí \emph{stín budoucnosti}.
Pozor na jednu výjimku: zná-li se předem přesný počet kol $N$, plyne indukcí od konce, že
optimální je stále zrazovat. V~praxi ale hráči konec neznají. Axelrod (1980) uspořádal
turnaje strategií a~vyšly z~nich tři experimenty:

\begin{itemize}
\item 1.\ turnaj sestával ze 14 strategií a~RANDOM, hrálo se 200 her, každý s~každým
  (a~i~sám se sebou) v~5 opakováních. Vyhrála \textbf{TFT} (\emph{tit for tat}): začni
  kooperací, pak opakuj poslední tah soupeře.
\item 2.\ turnaj měl 62 strategií a~všichni účastníci znali výsledky prvního; TFT
  vyhrála znovu.
\item 3.\ \uv{ekologický} turnaj byl simulací po vzoru GA: počet jedinců dané strategie
  v~další generaci byl úměrný počtu vítězství a~běželo se 1000 generací. TFT opět zvítězila.
\end{itemize}

\paragraph{Čtyři vlastnosti úspěšných strategií:}
\emph{slušnost} (niceness) znamená nezrazovat první;
\emph{dráždivost} (provocability/retaliation) žádá potrestat zradu okamžitě;
\emph{odpouštění} (forgiveness) velí se zase rychle uklidnit.
Nejméně samozřejmá je \emph{nezávistivost} (non-enviousness): úspěšná strategie neusiluje
o~to porazit \emph{konkrétního} soupeře. TFT nikdy nezískala víc bodů než její protihráč,
a~přesto turnaj vyhrála.
Jako pátou vlastnost Axelrod uvádí \emph{srozumitelnost} (clarity): buď jednoduchý, aby
s~tebou ostatní vůbec dokázali navázat spolupráci.
Žádná strategie nevyhraje proti všem, takže úspěch znamená obstát proti velmi rozmanitým
soupeřům (ALL-D, TFTT, RANDOM, TRIGGER) a~umět dobře hrát i~sám se sebou.

\paragraph{Poučení a~pozdější vývoj.} V~prostředích, kde výplaty přejí spolupráci a~kde je
vysoká pravděpodobnost opakování, se spolupráce obvykle vyvine. Není k~tomu potřeba
racionalita ani vědomí, stačí vyšší fitness.

Obrázek se ovšem mění, jakmile se do hry vloží šum. V~zašuměném prostředí, kde se tah může
provést chybně, je úspěšnější strategie \textbf{Pavlov} (\emph{win-stay, lose-shift}):
dostal-li jsem $R$ nebo $P$, zopakuj tah ($C$); dostal-li jsem $T$ nebo $S$, změň jej ($D$).
Turnaj opakovaný po 20 letech pak vyhrál \emph{tým} strategií z~University of Southampton.
Prvních zhruba 10 tahů sloužilo k~rozpoznání spoluhráče a~poté všichni členové týmu obětavě
pomáhali jedné \uv{otcovské} strategii získat vysoké skóre.

\subsubsection{Diverzita, druhy a~niky}

Malý selekční tlak stačí k~tomu, aby diverzita populace vymizela. (Podobnou úvahu rozvíjí
Flegrova \emph{zamrzlá evoluce}: geneticky proměnlivý druh se vůči selekci chová pružně jen
krátce, poté \uv{zamrzne} a~na změny prostředí přestane reagovat.) Existuje proto řada
mechanismů, jak diverzitu udržet; hodí se stejně pro koevoluci jako pro dynamickou fitness
a~vícekriteriální optimalizaci.

\begin{itemize}
\item \textbf{Fitness sharing} dělí fitness počtem podobných jedinců,
  $f'(i) = f(i)/\sum_j sh(d(i,j))$, kde
  $sh(d) = 1-(d/\sigma_{\mathrm{share}})^{\alpha}$ pro $d < \sigma_{\mathrm{share}}$
  a~jinak 0. Vytváří se tím tlak na obsazení různých nik.
\item \textbf{Crowding} nechává nového jedince nahradit \emph{nejpodobnějšího} jedince
  populace; patří sem deterministic crowding a~RTS.
\item \textbf{Nenáhodné páření} (\emph{non-random mating}) povoluje křížení jen mezi
  podobnými jedinci, čemuž se říká \emph{speciace}, případně je řídí topologií: jedinci
  žijí na 2D/3D mřížce a~kříží se jen se sousedy.
\item \textbf{Ostrovní model} (\emph{island model}) rozdělí populaci na několik subpopulací
  s~občasnou migrací; podrobněji o~něm níže.
\item \textbf{Explicitní druhy} se zavádějí značkovacími bity (\emph{tag bits}) nebo se
  odvozují ze strukturální podobnosti genomů (NEAT, odst.~\ref{sec:neuro}); páření je pak
  povoleno jen uvnitř druhu.
\end{itemize}
Slabinou všech těchto technik je, že se musí definovat podobnost, ať už jako Hammingova
vzdálenost, vzdálenost v~prostoru fenotypů, nebo počet shodných inovačních čísel, a~k~tomu
zvolit práh niky.

\paragraph{Paralelní modely EA.} Struktura populace úzce souvisí s~paralelizací a~rozlišují
se tři modely.

\begin{itemize}
\item \textbf{Master--slave} neboli globální paralelizace pracuje s~jedinou populací a~mezi
  uzly rozděluje pouze \emph{vyhodnocení fitness}. Algoritmus je matematicky totožný se
  sekvenčním a~vyplatí se, je-li fitness drahá.
\item \textbf{Ostrovní neboli hrubozrnný model} (\emph{coarse-grained}) udržuje několik
  nezávislých subpopulací, mezi nimiž občas \emph{migruje} několik jedinců. Ladí se u~něj
  topologie ostrovů, \emph{migrační interval} a~\emph{migrační míra} a~také výběr migrantů
  a~nahrazovaných jedinců. Vzácná migrace udržuje diverzitu, protože ostrovy konvergují
  jinam; příliš častá migrace model degeneruje na jedinou populaci.
\item \textbf{Buněčný neboli jemnozrnný model} (\emph{fine-grained}, difuzní) usadí jedince
  na 2D/3D mřížku a~kříží je jen se sousedy. Dobré řešení se pak šíří populací jako vlna,
  což silně zpomaluje předčasnou konvergenci; svou podstatou je to implicitní forma
  nenáhodného páření.
\end{itemize}
Ostrovní i~buněčný model tedy nejsou jen implementační trik. Mění samotnou dynamiku
algoritmu a~patří mezi standardní nástroje udržení diverzity.

\subsubsection{Dynamické fitness krajiny}

\begin{defbox}[Fitness krajina (\emph{fitness landscape})]
Pojem zavedl Sewall Wright (1932). Jde o~grafické znázornění závislosti fitness na genotypu,
případně na frekvenci alel či na rysu fenotypu. Je to teoretický nástroj pro uvažování o~EA;
ve vyšších dimenzích ovšem přestává být intuitivní.
\end{defbox}

Fitness se často mění v~čase. Robot pracuje v~místnosti, kde někdo rozsvítí; začne pršet;
na burze propukne krize; agenti se utkávají v~turnajích, tedy koevolvují. Klasické GA si
v~takových podmínkách vedou špatně, protože jsou \uv{přeučené} na statické úlohy: rychle
konvergují a~ztratí přitom právě tu diverzitu, kterou by pro reakci na změnu potřebovaly.
De Jong to doložil srovnáním: $(1+10)$-ES reaguje na náhlou změnu pomalu, protože si
samoadaptací zmenšila krok, kdežto generační GA s~ruletovou selekcí a~bez elitismu si udrží
více diverzity a~zotaví se rychleji.

Pro dynamické krajiny se proto EA upravují několika způsoby.

\begin{itemize}
\item \textbf{Diploidní reprezentace} (Goldberg, Smith) používá dvě sady chromozomů
  s~dominancí, které při oscilující fitness fungují jako dlouhodobá paměť.
\item \textbf{Hypermutace} (Cobb, Grefenstette) sleduje průměrnou fitness populace; klesne-li
  (pravděpodobně se změnilo prostředí), výrazně se zvýší pravděpodobnost mutace.
\item \textbf{Náhodná migrace} vloží do populace dávku nových náhodných jedinců.
\item \textbf{Udržování diverzity} zahrnuje nižší selekční tlak, crowding, niching
  a~subpopulace, ale i~volbu čárkové ES $(\mu,\lambda)$: rodiče v~ní nepřežívají, takže se
  populace snáz odpoutá od zastaralého optima.
\end{itemize}
Změny samotné se dělí na malé plynulé (opotřebení hardwaru), morfologické (vznikají
a~zanikají kopce), cyklické (roční období) a~katastrofické nespojité. Mění-li se fitness
příliš rychle, žádné řešení není trvale dobré (srov.\ No Free Lunch).

\subsection{Otevřená evoluce}
\label{ssec:otevrena}

\noindent\mimoq{celý oddíl kromě interaktivní evoluce} Otevřenou evoluci okruh výslovně
jmenuje, slidy EVA~I ani EVA~II o~ní ale nemají jedinou stránku. Jediná zmínka se najde
v~\emph{Evoluční robotice}, a~to v~jedné větě (\uv{evolúcia robotov je kontinuálny proces
s~otvoreným koncom}). Následující oddíl je proto zpracovaný stručně, na úrovni, která na
zkoušce obstojí. Výjimkou je \emph{interaktivní evoluce} na konci oddílu: ta ve slidech
Evoluční robotiky opravdu je, včetně subjektivní fitness a~Biomorphs.

\begin{defbox}[Otevřená evoluce (\emph{open-ended evolution})]
Evoluční proces, který \textbf{nemá pevný cíl} a~který neustále produkuje novou složitost
a~nové \uv{druhy} chování, aniž by konvergoval. Vzorem je pozemská biosféra: žádnou účelovou
funkci nemá, a~přesto trvale generuje novinky.
\end{defbox}

\paragraph{Znaky otevřenosti.} O~konkrétním systému je těžké rozhodnout, jestli je otevřený,
nebo jen dlouho běží. Proto se formulují \emph{znaky} (\emph{hallmarks}), podle kterých se to
dá posoudit. Prvním z~nich je \textbf{trvalá produkce novosti}: nové chování se objevuje
i~po libovolně dlouhém běhu, ne jenom v~úvodní fázi. Druhým je \textbf{neomezený růst
složitosti}, a~to jak u~samotných jedinců, tak u~vztahů mezi nimi. Třetím je \textbf{vznik
nových tříd entit} a~nových úrovní organizace (v~biologii \emph{major transitions}:
replikátor $\to$ buňka $\to$ mnohobuněčnost $\to$ společenství). Čtvrtým znakem je
\textbf{evoluce evolvability}, protože se mění i~samotný způsob, jakým systém novinky
generuje.

\paragraph{Systémy a~proč se zasekávají.} Pokusů postavit otevřený systém proběhla celá řada.
\emph{Tierra} (Ray) a~\emph{Avida} pracují se sebekopírujícími programy ve strojovém kódu
a~emergentně v~nich vznikli paraziti a~hyperparaziti. \emph{Polyworld} osazuje 2D svět agenty
s~neuronovými sítěmi, \emph{ECHO} (Holland) agenty, kteří si nesou \uv{tagy} a~soupeří
o~zdroje. Prakticky každý z~těchto systémů ale po čase dospěje do stavu, kdy se už nic
kvalitativně nového neobjevuje: prostředí je konečné, prostor možných \uv{vynálezů} se
vyčerpá a~dynamika zamrzne.

Proč se to děje, se zatím spíš odhaduje. Nabízejí se tři hypotézy. Chybí dostatečně bohaté
a~samo se rozšiřující prostředí; chybí skutečně otevřená reprezentace; a~chybí mechanismus,
který by průběžně zvyšoval nároky kladené na jedince. Dosažení skutečně otevřené evoluce
tak zůstává otevřeným výzkumným problémem.

\paragraph{Generativní (nepřímé) reprezentace.} S~otevřenou evolucí souvisí otázka, jak vůbec
zakódovat složitost, která má neomezeně růst. Přímý popis výsledku by musel růst spolu s~ním.
Odpovědí jsou proto kódování, v~nichž genotyp není popisem výsledku, ale \textbf{návodem, jak
výsledek vypěstovat}. Patří sem \emph{L-systémy}, tedy přepisovací gramatiky generující
větvící se struktury, dále \emph{celulární automaty}, \emph{buněčné kódování} (Gruau)
a~pro neuronové sítě \emph{CPPN/HyperNEAT} (odst.~\ref{sec:neuro}). Přínos mají všechna
společný: malý genom, velký fenotyp a~přirozené vyjádření symetrií i~opakujících se motivů.

\begin{defbox}[Novelty search a~řídkost]
Radikální myšlenku přinesli Lehman \& Stanley (2011) v~článku \emph{Abandoning Objectives}:
u~deceptivních úloh je účelová funkce zavádějící, protože vede do lokálních optim.
Nabízí se tedy postup přesně opačný, totiž \textbf{zahoďme cíl a~odměňujme pouze novost}.

Definuje se \emph{prostor chování} (např.\ koncová pozice robota v~bludišti) a~novost
jedince $x$ se měří řídkostí jeho okolí
\[
\rho(x) \;=\; \frac{1}{k}\sum_{i=1}^{k} \mathrm{dist}\big(x, \mu_i\big),
\]
kde $\mu_i$ je $k$ nejbližších sousedů $x$ v~aktuální populaci \textbf{a~v~permanentním
archivu}. Jako fitness slouží přímo $\rho$ a~původní účelová funkce se \emph{nepoužívá
vůbec}. Na úlohách typu \uv{robot v~bludišti} bývá takový postup úspěšnější než cílená
optimalizace. Navazuje na něj \textbf{quality-diversity}, konkrétně MAP-Elites: prostor
chování se pokryje mřížkou buněk indexovanou deskriptory chování a~v~každé buňce se udržuje
nejlepší nalezené řešení.
\end{defbox}

\paragraph{Interaktivní evoluce.} Ještě jedna cesta vede k~evoluci bez zapsaného cíle:
fitness funkci nahradí \textbf{člověk}. Uživateli se předloží populace kandidátů a~on vybere
ty, které se mu líbí. Uplatní se to tam, kde kritérium kvality neumíme zapsat, tedy
u~estetiky, designu nebo hudby. Klasickou ukázkou jsou Dawkinsovy \emph{Biomorphs}, dále
\textbf{Picbreeder} a~\textbf{EndlessForms}, kde se pomocí CPPN vyvíjejí obrázky a~3D tvary.
Omezením je pomalost a~únava uživatele, takže se nutně pracuje s~malými populacemi
(9--20 jedinců na obrazovku) a~s~malým počtem generací. Právě zkušenost z~Picbreederu byla
přímou motivací pro novelty search: nejzajímavější objekty tam nikdy nevznikly cíleným
hledáním.

\subsection{Shrnutí ke~zkoušce}

\begin{itemize}
\item Evoluční strategie pracují s~reálnými vektory: hlavním operátorem je mutace,
  rodičovská selekce probíhá náhodně a~environmentální selekce deterministicky,
  přičemž jedinec $= [\vec x, \vec\sigma]$.
\item U~dvojice $(\mu,\lambda)$ vs.\ $(\mu+\lambda)$ je nutné umět vyjmenovat rozdíly,
  tedy elitismus, rychlost konvergence, chování v~lokálních optimech a~vliv na samoadaptaci.
\item Velikost kroku v~$(1+1)$-ES řídí pravidlo 1/5: $p > 1/5 \Rightarrow$ zvětši $\sigma$,
  $p < 1/5 \Rightarrow$ zmenši. Jde o~řízení adaptivní, nikoli samoadaptivní.
\item Samoadaptace zmutuje nejdřív strategický parametr a~teprve pak proměnné:
  $\sigma' = \sigma \exp(\tau'N(0,1) + \tau N_i(0,1))$ a~poté $x' = x + \sigma' N(0,1)$,
  tedy \textbf{nejdřív $\sigma$, pak $x$}. Strategický parametr je jeden u~izotropní mutace,
  je jich $n$ u~nekorelované a~$n(n+1)/2$ u~korelované, kde přibývají úhly natočení.
\item Rekombinace se dělí podle počtu rodičů na lokální ($\rho=2$) a~globální ($\rho=\mu$)
  a~podle způsobu míchání na diskrétní (dominantní) a~intermediární (aritmetickou).
\item CMA-ES si nese stav $\theta = (\vec m,\sigma,\mathbf C,\vec p_\sigma,\vec p_c)$,
  vzorkuje z~$\mathcal N(\vec m, \sigma^2\mathbf C)$ a~kovarianční matici upravuje pomocí
  evolučních cest rank-1 a~rank-$\mu$ updatem; velikost kroku řídí délka $\vec p_\sigma$.
  Metoda je invariantní vůči monotónním transformacím fitness a~lineárním transformacím
  prostoru.
\item Schéma DE/rand/1/bin vytvoří dárce $\vec v = \vec x_a + F(\vec x_b - \vec x_c)$,
  zkříží ho binomicky s~prahem $C$ a~pojistkou $j = I_{\mathrm{rand}}$ a~potomka nechá
  soubojit s~vlastním rodičem, kterého nahradí jen při zlepšení.
\item Parametry DE jsou $F \in [0,2]$ (typicky $0{,}5$--$0{,}9$), $C \in [0,1]$
  a~velikost populace $N \approx 10D$.
\item Geometrická podstata DE spočívá v~tom, že se měřítko i~orientace kroků přebírají
  z~aktuálního rozptylu populace, takže adaptace probíhá automaticky a~bez strategických
  parametrů.
\item Nejčastější varianty jsou rand/1, rand/2, best/1, best/2 a~current-to-best,
  křížení pak \texttt{bin} nebo \texttt{exp}.
\item U~zkoušky se žádá předvedení jednoho kroku s~konkrétními čísly
  (mutace $\to$ křížení $\to$ selekce).
\item Koevoluce znamená kontextovou fitness a~dynamickou krajinu.
  \emph{Kompetitivní} podobu mají úlohy predátor--kořist, host--parazit a~Hillisovy řadicí
  sítě, \emph{kooperativní} pak dekompozice na komponenty, CCGA a~schéma
  moduly $+$ blueprinty.
\item Mezi patologie patří cyklení (in\-tran\-zi\-ti\-vi\-ta), odpojení, efekt Červené
  královny, průměrné stabilní stavy, přespecializace a~zapomínání. Lékem jsou archivy
  (hall of fame), sdílená fitness a~měření proti pevné sadě.
\item Ve vězňově dilematu platí $T>R>P>S$ a~$2R > T+S$; $D$ je dominantní strategie
  a~$DD$ Nashovo ekvilibrium, které jako jediné není Pareto-optimální. V~iterované verzi
  se prosadí TFT a~úspěšná strategie má čtyři vlastnosti (slušnost, dráždivost, odpouštění,
  nezávistivost), k~nimž se jako pátá přidává srozumitelnost; v~zašuměném prostředí
  se osvědčí Pavlov.
\item Diverzitu udržují fitness sharing, crowding, speciace, ostrovní model
  a~nenáhodné páření.
\item Na dynamickou fitness se odpovídá diploidií, hypermutací, náhodnou migrací,
  udržováním diverzity a~$(\mu,\lambda)$-ES.
\item Otevřená evoluce běží bez pevného cíle a~pozná se podle znaků: trvalé novosti,
  neomezeného růstu složitosti, vzniku nových tříd entit a~evoluce evolvability.
  Systémy Tierra, Avida, Polyworld a~ECHO se všechny dříve či později zaseknou.
  Neomezeně rostoucí složitost kódují generativní (nepřímé) reprezentace: L-systémy,
  celulární automaty, buněčné kódování a~CPPN.
\item Novelty search měří řídkost $\rho(x)$ v~prostoru chování, a~to vůči populaci
  \emph{a} permanentnímu archivu, zatímco účelovou funkci vůbec nepoužívá;
  navazuje na něj quality-diversity a~MAP-Elites s~mřížkou deskriptorů chování.
\item V~interaktivní evoluci hraje roli fitness člověk (Biomorphs, Picbreeder,
  EndlessForms); omezením je malá populace a~únava uživatele.
\end{itemize}

%=====================================================================
\section{Rojové optimalizační algoritmy}
\label{sec:roje}
%=====================================================================

Hejno špačků se dokáže srovnaně vyhnout dravci, aniž by mělo velitele, a~mravenčí kolonie
najde nejkratší cestu ke~zdroji potravy, přestože žádný jednotlivý mravenec tu cestu nezná.
Z~takových pozorování vychází rojová inteligence: nehledá se chytřejší jedinec, ale takové
lokální chování, ze kterého se rozumné řešení složí samo.

\begin{defbox}[Rojová inteligence (\emph{swarm intelligence})]
Přístup, kde globálně inteligentní chování \textbf{emerguje} z~jednoduchých
lokálních pravidel mnoha jedinců bez centrálního řízení. Typické rysy jsou čtyři: řízení je
decentralizované, jedinci jsou jednodušší, komunikace probíhá nepřímo prostřednictvím
prostředí (\emph{stigmergie}) a~systém je robustní vůči výpadku jedince, protože žádný
jedinec není nenahraditelný. Na rozdíl od EA zde obvykle není křížení, mutace ani \uv{smrt}
jedinců: jedinci místo toho trvale existují, pohybují se a~pamatují si.
\end{defbox}

\subsection{Optimalizace rojem částic (PSO)}

\emph{Particle swarm optimization} (Eberhart a~Kennedy, 1995) je inspirována hejny ptáků
a~ryb. Pták v~hejnu nemá mapu potravy. Řídí se dvěma věcmi: vzpomínkou na místo, kde se mu
dosud dařilo nejlépe, a~tím, kam míří hejno kolem něj. Právě tyto dva vlivy algoritmus
překládá do vzorce.

Jedinec se proto nejmenuje jedinec, ale \emph{částice}, a~je to bod v~$\mathbb{R}^n$
s~rychlostí. Nepoužívá se křížení ani mutace; namísto nich má algoritmus \textbf{paměť},
a~to hned dvojí:
\begin{itemize}
\item $\vec p_i$ (\emph{pBest}) je nejlepší pozice, kterou částice $i$ kdy navštívila,
  a~představuje tedy paměť kognitivní, individuální,
\item $\vec g$ (\emph{gBest}) je nejlepší pozice nalezená celým rojem, tedy paměť sociální
  a~kolektivní.
\end{itemize}

Pohyb částice pak vznikne složením tří tendencí: pokračovat setrvačně dosavadním směrem,
zatočit ke~své nejlepší vzpomínce a~zatočit k~nejlepšímu místu roje.

\begin{defbox}[Pohybové rovnice PSO]
\[
\vec v_i \;\gets\; \underbrace{w\,\vec v_i}_{\text{setrvačnost}}
\;+\; \underbrace{c_1 r_1 (\vec p_i - \vec x_i)}_{\text{kognitivní složka}}
\;+\; \underbrace{c_2 r_2 (\vec g - \vec x_i)}_{\text{sociální složka}},
\qquad
\vec x_i \;\gets\; \vec x_i + \vec v_i ,
\]
kde $r_1, r_2 \sim U(0,1)$ (losují se pro každou složku zvlášť),
$c_1, c_2$ jsou koeficienty učení (v~původní verzi $c_1 = c_2 = 2$)
a~$w$ je \emph{setrvačná váha} (\emph{inertia weight}).
Původní verze z~roku 1995 setrvačnou váhu ještě nemá, což odpovídá volbě $w = 1$. Doplnili ji
až Shi a~Eberhart a~doporučují nechat ji lineárně klesat, například z~$0{,}9$ na $0{,}4$:
velké $w$ podporuje exploraci, malé exploataci.
\end{defbox}

\begin{algorithm}[!ht]
\caption{PSO}
\begin{algorithmic}[1]
\State inicializuj pozice $\vec x_i$ a~rychlosti $\vec v_i$ náhodně; $\vec p_i \gets \vec x_i$
\Repeat
  \For{každou částici $i$}
    \State vyhodnoť $f(\vec x_i)$
    \If{$f(\vec x_i)$ je lepší než $f(\vec p_i)$} $\vec p_i \gets \vec x_i$ \EndIf
  \EndFor
  \State $\vec g \gets$ nejlepší z~$\{\vec p_i\}$
  \For{každou částici $i$}
    \State aktualizuj rychlost a~pozici dle pohybových rovnic
  \EndFor
\Until{je dosaženo maxima iterací nebo požadované přesnosti}
\end{algorithmic}
\end{algorithm}

\paragraph{Číselný příklad.} Spočítejme jeden krok pro částici v~$\mathbb{R}^2$. Ta má pozici
$\vec x = (1;\,2)$ a~rychlost $\vec v = (0{,}5;\,-0{,}5)$, její vlastní nejlepší pozice je
$\vec p = (0;\,3)$ a~nejlepší pozice roje $\vec g = (-1;\,1)$. Parametry jsou $w = 0{,}7$,
$c_1 = c_2 = 1{,}5$ a~náhodná čísla vyšla $r_1 = 0{,}4$, $r_2 = 0{,}8$.
\begin{align*}
\vec v' &= 0{,}7\,(0{,}5;\,-0{,}5) + 1{,}5\cdot 0{,}4\,\big((0;3)-(1;2)\big)
  + 1{,}5\cdot 0{,}8\,\big((-1;1)-(1;2)\big)\\
 &= (0{,}35;\,-0{,}35) + 0{,}6\,(-1;\,1) + 1{,}2\,(-2;\,-1)
  = (-2{,}65;\; -0{,}95),\\
\vec x' &= (1;\,2) + (-2{,}65;\,-0{,}95) = (-1{,}65;\; 1{,}05).
\end{align*}
Částice tedy neskončila u~$\vec g$, nýbrž přeletěla za ně. To je pro PSO typické a~je to
zdroj explorace: roj cíl obletuje a~přitom prohledává jeho okolí. Přestřelení ale nesmí být
neomezené, a~tak se rychlost buď shora ořezává podmínkou $|v_j| \le v_{\max}$, nebo se použije
\emph{constriction faktor} (Clerc a~Kennedy), který celou aktualizaci srazí jedním
koeficientem:
$\vec v \gets \chi[\vec v + c_1 r_1(\vec p - \vec x) + c_2 r_2 (\vec g - \vec x)]$,
kde $\chi = 2/|2 - \varphi - \sqrt{\varphi^2 - 4\varphi}|$ pro
$\varphi = c_1 + c_2 > 4$. Pro obvyklou volbu $c_1=c_2=2{,}05$ vychází $\chi \approx 0{,}729$
a~tato varianta má zaručenou konvergenci roje.

\paragraph{Topologie sousedství.} Sdílet jedno globální $\vec g$ znamená, že se každý dozví
o~každém zlepšení okamžitě, což roj rychle stáhne k~jednomu místu. Této topologii se říká
\emph{gbest} a~odpovídá úplnému grafu. Alternativa \emph{lbest} nechává částici ohlížet se jen
na nejlepšího jedince ve svém lokálním okolí:
\begin{itemize}
\item \textbf{Kruh}, kde každá částice má $k$ sousedů, šíří informaci pomalu. Konvergence je
  proto pomalejší, zato je lepší explorace a~menší riziko uvíznutí.
\item Používá se také \textbf{Von Neumannova mřížka}, \textbf{hvězda} nebo \textbf{náhodné
  grafy}.
\item Empiricky platí, že gbest je rychlejší na jednoduchých úlohách, kdežto lbest je
  robustnější na úlohách multimodálních.
\end{itemize}

\paragraph{Varianty a~vztah k~EA.}
Existuje i~diskrétní a~binární PSO, kde se rychlost interpretuje jako pravděpodobnost
překlopení bitu přes sigmoidu. Dále se používá PSO s~více roji, hybridizace s~lokálním
prohledáváním nebo \emph{grammatical swarm}. Následující tabulka shrnuje, čím se PSO liší od
genetického algoritmu.

\begin{center}
\begin{tabular}{@{}p{5cm}p{5cm}p{5.2cm}@{}}
\toprule
 & \textbf{genetický algoritmus} & \textbf{PSO} \\
\midrule
společné & náhodná inicializace, fitness, stochastické aktualizace &
  totéž \\
jedinci & vznikají a~zanikají & trvale existují, pohybují se \\
paměť & jen v~populaci & explicitní ($\vec p_i$, $\vec g$) \\
operátory & křížení, mutace & žádné, jen pohyb \\
tok informace & obousměrný mezi rodiči & jednosměrný: od lepších k~ostatním \\
\bottomrule
\end{tabular}
\end{center}

\subsection{Optimalizace mravenčí kolonií (ACO)}

\noindent\mimoq{ve slidech EVA je z~rojových algoritmů jen PSO; okruh ale mluví
o~rojových algoritmech v~množném čísle a~ACO je jejich kanonický druhý zástupce}

\emph{Ant colony optimization} (M.~Dorigo, 1991) napodobuje způsob, jakým hledají cestu
mravenci. Mravenec za sebou pokládá \emph{feromonovou stopu}. Po kratší cestě se vrátí dřív,
takže se na ní feromon hromadí rychleji, a~silnější stopa přitáhne další mravence: je to
kladná zpětná vazba, která krátkou cestu sama zvýrazní. Aby se kolonie nezasekla na první
slušné trase, feromon se navíc \textbf{odpařuje}, čímž systém postupně zapomíná zastaralou
informaci.

\begin{defbox}[Stigmergie]
Nepřímá komunikace prostřednictvím změn v~prostředí. Mravenci spolu
nekomunikují přímo; jediným kanálem je feromon na hranách. Je to obecný princip
reaktivních multiagentních architektur.
\end{defbox}

\paragraph{Ilustrace --- Steelsův průzkumník Marsu.} Že stigmergie stačí i~na netriviální
kolektivní úlohu, ukazuje známý myšlenkový příklad. Roj robotů má sbírat vzorky hornin
a~základna vysílá navigační signál, takže k~návratu stačí sledovat gradient. Robot má
k~dispozici \uv{drobky}, tedy radioaktivní značky sloužící k~nepřímé komunikaci. Chování
zadává pět pravidel s~prioritou. R1: vidíš-li překážku, změň směr. R2: neseš-li vzorek a~jsi
na základně, polož ho. R3: neseš-li vzorek, jdi po gradientu nahoru. R4: vidíš-li vzorek,
seber ho. R5: jinak se pohybuj náhodně. Druhá iterace návrhu přidá stigmergii. R7: neseš-li
vzorek a~nejsi na základně, polož 2~drobky a~jdi po gradientu nahoru. R8: cítíš-li drobek,
seber 1~drobek a~jdi po gradientu dolů. Tím vznikne efektivní kolektivní sběr z~bohatých
nalezišť, aniž by roboti měli mapu či vzájemnou komunikaci.

\paragraph{Obecné schéma ACO.} Aby se mravenci dali nasadit na optimalizační úlohu, převede
se úloha na hledání cesty v~ohodnoceném grafu; mravenci jsou pak agenti, kteří skládají
řešení po hranách. Jedna iterace má čtyři kroky:
\begin{enumerate}
\item Každý mravenec stochasticky zkonstruuje řešení, tedy posloupnost hran.
\item Volitelně se provedou \emph{akce démona} (\emph{daemon actions}), což je centralizovaný
  krok, typicky lokální prohledávání zlepšující nalezené cesty.
\item Porovnají se kvality nalezených řešení.
\item Aktualizují se feromony na hranách, a~to odpařením a~uložením nového feromonu.
\end{enumerate}

Volbu následujícího města řídí kompromis mezi tím, co se osvědčilo kolonii, a~tím, co vypadá
dobře lokálně.

\begin{defbox}[Ant System (AS) --- pravidlo přechodu]
Mravenec $k$ v~městě $i$ zvolí následující město $j$ z~množiny dosud
nenavštívených $\mathcal{N}_i^k$ s~pravděpodobností
\[
p_{ij}^{k} \;=\;
\frac{\big[\tau_{ij}\big]^{\alpha}\,\big[\eta_{ij}\big]^{\beta}}
     {\sum_{l \in \mathcal{N}_i^k}\big[\tau_{il}\big]^{\alpha}\,\big[\eta_{il}\big]^{\beta}},
\qquad \eta_{ij} = \frac{1}{d_{ij}},
\]
kde $\tau_{ij}$ je množství feromonu na hraně, $\eta_{ij}$ heuristická
\uv{viditelnost}, tedy převrácená vzdálenost, a~$\alpha, \beta$ váhy obou vlivů
(typicky $\alpha = 1$, $\beta \in [2,5]$).
\end{defbox}

\begin{defbox}[Ant System --- aktualizace feromonu]
\[
\tau_{ij} \;\gets\; (1-\rho)\,\tau_{ij} \;+\; \sum_{k=1}^{m} \Delta\tau_{ij}^{k},
\qquad
\Delta\tau_{ij}^{k} =
\begin{cases}
Q/L_k & \text{prošel-li mravenec } k \text{ hranou } (i,j),\\
0 & \text{jinak,}
\end{cases}
\]
kde $\rho \in (0,1)$ je míra odpařování a~$L_k$ délka trasy mravence $k$.
Příspěvek je nepřímo úměrný délce trasy, takže kratší trasa uloží více feromonu.
\end{defbox}

\paragraph{Číselný příklad (pravidlo přechodu).} Mravenec stojí ve městě $i$
a~může jít do měst $A, B, C$ se vzdálenostmi $d = (5;\,10;\,2)$
a~feromony $\tau = (2;\,4;\,1)$, přičemž $\alpha = 1$ a~$\beta = 2$.
Viditelnosti jsou převrácené vzdálenosti, tedy $\eta = (0{,}2;\,0{,}1;\,0{,}5)$, a~nenormované
váhy vyjdou takto:
\[
w_A = 2 \cdot 0{,}2^2 = 0{,}08,\quad
w_B = 4 \cdot 0{,}1^2 = 0{,}04,\quad
w_C = 1 \cdot 0{,}5^2 = 0{,}25,\quad \textstyle\sum = 0{,}37 .
\]
Po vydělení součtem dostaneme $p_A = 0{,}216$, $p_B = 0{,}108$ a~$p_C = 0{,}676$. Přestože
hrana do $B$ nese nejvíce feromonu, rozhodne krátká vzdálenost do $C$, protože se do váhy
promítne umocněná na $\beta = 2$.

\paragraph{Číselný příklad (aktualizace).} Nechť $\tau_{ij} = 2$, $\rho = 0{,}5$,
$Q = 1$ a~hranou prošli dva mravenci s~délkami tras $L_1 = 20$ a~$L_2 = 25$. Nejprve se
polovina feromonu odpaří, pak oba mravenci přispějí:
\[
\tau_{ij} \gets 0{,}5 \cdot 2 + \left(\tfrac{1}{20} + \tfrac{1}{25}\right)
= 1 + 0{,}09 = 1{,}09 .
\]

\paragraph{ACS (Ant Colony System).} Základní Ant System je spíše ilustrativní; teprve ACS je
vylepšení konkurenceschopné se specializovanými řešiči TSP. Inspiruje se Q-learningem a~mění
tři věci:
\begin{itemize}
\item \textbf{Globálně aktualizuje jen nejlepší řešení} (elitismus), takže feromon přidává
  pouze globálně, případně iteračně nejlepší mravenec.
\item \textbf{Lokální aktualizace} sníží feromon na hraně pokaždé, když po ní někdo projde
  ($\tau \gets (1-\xi)\tau + \xi\tau_0$). Hrana se tím stává méně atraktivní pro ostatní
  mravence a~roste rozmanitost tras.
\item \textbf{Pseudonáhodné pravidlo výběru} rozhoduje mezi exploatací a~explorací losem:
  s~pravděpodobností $q_0$ se zvolí deterministicky nejlepší hrana podle
  $\tau_{il}\eta_{il}^{\beta}$, jinak se použije standardní pravidlo AS.
\end{itemize}

\paragraph{Další varianty.}
V~\emph{Elitist AS} přispívá nejlepší trasa navíc oproti ostatním. \emph{Rank-based AS} nechá
přispět $N$ nejlepších mravenců a~váhu příspěvku určí podle pořadí. Algoritmus
\textbf{MAX--MIN AS} drží feromon v~intervalu $[\tau_{\min}, \tau_{\max}]$, aktualizuje jen
podle nejlepšího řešení a~inicializuje na $\tau_{\max}$, čímž brání předčasné stagnaci.
Existuje i~ACO pro spojitou optimalizaci, kde se prostor diskretizuje na podintervaly
a~feromon vede hledání do lepších podprostorů.

\paragraph{Vlastnosti a~aplikace.} ACO je decentralizované, řešení v~něm emerguje
a~komunikace je nepřímá, takže má blízko k~reaktivním architekturám agentů, jmenovitě
k~Brooksově subsumpční architektuře. Velkou předností je schopnost \textbf{optimalizovat
v~dynamickém prostředí}: změní-li se graf, feromony se změně samy přizpůsobí. Používá se na
TSP, směrování vozidel, směrování v~sítích, sekvenování DNA, skládání proteinů, detekci hran
v~obraze a~rozvrhování.

\subsection{Další rojové algoritmy}

\noindent\mimo{}

Rojových metod vznikly desítky a~každá si bere metaforu od jiného živočicha. Za pozornost
stojí zejména tyto:
\begin{itemize}
\item \textbf{Umělá včelí kolonie} (\emph{artificial bee colony}, ABC;
  Karaboga 2005) chápe populaci zdrojů potravy jako populaci řešení a~rozděluje včely do tří
  rolí. \emph{Dělnice} prohledávají okolí svého zdroje. \emph{Pozorovatelky} si zdroj vyberou
  ruletou podle kvality a~teprve pak prohledávají jeho okolí. \emph{Průzkumnice} řeší
  stagnaci: zdroj, který se $limit$ pokusů nezlepšil, se opustí a~nahradí náhodným, což je
  vestavěný restart.
\item \textbf{Firefly algorithm} (Yang, 2008) pracuje se světluškami, jejichž jasnost je
  úměrná fitness a~se vzdáleností klesá podle
  $I(r) = I_0 e^{-\gamma r^2}$. Méně jasná světluška se posouvá směrem
  k~jasnější:
  $\vec x_i \gets \vec x_i + \beta_0 e^{-\gamma r_{ij}^2}(\vec x_j - \vec x_i)
  + \alpha \vec\varepsilon$. Protože přitažlivost má omezený dosah, roj se přirozeně dělí na
  podskupiny, a~algoritmus se proto hodí na multimodální úlohy.
\item \textbf{Bat algorithm}, \textbf{cuckoo search}, \textbf{grey wolf
  optimizer} a~desítky dalších \uv{metaforických} algoritmů se objevují stále. Sklízejí ovšem
  kritiku, že mnohé z~nich jsou jen přejmenované varianty PSO či ES a~přinášejí spíše
  terminologický zmatek než nové principy.
\end{itemize}

\subsection{Shrnutí ke~zkoušce}

\begin{itemize}
\item Jádrem PSO jsou pohybové rovnice
  $\vec v \gets w\vec v + c_1r_1(\vec p - \vec x) + c_2r_2(\vec g - \vec x)$
  a~$\vec x \gets \vec x + \vec v$, kde se skládá setrvačnost $w$ (klesá 0{,}9 $\to$ 0{,}4)
  se složkou kognitivní a~sociální při $c_1=c_2=2$. Přestřelení cíle se krotí omezením
  $v_{\max}$ nebo constriction faktorem $\chi \approx 0{,}729$ a~roj může být propojen
  topologií gbest, nebo lbest (kruh, mřížka).
\item Proti GA se PSO liší tím, že nemá žádné genetické operátory, částice si pamatují
  $\vec p_i$ a~$\vec g$ a~informace teče jednosměrně od lepších k~horším.
\item V~ACO je pravděpodobnost přechodu úměrná
  $p_{ij} \propto \tau_{ij}^{\alpha}\eta_{ij}^{\beta}$ s~viditelností $\eta = 1/d$
  a~feromon se aktualizuje jako $\tau \gets (1-\rho)\tau + \sum_k Q/L_k$, tedy odpařováním
  a~uložením úměrným kvalitě trasy.
\item Vyplatí se umět odlišit původní AS, kde aktualizují všichni mravenci, od ACS
  s~globálním updatem jen podle nejlepšího, lokálním updatem při průchodu a~pseudonáhodným
  pravidlem s~parametrem $q_0$, a~od MAX--MIN AS, který drží meze feromonu. Akce démona jsou
  centralizovaný krok, prakticky lokální prohledávání.
\item Z~dalších metod stojí za zapamatování ABC s~rolemi dělnic, pozorovatelek
  a~průzkumnic, kde limit nezlepšení funguje jako restart, a~firefly algorithm, v~němž
  přitažlivost klesá exponenciálně se vzdáleností.
\item Klíčová slova celé kapitoly jsou stigmergie, emergence a~decentralizace; mechanismus
  stojí na kladné zpětné vazbě doplněné odpařováním a~díky tomu se metody hodí i~do
  dynamického prostředí.
\end{itemize}

%=====================================================================

\section{Lokální prohledávání a~memetické algoritmy}
\label{sec:lokalni}
%=====================================================================

\subsection{Lokální prohledávání a~horolezecký algoritmus}

Lokální metody pracují s~jediným řešením, které postupně vylepšují drobnými změnami.
Aby to bylo možné, musíme nejdřív říct, které změny jsou ještě drobné. K~tomu slouží
pojem okolí.

\begin{defbox}[Okolí a~lokální prohledávání]
Pro prostor řešení $X$ definujeme \emph{okolí} $N(x) \subseteq X$ jako množinu řešení,
do nichž se z~$x$ dostaneme jedinou elementární změnou; podle úlohy to bývá překlopení
bitu, prohození dvou měst nebo krok 2-opt. \emph{Lokální prohledávání} pak iterativně
přechází z~aktuálního řešení do některého souseda podle zvoleného pravidla.
\emph{Lokální optimum} vzhledem k~$N$ je takové $x$, že žádný jeho soused není lepší.
Tato vlastnost tedy závisí na definici okolí: co je lokálním optimem v~jednom okolí,
nemusí jím být v~jiném!
\end{defbox}

Nejjednodušší lokální metodou je horolezecký algoritmus, který se od nejbližších
příbuzných liší jen tím, jak souseda vybírá a~kdy jej přijme.

\begin{defbox}[Horolezecký algoritmus (\emph{hill climbing})]
\textbf{1. Steepest ascent (nejstrmější stoupání):} algoritmus projde \emph{celé} okolí
a~přejde do nejlepšího souseda, je-li lepší než $x$; jinak skončí.\\
\textbf{2. First-improvement (první zlepšení):} okolí se prochází jen do prvního
nalezeného lepšího souseda, do něhož se rovnou přejde. Je to levnější a~často stejně
dobré.\\
\textbf{3. Stochastic hill climbing:} vybere se náhodný soused a~přijme se, je-li lepší;
v~jiné variantě se přijímá s~pravděpodobností úměrnou zlepšení.\\
\textbf{4. Random-restart hill climbing:} po uvíznutí se začne z~nové náhodné pozice
a~nejlepší dosud nalezené řešení si algoritmus pamatuje. S~rostoucím počtem restartů
konverguje pravděpodobnost nalezení globálního optima k~1.
\end{defbox}

Slabiny má horolezec tři a~všechny plynou z~toho, že se dívá jen o~krok dopředu. Uvízne
v~lokálním optimu, zastaví se na \emph{plošinách} (plateau, kde jsou všichni sousedé
stejní a~pomáhají \emph{sideways moves}) a~má potíže na hřebenech. Proti tomu stojí jeho
přednosti: je jednoduchý, vystačí s~malou pamětí a~je rychlý. Právě proto se hodí jako
vylepšovací procedura uvnitř evolučního algoritmu.

\subsection{Simulované žíhání}

Nejpřímočařejší způsob, jak horolezce z~lokálního optima vyvést, je dovolit mu občas
krok dolů. Otázka pak zní, jak často a~jak hluboko.

\begin{defbox}[Simulované žíhání (\emph{simulated annealing}, SA)]
Metodu navrhli Kirkpatrick, Gelatt a~Vecchi (1983) na základě fyzikální analogie
s~pomalým chladnutím kovu, při němž atomy zaujmou uspořádání s~minimální energií.
Algoritmus je horolezec, který \textbf{připouští i~zhoršující kroky}, a~to
s~pravděpodobností klesající s~velikostí zhoršení a~s~teplotou.
\end{defbox}

Konkrétní podobu té pravděpodobnosti dává Metropolisovo kritérium.

\begin{defbox}[Metropolisovo kritérium]
Pro minimalizaci, kandidáta $y \in N(x)$ a~$\Delta = f(y) - f(x)$ platí:
\[
P(\text{přijmout } y) =
\begin{cases}
1, & \Delta \le 0 \quad (\text{zlepšení}),\\[2pt]
e^{-\Delta / T}, & \Delta > 0 \quad (\text{zhoršení}),
\end{cases}
\]
kde $T > 0$ je teplota.
\end{defbox}

\textbf{Číselný příklad.} Mějme zhoršení $\Delta = 2$. Při $T = 1$ vyjde
$P = e^{-2} = 0{,}135$, při $T = 0{,}5$ klesne na $P = e^{-4} = 0{,}018$ a~při $T = 5$
naopak stoupne na $P = e^{-0{,}4} = 0{,}67$. Na začátku běhu, kdy je horko, se tedy
algoritmus chová téměř jako náhodná procházka; na konci, kdy je chladno, jako čistý
horolezec.

\begin{algorithm}[!ht]
\caption{Simulované žíhání}
\begin{algorithmic}[1]
\State $x \gets$ počáteční řešení; $T \gets T_0$; $x^* \gets x$
\While{$T > T_{\min}$ a~není splněna ukončovací podmínka}
  \For{$k$ krát (délka Markovova řetězce při dané teplotě)}
    \State vyber náhodně $y \in N(x)$; $\Delta \gets f(y) - f(x)$
    \If{$\Delta \le 0$ nebo $\mathrm{rand}() < e^{-\Delta/T}$} $x \gets y$ \EndIf
    \If{$f(x) < f(x^*)$} $x^* \gets x$ \EndIf
  \EndFor
  \State $T \gets \mathrm{cool}(T)$
\EndWhile
\Ensure $x^*$
\end{algorithmic}
\end{algorithm}

O~tom, jak rychle teplota klesá, rozhoduje rozvrh chlazení.

\paragraph{Rozvrh chlazení (\emph{cooling schedule}).}
\begin{itemize}
\item \emph{Geometrický:} teplota se v~každém kroku vynásobí konstantou, tedy
  $T_{k+1} = \alpha T_k$ pro $\alpha \in [0{,}8; 0{,}99]$. V~praxi je to nejběžnější volba.
\item \emph{Lineární:} teplota klesá o~pevný krok podle $T_k = T_0 - \beta k$.
\item \emph{Logaritmický:} $T_k = c/\log(k+1)$. Je to jediný rozvrh, pro který je
  dokázána konvergence.
\item Adaptivní rozvrhy se řídí průběhem běhu; patří k~nim \emph{reheating}, kdy se při
  stagnaci teplota znovu zvýší.
\end{itemize}

\begin{veta}[Konvergence SA]
Při logaritmickém rozvrhu $T_k \ge c/\log(k+1)$ s~dostatečně velkou konstantou
$c$ (souvisí s~výškou nejhlubší \uv{bariéry} v~krajině) konverguje simulované
žíhání k~množině globálních optim s~pravděpodobností 1.
\end{veta}
Prakticky je ovšem takový rozvrh nepoužitelně pomalý. Simulované žíhání je proto
teoreticky uspokojivé, ale v~praxi se sahá po rychlejších heuristických rozvrzích.

Formálně je SA nehomogenní Markovův řetězec. Při pevné $T$ konverguje k~Boltzmannovu
rozdělení $\pi(x) \propto e^{-f(x)/T}$, které se pro $T \to 0$ koncentruje na globálních
optimech. Tatáž myšlenka se objevuje i~v~evolučních algoritmech: \emph{Boltzmannův
turnaj} (odst.~\ref{ssec:selekce}) ji přenáší do selekce.

\subsection{Tabu prohledávání}

Simulované žíhání se z~lokálního optima dostává náhodou. Tabu prohledávání volí opačnou
cestu: cíleně si zakáže tahy, které by je vrátily tam, odkud právě přišlo.

\begin{defbox}[Tabu search (Glover, 1986)]
Lokální prohledávání doplněné o~\emph{krátkodobou paměť}. Uchovává se \emph{tabu list}
nedávno provedených tahů, což je efektivnější než ukládat celá řešení, a~tyto tahy jsou
po dobu zvanou \emph{tabu tenure} zakázány. V~každém kroku algoritmus přejde do
\textbf{nejlepšího nezakázaného souseda}, a~to i~tehdy, je-li horší než aktuální řešení.
Právě tím se dostane z~lokálního optima a~zároveň necyklí.
\end{defbox}

\begin{itemize}
\item \textbf{Tabu tenure} $T$ může být statická, tedy pevný počet iterací, nebo
  dynamická, kdy se losuje z~nějakého rozsahu.
\item \textbf{Aspirační kritérium} dovoluje tabu porušit, vede-li tah k~řešení lepšímu
  než dosud nejlepší nalezené. Bez něj by tabu mohlo blokovat postup.
\item \textbf{Střednědobá paměť} jsou pravidla, která hledání směrují do slibných
  oblastí, tedy zajišťují intenzifikaci.
\item \textbf{Dlouhodobá paměť}, jinak též frekvenční, je počítadlo toho, jak často byl
  který prvek řešení použit. Často navštěvované konfigurace se penalizují, což slouží
  k~diverzifikaci, případně se restartuje z~málo navštívené oblasti.
\end{itemize}

Tabu search se používá na TSP (výměny hran, \emph{ejection chains}), na rozvozní úlohy,
na sekvenování DNA a~na rozvrhování. Mluví pro něj únik z~lokálních optim, absence
cyklení a~to, že funguje v~diskrétních i~spojitých doménách. Proti němu stojí velký počet
parametrů a~vyšší výpočetní náročnost daná tím, že se prochází celé okolí.

\subsection{Scatter search}

Scatter search je populační metaheuristika zaměřená na \textbf{diverzitu}. Udržuje malou
\emph{referenční množinu} $R$ o~řádově desítkách jedinců a~vybírá do ní podle kombinace
kvality a~vzájemné vzdálenosti: u~nově přidávaného jedince se maximalizuje jeho minimální
vzdálenost od zbytku $R$.

Cyklus má pět kroků. Začíná (1)~\emph{diverzifikací}, tedy chytrou inicializací. Pak
přijde (2)~\emph{generování podmnožin}, typicky všech dvojic z~$R$, a~jejich
(3)~\emph{kombinace} aritmetickým či permutačním křížením. Výsledek projde
(4)~\emph{vylepšením} lokálním prohledáváním, mutací nebo tabu search a~nakonec se
(5)~\emph{aktualizuje $R$}.

Aktualizovat referenční množinu lze třemi způsoby. \emph{Inkrementálně}, kdy za generaci
přibude jeden či několik jedinců; \emph{úplnou obnovou} každou generaci; nebo \emph{průběžně}
už ve vnitřním cyklu, kdy se noví jedinci do $R$ zařazují okamžitě a~rekombinuje se
rovnou s~nimi. Metoda se hodí pro úlohy s~velmi drahou fitness, například pro black-box
optimalizátor OptQuest.

\subsection{Memetické algoritmy}

Všechny dosavadní lokální metody mají společné to, že jeden běh nijak netěží ze
zkušenosti běhu předchozího. Memetické algoritmy tuhle mezeru zaplňují tím, že lokální
prohledávání zasadí dovnitř evoluce.

\begin{defbox}[Mem a~memetický algoritmus]
\emph{Mem} (Dawkins, 1976) je jednotka kulturní informace, která se šíří napodobováním;
bývá popsán jako \uv{gen kultury}. \emph{Memetický algoritmus} (Moscato, 1989) je
hybridní metoda: \textbf{evoluční algoritmus s~vnořeným lokálním prohledáváním}, které
vylepšuje jednotlivé jedince, a~vykládá se proto jako \uv{individuální učení}. Setkáme
se i~se synonymy hybridní EA, genetické lokální prohledávání, lamarckovské či
baldwinovské EA a~kulturní algoritmy.
\end{defbox}

\begin{algorithm}[!ht]
\caption{Memetický algoritmus}
\begin{algorithmic}[1]
\State inicializuj populaci
\While{není splněna ukončovací podmínka}
  \State ohodnoť jedince v~populaci
  \State vytvoř novou populaci stochastickými operátory (selekce, křížení, mutace)
  \State vyber podmnožinu $O$ jedinců, kteří projdou vylepšením
  \For{každého $x \in O$}
    \State proveď individuální učení $x$ pomocí memu (lokálního prohledávače),
      s~pravděpodobností $p$ po dobu $t$
    \State výsledek zpracuj \textbf{lamarckovsky} nebo \textbf{baldwinovsky}
  \EndFor
\EndWhile
\end{algorithmic}
\end{algorithm}

Podstatou metody je lokální prohledávání vnořené do evolučního cyklu. Genetické operátory
přenášejí informaci \emph{mezi} jednotlivými běhy lokálního prohledávače, takže celkové
hledání je mnohem efektivnější než opakovaný restart samotné lokální metody
(Krasnogor \& Smith). V~roli memu může vystupovat hill climbing, simulované žíhání, tabu
search, kombinatorická heuristika jako 2-opt, ale i~gradientní metoda, typicky zpětné
šíření chyby při učení neuronové sítě.

\begin{defbox}[Lamarckismus vs.\ baldwinismus]
Co udělat s~vylepšeným jedincem $x' = \mathrm{LS}(x)$?
\begin{itemize}
\item \textbf{Lamarckovsky:} do populace se vloží $x'$ i~s~jeho fitness, takže
  \emph{genotyp se změní podle fenotypu}. Z~darwinistického hlediska je to
  \uv{nekošer}, prakticky je to ale rychlejší.
\item \textbf{Baldwinovsky:} jedinci $x$ se přiřadí \emph{fitness} vylepšeného $x'$, ale
  genotyp $x$ zůstane \emph{nezměněn}. To je darwinisticky korektní a~fitness se pak
  interpretuje jako \uv{potenciál} jedince. V~tom spočívá \emph{Baldwinův efekt}:
  schopnost učit se jedince zvýhodňuje a~evoluce se totéž posléze \uv{doučí} sama,
  což se označuje jako genetická asimilace.
\end{itemize}
Lamarckismus konverguje rychleji, zato více snižuje diverzitu a~může zkreslit krajinu.
Baldwinismus diverzitu zachovává a~fitness krajinu vyhlazuje, je však pomalejší.
\end{defbox}

\paragraph{Návrhové otázky.} Při návrhu memetického algoritmu se rozhoduje o~třech
věcech. Které jedince vylepšovat: všechny, jen elitu, nebo náhodný podíl $p$? Jak dlouho,
tedy kolik kroků $t$? A~jak silně, tedy do jaké hloubky lokální prohledávání pustit?
Přílišné lokální prohledávání znamená zbytečné náklady a~ztrátu diverzity, příliš malé
nepřinese žádný užitek.

\paragraph{Memetické počítání (\emph{memetic computing}).} Zobecnění, v~němž memy nejsou
známy předem, ale jsou samy předmětem hledání.
\begin{itemize}
\item \textbf{Multi-mem MA:} mem je zakódován \emph{jako součást genotypu} jedince.
  Dědí se i~kříží (potomek zdědí mem lepšího rodiče, nebo se mem vybere náhodně)
  a~používá se k~vylepšení právě tohoto jedince.
\item \textbf{Meta-lamarckovský MA:} memy mají \emph{vlastní populaci} a~koevolvují
  s~jedinci. Při vylepšení se jedinci přiřadí mem, úspěch jedince je pro mem odměnou
  a~nad memy běží druhý evoluční algoritmus.
\item \emph{Memetický prostor} je protějškem prostoru fenotypů. Memy v~něm lze vybírat
  heuristikami, nechat je soutěžit podle úspěšnosti, dělit do subpopulací nebo je řídit
  meta-memetickým algoritmem.
\end{itemize}

\subsection{Ladění a~řízení parametrů}

S~memetickými algoritmy úzce souvisí otázka, odkud se berou hodnoty parametrů. Evoluční
algoritmus jich má mnoho: velikost populace, $p_C$, $p_M$, velikost turnaje, míru
elitismu. Nejsou na sobě nezávislé a~vyčerpávající prohledání všech kombinací není možné.
Podle toho, kdy se o~nich rozhoduje, se rozlišuje ladění od řízení.
\begin{itemize}
\item \textbf{Ladění} (\emph{tuning}) určí parametry \emph{před} během. Používá se
  pokus-omyl, grid search, meta-EA nebo iterované závodění (irace).
\item \textbf{Řízení} (\emph{control}) mění parametry \emph{během} běhu a~má tři podoby:
  \begin{itemize}
  \item \emph{Pevná strategie} (deterministická) závisí jen na čase, například
    $\sigma(t) = 1 - 0{,}9\,t/T$ nebo pokles $p_M$ o~0{,}1 \% za generaci; její
    stochastickou variantou je simulované žíhání.
  \item \emph{Adaptivní} strategie se řídí zpětnou vazbou z~běhu, tedy kvalitou řešení,
    rozptylem fitness či diverzitou. Patří sem pravidlo 1/5, hypermutace při poklesu
    fitness a~adaptivní zjemňování reprezentace, kdy se zvyšuje počet bitů na reálné
    číslo, klesne-li diverzita měřená Hammingovou vzdáleností nejlepšího a~nejhoršího
    jedince.
  \item \emph{Samoadaptivní} strategie dělá z~parametrů součást genomu a~nechává je
    evolvovat; tak pracují ES, EP i~multi-mem MA.
  \end{itemize}
\item Řídit lze i~\emph{velikost populace}, byť jen nepřímo: každému novému jedinci se
  přidělí \uv{délka života} úměrná jeho fitness a~po jejím vypršení jedinec zemře. Lepší
  jedinci žijí déle, takže velikost populace je odvozená dynamická veličina.
\end{itemize}

\subsection{Shrnutí ke~zkoušce}

\begin{itemize}
\item Horolezecký algoritmus má čtyři varianty: steepest ascent, first improvement,
  stochastickou a~s~náhodnými restarty. Naráží na lokální optima, plošiny a~hřebeny.
\item Simulované žíhání přijímá kroky podle Metropolisova kritéria
  $P = \min(1, e^{-\Delta/T})$. Z~rozvrhů chlazení se v~praxi používá geometrický,
  zatímco logaritmický je ten, pro který je dokázána konvergence; pro pevné $T$ řetězec
  konverguje k~Boltzmannovu rozdělení $\propto e^{-f/T}$.
\item Tabu search si drží tabu list tahů, jejich tenure, aspirační kritérium
  a~střednědobou i~dlouhodobou (frekvenční) paměť. Vždy jde do nejlepšího nezakázaného
  souseda, i~když je horší než aktuální řešení.
\item Scatter search udržuje referenční množinu kombinující kvalitu a~diverzitu a~pracuje
  s~aritmetickými kombinacemi doplněnými o~vylepšení.
\item Memetický algoritmus je evoluční algoritmus doplněný o~lokální prohledávání nad
  jedinci. \textbf{Lamarckismus} mění genotyp a~je rychlejší, \textbf{baldwinismus} mění
  jen fitness a~zachovává diverzitu. Memetic computing jde dál a~umisťuje memy do genomu
  nebo do koevolvující populace.
\item U~parametrů se rozlišuje ladění před během od řízení během běhu, přičemž strategie
  řízení mohou být pevné, adaptivní, nebo samoadaptivní.
\end{itemize}


%=====================================================================

\section{Aplikace evolučních algoritmů}
\label{sec:aplikace}
%=====================================================================

Poslední věta okruhu vyjmenovává čtyři aplikační oblasti a~každá z~nich tvoří jeden oddíl
této kapitoly. Přestože se úlohy liší, mají jedno společné: o~úspěchu v~nich rozhoduje
\textbf{volba reprezentace a~operátorů}, nikoli volba \uv{lepšího} evolučního algoritmu.
Právě proto stojí za to procházet aplikace jednu po druhé a~sledovat u~každé, co je v~ní
jedincem a~jak se s~ním smí zacházet.

\subsection{Evoluce pravidlových a~expertních systémů}
\label{sec:lcs}

\subsubsection{Michigan vs.\ Pittsburgh}

Úloha zní: naučit se celou \textbf{sadu pravidel} tvaru \texttt{IF podmínka THEN akce},
a~to buď z~trénovacích dat, nebo z~interakce s~prostředím. Takový výstup se hodí do dobývání
znalostí a~do expertních systémů, ale stejně dobře poslouží při řízení robotů a~ve
zpětnovazebním učení. Zbývá rozhodnout jedinou, zato zásadní věc: co bude jedincem. Odpovědi
jsou dvě a~vedly ke~dvěma základním přístupům, které se liší prakticky ve všem ostatním:

\begin{center}
\begin{tabular}{@{}p{2.8cm}p{6.2cm}p{6.2cm}@{}}
\toprule
 & \textbf{Michigan} (Holland) & \textbf{Pittsburgh} (Smith) \\
\midrule
jedinec & \textbf{jedno pravidlo} & \textbf{celá sada pravidel} \\
řešení & celá populace dohromady & jeden jedinec \\
fitness & síla/přesnost pravidla (nutné rozdělení odměny) & kvalita celého
  systému (přímočará) \\
operátory & standardní, evoluce probíhá jen občas a~na části populace &
  složité, desítky operátorů nad sadami i~jednotlivými pravidly \\
výhody & on-line učení, malý prostor, přirozené pro RL & jednoznačná fitness,
  žádný credit assignment \\
nevýhody & přiřazení zásluh, pravidla musí spolupracovat, ale zároveň soutěží &
  velcí jedinci, drahé vyhodnocení, riziko bloatu \\
\bottomrule
\end{tabular}
\end{center}

\subsubsection{Learning classifier systems (Michigan)}

Michiganskou větev reprezentují \emph{learning classifier systems}, tedy systémy, v~nichž se
znalost neschovává do jednoho jedince, ale do celé populace.

\begin{defbox}[LCS (\emph{learning classifier system})]
Systém, kde populace jednoduchých pravidel (\emph{klasifikátorů}) tvoří
dohromady expertní či řídicí systém. Klasifikátor má levou stranu
(podmínku) jako řetězec nad $\{0,1,*\}$ porovnávaný se vstupem, pravou stranu
(akci či klasifikační třídu) a~\textbf{váhu / sílu} odrážející jeho úspěšnost,
která slouží jako fitness. Evoluce neprobíhá generačně, ale průběžně a~jen na
části populace.
\end{defbox}

\paragraph{Architektura LCS.} Klíčové je uvědomit si, že \emph{expertním
systémem je celá populace}, nikoli jedinec. Systém pracuje ve smyčce, která má šest kroků:
\begin{enumerate}
\item \textbf{Detektory} sejmou stav prostředí a~převedou jej na vstupní zprávu.
\item Zpráva se uloží do \textbf{bufferu zpráv} a~porovná se s~levými stranami
  všech pravidel; ta pravidla, která jí vyhoví, tvoří dohromady \emph{match set}.
\item Mezi vyhovujícími pravidly proběhne \textbf{aukce/soutěž} o~právo jednat,
  v~níž rozhoduje síla pravidla.
\item Vítězné pravidlo zapíše svou pravou stranu buď do bufferu jako vnitřní
  zprávu, nebo ji přes \textbf{efektory} pošle do prostředí jako akci.
\item Prostředí vrátí \textbf{odměnu}, kterou mechanismus \emph{bucket brigade}
  rozdělí mezi pravidla, jež k~ní vedla.
\item Občas a~jen na části populace se spustí \textbf{GA}, který pravidla obměňuje.
\end{enumerate}

\paragraph{Problém reaktivity.} Čistě reaktivní systém nemá vnitřní paměť,
takže se nedokáže rozhodovat podle ničeho jiného než podle okamžitého vjemu.
Holland proto zavedl \emph{zprávy} (\emph{messages}): pravá strana pravidla
může místo akce zapsat vnitřní zprávu do bufferu a~levá strana má
\emph{receptory}, jimiž ji zachytí. Pravidla se tak dají řetězit za sebe
a~systém tím získá vnitřní stav, a~s~ním i~výpočetní sílu. Za to ovšem platí
novým problémem: jak rozdělit jedinou odměnu mezi celý řetězec pravidel, který
k~ní vedl.

\begin{defbox}[Bucket brigade (algoritmus \uv{požárního řetězu})]
Mechanismus rozdělení odměny v~LCS. Pravidlo, které chce být aplikováno,
musí \uv{zaplatit} část své síly pravidlu, jež ho aktivovalo (jako by kupovalo
právo jednat). Odměna z~prostředí připadne poslednímu článku řetězu a~postupně
\uv{proteče} zpět k~pravidlům, která k~ní vedla. Prakticky je obtížné
vyvážit tuto ekonomiku pravidel, dnes se proto téměř nepoužívá.
\end{defbox}

\paragraph{ZCS (Wilson, 1994) --- \uv{zero-level} LCS.} Wilson šel opačnou cestou než
Holland a~systém radikálně zjednodušil: žádné vnitřní zprávy, žádná složitá redistribuce
odměny. Zůstane tohle:
\begin{itemize}
\item Populaci pravidel označíme $P$, vstup přicházející z~detektorů je $x$.
\item Do \emph{match setu} $M$ patří ta pravidla, jejichž podmínka odpovídá vstupu $x$.
\item Akce $a$ se z~$M$ vybere ruletou podle síly pravidel; pravidla, která prosazují
  právě zvolenou akci, tvoří \emph{action set} $A \subseteq M$.
\item \textbf{Posílení:} z~pravidel v~$A$ se nejdřív odebere podíl $\beta$ jejich síly
  do \uv{kbelíku} $B$. Jakmile dorazí odměna $r$, přičte se každému pravidlu v~$A$ hodnota
  $\beta r/|A|$ a~pravidlům z~předchozího action setu $A_{-1}$ hodnota
  $\gamma B/|A_{-1}|$. Pravidla z~$M \setminus A$, tedy ta, která vstupu odpovídala,
  ale prosazovala jinou akci, jsou zdaněna.
\item \textbf{Cover operátor:} pokud pro aktuální vstup neexistuje žádné vyhovující
  pravidlo, vygeneruje se ad hoc --- do podmínky se náhodně přidají hvězdičky a~akce se
  zvolí náhodně.
\end{itemize}

\paragraph{XCS (Wilson, 1995) --- fitness založená na přesnosti.}
ZCS má tři slabiny. Neevolvuje \emph{úplný} systém pravidel, který by pokrýval všechny
situace. Pravidla stojící na začátku řetězců jsou odměňována jen zřídka, a~proto vymírají.
A~vymírají i~pravidla, která vedou k~malým, zato správně předpovězeným odměnám. Společnou
příčinou je, že síla plnila dvě role najednou: byla fitness pro GA a~zároveň kritériem
výběru akce. Wilsonova myšlenka rozdělila obojí. \emph{Predikce} má odhadovat odměnu, kdežto
\emph{evoluce} má preferovat \textbf{spolehlivé (přesné)} klasifikátory, tedy ty, jejichž
predikce se trefuje, i~když slibuje málo.

\begin{defbox}[XCS --- klasifikátor jako pětice]
$(c, a, p, \epsilon, F)$: podmínka, akce, predikce odměny $p$, chyba predikce
$\epsilon$, fitness $F$. Přesnost se počítá jako klesající funkce chyby,
$\kappa = \alpha\,(\epsilon/\epsilon_0)^{-\nu}$ (pro $\epsilon > \epsilon_0$,
jinak $\kappa = 1$; $\alpha, \epsilon_0, \nu$ jsou parametry);
fitness je \emph{relativní} přesnost v~rámci action setu.
Predikce a~chyba se aktualizují pravidly Q-learningu.
\end{defbox}

Běh XCS pak vypadá takto. Podle vstupu se sestaví match set $M$ a~ten se rozdělí na action
sety $A_i$ podle jednotlivých akcí. Pro každou akci se předpoví výplata jako průměr predikcí
vážený fitness. Akce se zvolí buď maximem, což je exploatace, nebo ruletou, což je explorace,
a~provede se. Odměna se rozdělí do action setu předchozího kroku: nejprve se aktualizuje
predikce $p$ pomocí $r + \gamma \max_a p_a$, pak chyba $\epsilon$ a~nakonec fitness $F$. GA je
zde \textbf{nikový} --- běží pouze uvnitř action setu, nikoli nad celou populací.

Výsledek stojí za to shrnout. XCS propojuje LCS se zpětnovazebním učením, protože roli
mechanismu pro přiřazení zásluh přebírá Q-learning. Evolvuje \textbf{úplnou a~minimální} mapu
vstup--výstup, takže výstupem není jen optimální politika, ale i~kompaktní a~srozumitelný
popis chování. A~má za sebou řadu praktických aplikací v~dobývání znalostí, v~řízení
i~v~bioinformatice.

\paragraph{Dnešní rodina LCS.} Varianty se dnes rozlišují podle toho, na jakou úlohu míří.
\emph{XCS} slouží pro zpětnovazební učení a~predikuje odměnu.
\textbf{UCS} (\emph{sUpervised Classifier System}) je varianta pro učení
s~učitelem: správná třída je známa, takže se místo predikce odměny rovnou měří přesnost
klasifikace. \emph{XCSF} míří na aproximaci spojitých funkcí a~pravá strana pravidla v~něm
není konstanta, nýbrž lineární model. Podmínky se navíc dnes běžně zapisují nejen nad
$\{0,1,*\}$, ale i~intervaly pro reálné atributy. Tím se z~LCS stává plnohodnotný
nástroj dobývání znalostí, a~to s~jednou předností navíc: výsledkem je
\textbf{čitelná sada pravidel}.

\subsubsection{Pittsburghský přístup a~systém GIL}

V~pittsburghském pojetí je jedincem celá sada pravidel, tedy rovnou kompletní klasifikační
systém. Hodnocení je proto složitější, protože se musí vypořádat s~prioritami pravidel,
s~konflikty mezi nimi a~s~falešně pozitivními i~negativními odpověďmi. Operátorů je celá
řada a~pracují na několika úrovních. Důraz se přitom klade na bohatou doménovou reprezentaci,
tedy na množiny, výčty a~intervaly.

Klasickým představitelem je \textbf{GIL}. Řeší binární klasifikaci a~jedinec v~něm klasifikuje
implicitně do jedné třídy, takže pravidla vůbec nemají pravou stranu.
Struktura jedince je třípatrová: jedinec je \emph{disjunkce komplexů}, komplex je
\emph{konjunkce selektorů} (po jednom pro každou proměnnou) a~selektor je \emph{disjunkce
hodnot} z~domény dané proměnné, zakódovaná bitovou maskou. Například
\[
\big((X{=}A_1) \wedge (Z{=}C_3)\big) \vee \big((X{=}A_2) \wedge (Y{=}B_2)\big)
\;\equiv\;
[\,001|11|0011 \;\vee\; 010|10|1111\,].
\]
Každému patru odpovídají vlastní operátory:
\begin{itemize}
\item na úrovni \emph{jedince} se pravidla vyměňují, kopírují, generalizují, specializují
  a~mažou, případně se do jedince zahrne pozitivní příklad;
\item na úrovni \emph{komplexu} se komplex rozdělí na jednom selektoru, selektor se
  generalizuje (nahradí se $11\ldots1$) nebo specializuje, případně se zahrne negativní
  příklad;
\item na úrovni \emph{selektoru} jde o~mutaci $0 \leftrightarrow 1$, rozšíření $0 \to 1$
  a~zúžení $1 \to 0$.
\end{itemize}
Tyto operátory přímo odpovídají operacím induktivního učení, tedy generalizaci
a~specializaci. GIL je proto vlastně evolučně řízený induktivní algoritmus.

\subsubsection{Připomenutí: metriky klasifikace}

Kvalitu naučeného pravidlového systému je čím měřit. Z~matice záměn (TP, FP, FN, TN) se
počítá
$\text{přesnost} = (TP+TN)/(TP+TN+FP+FN)$,
$\text{senzitivita} = TP/(TP+FN)$ a~$\text{specificita} = TN/(TN+FP)$.
Fitness pravidlového systému bývá kombinací přesnosti, pokrytí a~jednoduchosti, kterou měří
počet pravidel. Jde tedy přirozeně o~vícekriteriální úlohu, i~když se obvykle řeší skalarizací.

\subsection{Neuroevoluce}
\label{sec:neuro}

Druhou velkou aplikační oblastí je učení neuronových sítí evolucí. Vymezení pojmu se cituje
takto:

\begin{defbox}[Neuroevoluce (\emph{neuroevolution})]
\uv{Neuroevoluce je metoda pro úpravu vah, topologií nebo ansámblů neuronových
sítí za účelem naučení konkrétní úlohy. Evoluční výpočty se používají
k~hledání parametrů sítě maximalizujících fitness, která měří výkon v~úloze.
Ve srovnání s~jinými metodami učení je neuroevoluce velmi obecná --- umožňuje
učení bez explicitních cílových hodnot, s~nediferencovatelnými aktivačními
funkcemi a~s~rekurentními sítěmi.} (Miikkulainen, 2017)
\end{defbox}

Evolvovat se dá v~podstatě cokoli, co síť určuje: samotné váhy, další parametry jako
aktivační funkce či konstanty učení, architektura ve smyslu topologie a~propojení,
architektura a~váhy zároveň, a~dokonce i~parametry samotného učicího algoritmu. Hlavní
doménou neuroevoluce je \textbf{zpětnovazební učení} a~robotika, tedy úlohy, kde nejsou
k~dispozici cílové výstupy potřebné pro učení s~učitelem.

\subsubsection{Evoluce vah}

Nejjednodušší je evolvovat váhy pevně dané sítě. Váhy se zakódují do vektoru pevné délky
a~pustí se na něj reálný GA, ES nebo DE se standardními operátory; nic dalšího není potřeba
vymýšlet. Zajímavější jsou vlastnosti takového postupu:
\begin{itemize}
\item Pro učení s~učitelem je evoluce pomalejší než specializované gradientní metody, tedy
  než zpětné šíření chyby.
\item Výhodou je naopak snadná paralelizace a~nezávislost na gradientu: postup funguje
  i~pro nediferencovatelné a~rekurentní sítě a~dá se použít pro RL.
\item Naráží ale na \textbf{problém konkurujících konvencí}
  (\emph{competing conventions}), jinak též permutační problém: tatáž funkce sítě odpovídá
  mnoha různým genotypům, protože skryté neurony lze libovolně přeuspořádat. Křížení dvou
  funkčně shodných rodičů proto často vytvoří nefunkčního potomka.
\item V~posledních letech se zájem o~evoluci vah vrátil. Ukazuje se, že je
  konkurenceschopná i~pro velké sítě, pokud se hodnocení provádí na minibatchích.
\end{itemize}

\subsubsection{Evoluce architektury}

Jakmile se místo vah evolvuje architektura, změní se povaha fitness. Fitness architektury je
totiž odhad výkonu sítě, a~aby se dal spočítat, musí se síť postavit, inicializovat a~naučit,
ať už zpětným šířením, nebo jiným EA. Nejlépe navíc vícekrát. Vyhodnocení fitness je proto
drahé a~zašuměné.

\begin{itemize}
\item \textbf{Přímé kódování} (\emph{direct encoding}) popisuje strukturu propojení jako
  binární matici. Jedinec je pak buď linearizovaná matice, na kterou stačí standardní
  operátory, nebo matice sama a~pracuje se s~ní speciálními 2D operátory. Je to jednoduché,
  ale neškáluje, protože velikost genomu roste kvadraticky.
\item \textbf{Gramatické kódování} (Kitano, 1990) matici propojení negeneruje přímo:
  vyrobí ji jednoduchá 2D formální gramatika a~evolvují se \emph{pravidla gramatiky}.
  Kódování je kompaktní, tedy logaritmické, a~umožňuje opakující se motivy. Za to platí tím,
  že je nepřímé a~hůře se ladí.
\item \textbf{Buněčné kódování} (\emph{cellular encoding}, Gruau) jde ještě dál: jedincem je
  \textbf{program v~GP}, který popisuje, jak má síť \emph{vyrůst} z~jediné buňky pomocí
  operací přidej neuron, rozděl neuron sériově, rozděl paralelně, přepoj synapsi a~změň
  váhu. Jde tedy o~vývojový (developmental) přístup.
\item \textbf{Simulovaný růst} spojů ve 2D rovině patřil mezi rané pokusy v~evoluční
  robotice, ale neškáloval.
\end{itemize}

\paragraph{GNARL (Angeline a~kol., 1993).} Systém staví na evolučním programování nad grafy.
Vstupní a~výstupní neurony jsou pevné, kdežto skryté neurony a~spoje se evolvují, a~to
architektura i~váhy \emph{současně}. Mutace jsou dvojího druhu: \emph{strukturální} přidá
neuron bez spojů, přidá spoj s~nulovou váhou, smaže neuron nebo smaže spoj, zatímco
\emph{parametrická} je zaujatá mutace váhy. Sílu mutace řídí \emph{teplota}, což je přístup
převzatý ze simulovaného žíhání; strukturální změny mají zůstat malé a~nepříliš destruktivní.
Rodiči se stává lepší polovina populace a~křížení se nepoužívá, jak je v~EP zvykem.

\subsubsection{NEAT}

Nejvlivnější odpovědí na problém konkurujících konvencí je NEAT.

\begin{defbox}[NEAT (\emph{NeuroEvolution of Augmenting Topologies},
K.~Stanley, 2002)]
Síť je reprezentována jako \textbf{seznam hran}; každý gen-hrana obsahuje:
vstupní a~výstupní neuron, váhu, příznak \emph{enabled/disabled}
a~\textbf{globální inovační číslo} (\emph{innovation number}), které se
přiděluje při prvním vzniku dané strukturální novinky.
\end{defbox}

Celý systém stojí na třech klíčových myšlenkách:
\begin{enumerate}
\item \textbf{Inovační čísla řeší problém konkurujících konvencí.}
  Kříží se pouze geny se \emph{shodným inovačním číslem}, tedy geny se společným evolučním
  původem. Geny, které přebývají u~jednoho rodiče
  (\emph{disjoint} a~\emph{excess}), se přeberou od zdatnějšího z~rodičů.
  Křížení tak strukturu nerozbíjí a~nevzniká z~něj žádné \uv{bezhlavé kuře}.
\item \textbf{Speciace (niching) chrání inovace.} Na vektorech inovačních
  čísel se definuje míra podobnosti
  \[
  \delta \;=\; \frac{c_1 E}{N} + \frac{c_2 D}{N} + c_3 \overline{\Delta W},
  \]
  kde $E$ (\emph{excess}) je počet genů, jejichž inovační číslo leží
  \emph{za} nejvyšším inovačním číslem druhého genomu, $D$ (\emph{disjoint})
  počet genů chybějících u~druhého rodiče \emph{uvnitř} společného rozsahu
  inovačních čísel, $\overline{\Delta W}$ průměrný rozdíl vah genů se
  \emph{shodným} inovačním číslem a~$N$ velikost většího genomu, která slouží
  k~normalizaci. Sítě, jejichž $\delta$ je pod prahem, tvoří jeden \emph{druh}
  a~fitness se počítá jako sdílená, tedy relativní v~rámci druhu. Nová topologie
  tak dostane čas \uv{doladit váhy\,} dřív, než ji vytlačí zavedená řešení.
  Bez toho by strukturální novinka, která fitness zpočátku zhoršuje, okamžitě
  vymřela.
\item \textbf{Postupné narůstání složitosti} (\emph{complexification}) znamená, že
  populace startuje z~minimálních sítí, které mají jen vstupy a~výstupy, a~roste teprve
  mutacemi \emph{přidej spoj} a~\emph{přidej neuron}. Druhá z~nich rozdělí existující spoj:
  původní se zakáže a~vzniknou dva nové. Prohledávání tedy začíná
  v~prostoru nejnižší dimenze a~zvyšuje ji jen tehdy, když se to vyplatí.
\end{enumerate}

\subsubsection{HyperNEAT a~CPPN}

\noindent\mimoq{slidy zmiňují jednou větou, CPPN vůbec}

\textbf{HyperNEAT} rozšiřuje NEAT o~dvě myšlenky. První je \emph{substrát}: váhy sítě se
nekódují jako vektor čísel, ale \emph{generuje je funkce} $S : \mathbb{R}^4 \to \mathbb{R}$,
která dostane souřadnice dvou neuronů v~rovině a~vrátí váhu jejich spojení. Odtud pochází
předpona \uv{hyper}. Druhou myšlenkou je \emph{CPPN}
(\emph{compositional pattern-producing network}), což je způsob, jak funkci $S$ zapsat:
je reprezentovaná sítí ve tvaru DAG, jejíž uzly počítají různé funkce ze slovníku, jako je
sigmoida, gaussovka, sinus či absolutní hodnota. Tuto CPPN pak evolvuje samotný NEAT.
Přínos je v~tom, že reprezentace vah funkcí umí vyjádřit
\textbf{symetrie a~opakující se motivy} a~vygenerovat sítě s~milióny spojů z~malého
genomu.

\subsubsection{Neuroevoluce v~éře hlubokých sítí}

\noindent\mimo{}

Hledání architektury hlubokých sítí dostalo vlastní jméno: \textbf{NAS}
(\emph{neural architecture search}). Popisuje se třemi osami. \emph{Prohledávaný prostor}
může být posloupnost vrstev, DAG opakovaných buněk, nebo podarchitektura jedné velké
\uv{super sítě}. \emph{Prohledávací strategií} bývají evoluční algoritmy, náhodné
prohledávání jako překvapivě silná baseline, bayesovská optimalizace, zpětnovazební učení
nebo gradientní DARTS. \emph{Odhad výkonu} se zlevňuje zkráceným učením a~sdílením vah.
Z~konkrétních systémů stojí za zapamatování \textbf{ENAS}, který sdílí váhy mezi kandidáty,
\textbf{DeepNEAT}, kde je uzlem chromozomu celá vrstva, a~\textbf{CoDeepNEAT}, který
koevolvuje populaci modulů a~populaci \uv{blueprintů}. Samostatnou linií je
\textbf{ES pro RL} (OpenAI 2017): jednoduchá evoluční strategie \emph{pouze s~mutacemi}
nad milióny vah, kterou lze triviálně paralelizovat, protože uzly si mezi sebou vyměňují
jen semínka generátoru náhodných čísel.

\subsection{Kombinatorické úlohy a~práce s~omezeními}
\label{sec:kombi}

Na NP-těžké kombinatorické úlohy jsou evoluční algoritmy oblíbeným nástrojem. Optimum sice
nezaručí, zato v~rozumném čase poskytnou dobré řešení a~snadno se přizpůsobí dodatečným
omezením reálné úlohy, i~když jsou tato omezení sebeošklivější.

\subsubsection{Problém batohu}

\textbf{Zadání.} Batoh má kapacitu $C_{\max}$ a~k~dispozici je $N$ předmětů, z~nichž každý
má svou cenu $v_i$ a~svůj objem $c_i$. Úkolem je maximalizovat $\sum_i x_i v_i$ za podmínky
$\sum_i x_i c_i \le C_{\max}$, kde $x_i \in \{0,1\}$ vyjadřuje, jestli předmět bereme, nebo
ne.

\textbf{Kódování} je triviální. Jedincem je bitová mapa, takže třeba $0110010$ znamená
\uv{vezmi předměty 2, 3 a~6}, a~fungují na ni standardní operátory, tedy křížení a~bit-flip
mutace. \textbf{Problém je fitness}, protože jedinci mohou porušovat kapacitu.
Stojí tu proti sobě dvě protichůdná kritéria: maximalizovat $\sum v_i$ a~zároveň splnit
$\sum c_i \le C_{\max}$.

\subsubsection{Tři obecné strategie pro omezení}

Batoh je tu jen záminkou: potíž s~nepřípustnými jedinci má každá úloha s~omezeními a~odpovědi
na ni jsou tři. Liší se tím, kde se omezení uplatní --- zda až ve fitness, po vygenerování
jedince, nebo už v~samotném kódování.

\begin{defbox}[Práce s~omezeními v~EA]
\textbf{1. Penalizace} (\emph{penalty}): nepřípustná řešení zůstávají
v~populaci, ale jejich fitness se sníží, např.
$f'(x) = \sum_i x_i v_i - \alpha \max\big(0, \sum_i x_i c_i - C_{\max}\big)$.
Volba $\alpha$ je zásadní: příliš malá $\Rightarrow$ populace se plní
nepřípustnými řešeními, příliš velká $\Rightarrow$ ztrácí se možnost projít
nepřípustnou oblastí ke slibnému řešení. Používá se i~dynamická
($\alpha$ roste v~čase) či adaptivní penalizace.\\
\textbf{2. Oprava} (\emph{repair}): nepřípustné řešení se opraví
(u~batohu: vyhazuj předměty s~nejhorším poměrem $v_i/c_i$, dokud se vejde).
Opravu lze provést jen pro výpočet fitness (baldwinovsky) nebo trvale
zapsat do genotypu (lamarckovsky).\\
\textbf{3. Dekodér} (\emph{decoder}) / uzavřené operátory: kódování se navrhne
tak, že \emph{každý} genotyp reprezentuje přípustné řešení. U~batohu: bit
1 znamená \uv{přidej předmět, pokud se ještě vejde}, jinak jej přeskoč.
Elegantní a~účinné; nevýhodou je ne\-jed\-no\-znač\-nost mapování (více
genotypů dává tentýž fenotyp) a~jeho horší lokalita.
\end{defbox}

Nabízí se ještě čtvrtá možnost: chápat porušení omezení jako \textbf{druhé kritérium}
a~sáhnout po vícekriteriálním EA (odst.~\ref{sec:moea}). Výsledkem pak není jediné řešení,
ale celá Pareto fronta kompromisů mezi kvalitou řešení a~mírou porušení omezení.

\subsubsection{Problém obchodního cestujícího}

Úloha TSP zní jednoduše: pro $N$ měst najdi nejkratší okružní jízdu. Fitness je zde
triviální, protože jí je délka cesty, a~o~nastavení evolučního algoritmu se tentokrát mluvit
nemusí. \textbf{Problémem je kódování a~zejména křížení}, protože z~dvojice platných cest se
ledabylým operátorem snadno stane cesta, která žádnou cestou není.

\paragraph{Reprezentace.}
\begin{itemize}
\item \textbf{Sousednostní} (\emph{adjacency}) reprezentace má na pozici $i$ město $j$
  právě tehdy, vede-li hrana $i \to j$. Například $(248397156)$ odpovídá cestě
  $1{-}2{-}4{-}3{-}8{-}5{-}9{-}6{-}7$. Každá cesta má jedinou reprezentaci,
  ale ne každý seznam je platná cesta a~klasické křížení nefunguje.
  Zajímavé je, že schémata tu dávají smysl: $(*3*\ldots)$ jsou třeba všechny cesty
  s~hranou 2--3. Přesto platí: \emph{nepoužívat.}
\item \textbf{Ordinální} (\emph{buffer}) reprezentace si udržuje seznam měst, gen $i$ je
  \emph{pozice} města v~aktuálním seznamu a~použité město se ze seznamu
  vyjme. Cesta $1{-}2{-}4{-}3{-}8{-}5{-}9{-}6{-}7$ se tedy pro seznam
  $(123456789)$ zakóduje jako $(112141311)$. Motivací bylo umožnit standardní
  jednobodové křížení. To sice produkuje platné cesty, jenže výsledek nemá
  s~rodiči nic společného. \emph{Také nepoužívat.}
\item \textbf{Cestová (permutační)} reprezentace popisuje cestu prostě pořadím měst,
  takže cesta $5{-}1{-}7{-}8{-}9{-}4{-}6{-}2{-}3$ se zapíše jako $(517894623)$.
  Je nejpřirozenější,
  ale vyžaduje speciální operátory PMX, OX, CX a~ER (odst.~\ref{sec:reprez}).
\item \textbf{Maticová} reprezentace používá binární matici, kde $1$ na pozici $(i,j)$
  znamená buď hranu $i \to j$, nebo, což je častější, že $i$ předchází $j$. Operátory jsou
  opět speciální: \emph{konjunkce} udělá bitové AND obou rodičů a~náhodně doplní chybějící
  hrany, \emph{disjunkce} rozdělí matici na kvadranty, dva z~nich zahodí, odstraní spory
  a~zbytek doplní náhodně.
\end{itemize}

\paragraph{Inicializace.} Populaci lze naplnit náhodnými permutacemi, nebo sáhnout po
konstrukčních heuristikách. \emph{Nejbližší soused} začne náhodným městem a~vždy jde do
nejbližšího nenavštíveného. \emph{Vkládání hran} k~částečné cestě $T$ vybere nejbližší
město $c$ mimo $T$, najde v~$T$ hranu $k$--$j$ minimalizující
$d(k,c) + d(c,j) - d(k,j)$ a~tuto hranu nahradí dvojicí $k$--$c$, $c$--$j$.

\paragraph{Mutace.} Nejdůležitější je inverze úseku, protože odpovídá 2-opt kroku. Dále se
používá vložení města jinam, posun podcesty, prohození dvou měst, prohození podcest
a~heuristické mutace 2-opt a~3-opt, které vezmou dvě hrany a~přepojí čtyři města jinak.
Kombinace ES či GA s~\uv{chytrými mutacemi} typu 2-opt je v~praxi velmi
účinná; de facto jde o~memetický algoritmus.

\subsubsection{SAT}

\noindent\mimoq{slidy řeší batoh, TSP a~rozvrhování; SAT je z~Michalewicze}

\textbf{Zadání.} Je dána formule $f: \mathbb{B}^n \to \mathbb{B}$ v~CNF a~hledá se
ohodnocení $x$, pro které platí $f(x)=1$. Zatímco 2-SAT je v~P, 3-SAT a~všechno nad ním
je NP-úplné.

\textbf{Fitness.} Samotné $f$ je nepoužitelné, protože nabývá jen hodnot 0 a~1 a~krajina
tak vypadá jako \uv{jehla v~kupce sena}. Standardem je proto \textbf{počet splněných
klauzulí}. Zlepšit se dá vážením problematických klauzulí a~\emph{zjemňující funkcí}, tedy
regularizačním členem, který rozliší jedince se stejným počtem splněných klauzulí a~pomůže
tím z~plošin.

\textbf{Reprezentace} bývá přímo bitový řetězec. Zajímavá je \emph{reálná} varianta:
proměnná $x_j$ se nahradí reálným $y_j$ (pravda $=1$, nepravda $=-1$), literály se přepíší
jako $x \mapsto (1-y)^2$ a~$\neg x \mapsto (1+y)^2$ a~výraz se \emph{minimalizuje}. Hodnota
literálu je tedy \emph{penalta}, a~proto se \textbf{disjunkce} (klauzule) přepíše jako
\textbf{součin} a~\textbf{konjunkce} (formule) jako \textbf{součet}. Klasickou lokální
heuristikou je \textbf{WalkSAT} a~kombinace EA + WSAT je typické memetické řešení.

\subsubsection{Rozvrhování a~plánování výroby}

\textbf{Školní rozvrh.} Jedincem je celý rozvrh zapsaný jako matice, jejíž řádky jsou
učitelé, sloupce třídy a~hodnoty kódy předmětů. Mutace míchá předměty a~křížení vyměňuje
mezi jedinci lepší řádky. Fitness kombinuje kvalitu jednotlivých řádků s~měkkými kritérii,
jako jsou okna a~rovnoměrnost. \textbf{Tvrdá omezení}, například že učitel nemůže učit dvě
třídy naráz, je nutné respektovat \emph{přímo v~operátorech}. Jinak vzniká příliš mnoho
nepřípustných jedinců.

\textbf{Job shop scheduling.} Výrobky se skládají z~dílů, pro každý díl existuje více plánů
výroby na strojích a~přestavení stroje stojí čas; fitness je celkový čas výroby
(\emph{makespan}). Kritické je i~zde kódování. \emph{Permutační} jedinec nese jen pořadí
výrobků a~zbytek dourčí dekodér, což dovolí použít křížení z~TSP, ale ukazuje se to jako
málo efektivní, protože složitou část práce odvede dekodér. \emph{Přímé} kódování celého
plánu je výkonnější, zato vyžaduje specializované operátory.

\subsection{Vícekriteriální optimalizace}
\label{sec:moea}

\subsubsection{Dominance a~Pareto fronta}

Místo jediné fitness teď máme celý vektor kritérií $f_1,\dots,f_k$; pro jednoduchost
budeme všechna minimalizovat. Taková kritéria si obvykle odporují: levnější řešení bývá
horší kvality, přesnější model bývá větší. Jakmile jsou kritéria v~konfliktu, přestává
existovat jediné nejlepší řešení, protože zlepšení v~jednom kritériu se platí zhoršením
v~jiném. Potřebujeme tedy jiné vymezení toho, co znamená \uv{lepší}, než je porovnání
dvou čísel.

\begin{defbox}[Dominance]
Pro dvojici jedinců $x, y$ zavádíme tři vztahy:
\begin{itemize}
\item Jedinec $x$ \textbf{slabě dominuje} $y$ ($x \preceq y$), právě když
  $f_i(x) \le f_i(y)$ pro všechna $i = 1,\dots,k$, tedy není-li $x$ v~žádném
  kritériu horší.
\item Jedinec $x$ \textbf{dominuje} $y$ ($x \prec y$), právě když $x \preceq y$
  a~navíc existuje $j$ s~$f_j(x) < f_j(y)$: v~ničem není horší a~alespoň v~jednom
  kritériu je opravdu lepší.
\item Jedinci $x$ a~$y$ jsou \textbf{neporovnatelné}, pokud neplatí ani $x \prec y$,
  ani $y \prec x$; každý z~nich je pak lepší v~něčem jiném.
\end{itemize}
Relace $\prec$ je \emph{ostré} (striktní) částečné uspořádání: je ireflexivní
a~tranzitivní, ale nikoli úplné. Právě z~té neúplnosti plyne, že optimem obecně není
jediný bod, nýbrž celá množina navzájem neporovnatelných řešení.
\end{defbox}

\begin{defbox}[Pareto optimalita a~Pareto fronta]
Řešení $x^*$ je \emph{Pareto-optimální}, není-li dominováno žádným přípustným řešením.
Množina všech Pareto-optimálních řešení v~prostoru proměnných se nazývá \emph{Pareto
množina} a~její obraz v~prostoru kritérií je \textbf{Pareto fronta}. Cílem
vícekriteriálních EA je tuto frontu populací \emph{aproximovat}.
\end{defbox}

Od algoritmu se proto chtějí dvě věci najednou. \textbf{Konvergence} znamená dostat
populaci co nejblíže ke skutečné frontě, kdežto \textbf{diverzita} žádá pokrýt frontu
rovnoměrně v~celém jejím rozsahu. Právě proto se na tuhle úlohu evoluční algoritmy hodí
lépe než cokoli jiného: pracují s~populací, a~mohou tedy celou aproximaci fronty vrátit
v~jediném běhu.

\textbf{Příklad.} Uvažujme body $A(1,10)$, $B(2,7)$, $C(4,5)$, $D(8,2)$ a~$E(5,8)$
a~minimalizujme obě kritéria. Bod $E$ je dominován bodem $C$, protože platí $4 \le 5$
i~$5 \le 8$ a~alespoň jedna z~nerovností je ostrá; $E$ tedy do fronty nepatří. Zbylé
body $A, B, C, D$ jsou navzájem neporovnatelné a~tvoří Pareto frontu.

\subsubsection{Skalarizace --- jednoduchá řešení}

Nejsnazší způsob, jak si s~více kritérii poradit, je převést je zpátky na jediné číslo
a~pustit na úlohu obyčejný jednokriteriální algoritmus. Podob takového převodu je několik
a~liší se hlavně tím, na které části fronty vůbec dosáhnou.

\begin{itemize}
\item \textbf{Lineární skalarizace} sečte kritéria s~vahami, $f = \sum_i w_i f_i$,
  a~výsledek předá standardnímu jednokriteriálnímu EA (v~kontextu MOEA se označuje jako
  SOEA). Slabiny má tři: neumíme volit váhy, jeden běh dá jediný bod fronty a~zásadně
  \textbf{nekonvexní části fronty nelze lineární skalarizací nikdy nalézt}.
\item \textbf{$\varepsilon$-constraint} minimalizuje jen kritérium $f_1$ a~zbylá kritéria
  převede na omezení $f_i \le \varepsilon_i$ pro $i \ge 2$. Na nekonvexní části fronty už
  dosáhne, zaplatí se za to ale volbou konstant $\varepsilon_i$ a~tím, že se řeší úloha
  s~omezeními.
\item \textbf{Čebyševova (Chebyshev) skalarizace} minimalizuje
  $\max_i \lambda_i |f_i(x) - z_i^*|$, kde $z^*$ je referenční (ideální) bod. Na rozdíl
  od lineární varianty umí dosáhnout na libovolný bod fronty včetně nekonvexních částí.
\end{itemize}

\subsubsection{Historické algoritmy}

\paragraph{VEGA (Schaffer, 1985)} patří mezi vůbec první MOEA. Populace $N$ jedinců se
pro každé kritérium $f_i$ zvlášť seřadí a~vybere se z~ní $N/k$ jedinců nejlepších podle
$f_i$; takto vzniklé podmnožiny se pak sloučí dohromady, zkříží a~zmutují. Slabina je
z~toho vidět hned: každá podpopulace táhne evoluci ke svému vlastnímu extrému, takže se
špatně udržuje diverzita a~celek konverguje ke krajním bodům optimálním pro jednotlivá
kritéria, zatímco \uv{středy} fronty se ztrácejí.

\paragraph{NSGA (1994)} poprvé použil dominanci přímo jako fitness. Populace se rozdělí
do front: $F_1$ jsou nedominovaní jedinci z~$P$, $F_2$ nedominovaní z~$P \setminus F_1$,
$F_3$ nedominovaní z~$P \setminus (F_1 \cup F_2)$ a~tak dál. Jedinci z~první fronty
dostanou \uv{fiktivní} fitness, kterou algoritmus vydělí \emph{nikovým faktorem}
$\sum_j sh(i,j)$, kde
$sh(i,j) = 1 - (d(i,j)/d_{\mathrm{share}})^2$ pro $d(i,j) < d_{\mathrm{share}}$
a~0 jinak. Druhá fronta dostane fitness menší, než má nejhorší jedinec první fronty,
opět dělenou nikovým faktorem, a~stejně se pokračuje dál. Nevýhody jsou tři: je nutné
nastavit $d_{\mathrm{share}}$, chybí elitismus a~výpočetní složitost $O(kN^3)$ je vysoká.

\subsubsection{NSGA-II}

Právě tyto nedostatky odstraňuje nástupce, který se dodnes používá jako standardní
srovnávací algoritmus vícekriteriální optimalizace.

\begin{defbox}[NSGA-II (Deb a~kol., 2000)]
Opravuje nedostatky NSGA a~stojí na třech pilířích.
\textbf{(1) Rychlé nedominované třídění} rozdělí populaci do front se složitostí
$O(kN^2)$.
\textbf{(2) Crowding distance} nahradí parametr $d_{\mathrm{share}}$, takže odpadá jeho
ruční nastavování.
\textbf{(3) Elitismus} spočívá v~tom, že se rodiče i~potomci spojí dohromady, setřídí
a~novou populaci tvoří lepší polovina.
\end{defbox}

\begin{defbox}[Crowding distance]
Počítá se pro každou frontu zvlášť. Pro každé kritérium $m$ se jedinci fronty seřadí
podle $f_m$, krajním jedincům se přiřadí $\infty$ a~ostatním se přičte normalizovaná
vzdálenost jejich sousedů:
\[
d_i \;=\; \sum_{m=1}^{k}
\frac{f_m(i+1) - f_m(i-1)}{f_m^{\max} - f_m^{\min}} .
\]
Výsledek vyjadřuje \uv{obvod kvádru} tvořeného nejbližšími sousedy, takže velká hodnota
znamená řídce obsazenou oblast, kterou chceme zachovat.
\end{defbox}

\begin{defbox}[Crowded comparison operator]
Jedinec $i$ je lepší než $j$ ve dvou případech: buď $i$ leží v~\emph{nižší} frontě
($rank_i < rank_j$), nebo se fronty shodují a~$i$ má \emph{větší} crowding distance.
Žádná skalární fitness se tedy vůbec nepočítá, porovnávají se jen tato dvě čísla,
typicky v~binárním turnaji.
\end{defbox}

\begin{algorithm}[!ht]
\caption{NSGA-II (jedna generace)}
\begin{algorithmic}[1]
\State $R_t \gets P_t \cup Q_t$ \Comment{rodiče a~potomci, celkem $2N$ jedinců}
\State $(F_1, F_2, \dots) \gets$ nedominované třídění $R_t$
\State $P_{t+1} \gets \emptyset$; $i \gets 1$
\While{$|P_{t+1}| + |F_i| \le N$}
  \State spočítej crowding distance ve $F_i$; $P_{t+1} \gets P_{t+1} \cup F_i$; $i \gets i+1$
\EndWhile
\State seřaď $F_i$ sestupně podle crowding distance a~doplň $P_{t+1}$ na $N$ jedinců
\State $Q_{t+1} \gets$ turnajová selekce (crowded comparison) $+$ křížení (SBX) $+$ mutace
\end{algorithmic}
\end{algorithm}

\paragraph{Číselný příklad crowding distance.} Vraťme se k~frontě
$A(1,10)$, $B(2,7)$, $C(4,5)$, $D(8,2)$. První kritérium se pohybuje v~rozsahu
$f_1 \in [1,8]$, tedy s~rozpětím 7, druhé v~rozsahu $f_2 \in [2,10]$ s~rozpětím 8.
Krajní body $A$ a~$D$ dostanou $\infty$, protože jsou krajní v~obou kritériích. Pro
vnitřní body vyjde
\[
d_B = \frac{f_1(C)-f_1(A)}{7} + \frac{f_2(A)-f_2(C)}{8}
    = \frac{4-1}{7} + \frac{10-5}{8} = 0{,}43 + 0{,}63 = 1{,}05,
\]
\[
d_C = \frac{f_1(D)-f_1(B)}{7} + \frac{f_2(B)-f_2(D)}{8}
    = \frac{8-2}{7} + \frac{7-2}{8} = 0{,}86 + 0{,}63 = 1{,}48 .
\]
Musí-li algoritmus vybrat jednoho z~dvojice $B$, $C$, zvítězí $C$, protože leží
v~řidčeji obsazené části fronty.

\paragraph{NSGA-III.} Tato varianta míří na úlohy s~mnoha kritérii ($k \ge 4$,
\emph{many-objective}), kde crowding distance selhává, protože při tolika kritériích je
nedominované téměř všechno. Crowding distance je zde nahrazena množinou rovnoměrně
rozmístěných \textbf{referenčních bodů}, jejichž počet zhruba odpovídá velikosti
populace a~pro $p$ dělení vychází $N = \binom{p+k-1}{p}$. Po nedominovaném třídění se
přednostně doplňují ty referenční směry, které mají nejméně přiřazených řešení;
nemá-li nějaký směr řešení žádné, přežije jedinec s~nejmenší kolmou vzdáleností od
příslušného referenčního vektoru.

\subsubsection{Další významné algoritmy}

\noindent\mimoq{slidy končí u~NSGA-II}

\textbf{SPEA2} si vedle populace udržuje externí \textbf{archiv} nedominovaných řešení.
Fitness jedince je součet \emph{sil} všech jedinců, kteří ho dominují, přičemž síla
jedince je počet jedinců, které sám dominuje; k~tomu se přičítá hustota odvozená ze
vzdálenosti k~$\sigma$-tému nejbližšímu sousedovi.
\textbf{MOEA/D} rozloží úlohu na $\mu$ \emph{skalárních} podúloh s~rovnoměrně
rozloženými váhovými vektory, které se na jedno číslo převádějí Čebyševovou skalarizací.
Každý jedinec odpovídá jedné podúloze a~spolupracuje jen se svým sousedstvím, odkud
plyne nízká složitost a~dobrá škálovatelnost.
\textbf{SMS-EMOA} vybírá přímo podle příspěvku jedince k~\emph{hypervolume}; teoreticky
je nejlépe podložený, ale s~počtem kritérií exponenciálně zdražuje.
\textbf{NSGA-III} už padl výše: nahrazuje crowding distance rovnoměrně rozmístěnými
\textbf{referenčními body} a~cílí na úlohy s~mnoha kritérii ($k \ge 4$), kde crowding
distance selhává.

\subsubsection{Metriky kvality aproximace}

\noindent\mimo{} Když algoritmus vrací celou aproximaci fronty místo jediného čísla,
není samo sebou, jak dva běhy porovnat. K~tomu slouží následující indikátory.

\begin{defbox}[Hypervolume ($S$-metrika)]
Objem oblasti dominované nalezenou frontou a~ohraničené \emph{referenčním bodem} $r$,
který se volí tak, aby byl dominován všemi řešeními. Je to jediná běžná metrika, která
je \emph{striktně monotónní vůči dominanci} a~nevyžaduje znalost skutečné Pareto fronty.

\emph{Příklad:} mějme frontu $\{(2,7), (4,5), (8,2)\}$ a~$r = (10,10)$. Objem sečteme po
svislých pruzích: bod $(8,2)$ dá $2\cdot8 = 16$, bod $(4,5)$ dá $4\cdot5 = 20$ a~bod
$(2,7)$ dá $2\cdot3 = 6$, celkem tedy $HV = 42$. Přidání dominovaného bodu hodnotu
nezmění.
\end{defbox}

Používají se i~další indikátory. \emph{GD} je průměrná vzdálenost nalezených bodů od
skutečné fronty, a~měří tedy konvergenci; \emph{IGD} počítá vzdálenost opačným směrem
a~postihne kromě konvergence i~pokrytí fronty; \emph{spread} hodnotí rovnoměrnost
rozložení. GD i~IGD ovšem vyžadují znalost skutečné fronty.

\subsubsection{Souhrnné srovnání MOEA}

Následující tabulka staví vedle sebe to, čím se jednotlivé algoritmy liší v~selekci
a~v~udržování diverzity.

\begin{center}
\small
\begin{tabular}{@{}llll@{}}
\toprule
\textbf{algoritmus} & \textbf{princip selekce} & \textbf{diverzita} & \textbf{poznámka} \\
\midrule
VEGA & podpopulace podle $f_i$ & žádná & konverguje ke krajům \\
NSGA & fronty + fiktivní fitness & fitness sharing ($d_{\mathrm{share}}$) & bez elitismu, $O(kN^3)$ \\
NSGA-II & fronty + crowding & crowding distance & elitismus, $O(kN^2)$, standard \\
NSGA-III & fronty + referenční body & referenční směry & many-objective \\
SPEA2 & síla dominujících & $\sigma$-tý soused, truncation & externí archiv \\
MOEA/D & skalarizované podúlohy & rovnoměrné váhové vektory & sousedství, škáluje \\
SMS-EMOA & příspěvek k~hypervolume & implicitní & teoreticky nejlépe podložené, drahé \\
\bottomrule
\end{tabular}
\end{center}

\subsection{Shrnutí ke~zkoušce}

\begin{itemize}
\item V~\textbf{michiganském} přístupu je jedinec jedno pravidlo a~řešením je celá
  populace, takže se musí řešit rozdělení odměny mezi pravidla (credit assignment).
  V~\textbf{pittsburghském} přístupu je jedinec celá sada pravidel, fitness se proto
  spočítá přímočaře, zato jsou složité operátory (GIL).
\item Hollandův LCS má podmínky zapsané nad abecedou $\{0,1,*\}$, používá vnitřní zprávy
  a~receptory a~odměnu rozděluje mechanismem bucket brigade, v~němž pravidlo platí svému
  předchůdci za právo jednat.
\item ZCS se obejde bez vnitřních zpráv: z~množiny odpovídajících pravidel (match set
  $M$) se ruletou vybere akce a~tím vznikne action set $A$, posiluje se $A$ i~$A_{-1}$
  a~pravidla z~$M \setminus A$ platí daň. Chybí-li pro vstup vhodné pravidlo, zasáhne
  cover operátor.
\item XCS pracuje s~klasifikátorem $(c,a,p,\epsilon,F)$ a~jeho hlavní myšlenkou je
  \textbf{fitness podle přesnosti predikce} $\kappa = \alpha(\epsilon/\epsilon_0)^{-\nu}$;
  predikce se aktualizuje Q-learningem a~nikový GA běží v~action setu.
  Výsledkem je úplná a~minimální mapa vstup--výstup. Varianty jsou UCS pro učení
  s~učitelem a~XCSF pro aproximaci funkcí.
\item Je dobré umět vysvětlit, proč přesnost jako fitness řeší nedostatky
  \uv{strength-based} systémů.
\item Evolvovat se dají váhy, topologie, obojí zároveň, nebo hyperparametry učení.
  Výhodou EA je, že se obejdou bez gradientu, zvládnou rekurentní sítě i~RL a~dobře se
  paralelizují.
\item Problém konkurujících konvencí (permutační problém) znamená, že křížení sítí je
  bez zvláštních opatření destruktivní.
\item Architekturu lze kódovat přímo maticí, gramaticky (Kitano), nebo buněčně (Gruau,
  kde genotypem je GP program pro růst sítě). GNARL je EP nad grafy se strukturálními
  i~parametrickými mutacemi, jejichž míru řídí teplota.
\item NEAT reprezentuje síť seznamem hran opatřených \textbf{inovačními čísly}, díky
  nimž se kříží jen odpovídající si geny. \textbf{Speciace} podle podobnosti genomů se
  sdílenou fitness chrání nově vzniklé topologie a~\textbf{complexification} znamená, že
  se začíná od minimální sítě.
\item HyperNEAT používá substrát $S:\mathbb{R}^4 \to \mathbb{R}$, který generuje váhy
  ze souřadnic neuronů; ten je reprezentovaný pomocí CPPN, tedy DAG s~různými aktivačními
  funkcemi, a~evolvuje se algoritmem NEAT. Díky tomu
  zachycuje symetrie a~regularity.
\item U~NAS se rozlišuje prostor architektur (lineární, DAG buněk, super síť),
  vyhledávací strategie (EA, RL, bayesovská, gradientní) a~způsob odhadu výkonu, kde
  pomáhá sdílení vah (ENAS). Sem patří i~DeepNEAT a~CoDeepNEAT s~moduly a~blueprinty
  a~ES pro deep RL.
\item Batoh se kóduje bitovou mapou, operátory jsou triviální a~celý problém je
  v~omezení kapacity. Nabízejí se tři strategie: \textbf{penalizace}, u~níž je nutné
  zvolit váhu $\alpha$, \textbf{oprava} jedince (lamarckovsky, nebo baldwinovsky)
  a~\textbf{dekodér} typu \uv{přidej, pokud se vejde}, po němž je každý genotyp
  přípustný; čtvrtou možností je vícekriteriální formulace.
\item U~TSP je fitness triviální a~problém je v~kódování. Sousednostní ani ordinální
  reprezentace se nepoužívají, osvědčila se permutační s~operátory PMX, OX, CX a~ER,
  případně maticová. Inicializovat se dá nejbližším sousedem nebo vkládáním hran,
  mutovat inverzí, což odpovídá kroku 2-opt.
\item U~SAT je fitness počet splněných klauzulí, samotné $f$ ovšem nestačí, a~proto se
  přidává vážení klauzulí a~zjemňující funkce; reprezentace může být bitová, reálná,
  klauzulová i~cestová a~jako lokální heuristika slouží WalkSAT. U~reálné reprezentace
  ($x \mapsto (1-y)^2$, $\neg x \mapsto (1+y)^2$, minimalizace) pozor na to, že hodnoty
  jsou \emph{penalty}, takže disjunkce (klauzule) se vyhodnotí jako \textbf{součin}
  a~konjunkce (formule) jako \textbf{součet}.
\item Rozvrhování se kóduje přímo maticí, tvrdá omezení se řeší už v~operátorech
  a~měkká až ve fitness; u~job shopu stojí proti sobě permutace s~dekodérem a~přímé
  kódování.
\item Obecně platí, že u~kombinatorických úloh rozhoduje kódování a~operátory
  respektující strukturu úlohy a~že hybridizace s~lokální heuristikou (2-opt, WSAT) je
  téměř vždy výhodná.
\item Dominance $x \prec y$ znamená, že $x$ není v~žádném kritériu horší a~alespoň
  v~jednom je lepší; Pareto množina leží v~prostoru proměnných, Pareto fronta v~prostoru
  kritérií a~dvojím cílem algoritmu je konvergence a~diverzita.
\item Ze skalarizací je nejjednodušší lineární, ta ale nenajde nekonvexní části fronty;
  dále se používá $\varepsilon$-constraint a~Čebyševova skalarizace.
\item NSGA-II třídí populaci nedominovaným tříděním do front, uvnitř fronty počítá
  crowding distance $d_i = \sum_m (f_m(i+1)-f_m(i-1))/(f_m^{\max}-f_m^{\min})$
  s~$\infty$ pro krajní body a~jedince porovnává operátorem crowded comparison (vyhrává
  nižší fronta, při shodě větší crowding distance). Elitismus zajišťuje spojení
  $P_t \cup Q_t$. Crowding distance je třeba umět spočítat na příkladu.
\item SPEA2 si udržuje archiv a~fitness staví na síle dominujících jedinců a~na hustotě
  přes $\sigma$-tého souseda, MOEA/D používá Čebyševovu skalarizaci s~váhovými vektory
  a~sousedstvím, SMS-EMOA vybírá podle příspěvku k~hypervolume a~NSGA-III nasazuje
  referenční body na many-objective úlohy.
\item Hypervolume je objem dominovaný frontou vůči referenčnímu bodu a~jako jediná běžná
  metrika je monotónní vůči dominanci a~obejde se bez znalosti skutečné fronty. Vedle něj
  se používají GD, IGD, spread a~$\varepsilon$-indikátor.
\end{itemize}

%=====================================================================
\section*{Souhrnné srovnání a~typické zkouškové otázky}
\addcontentsline{toc}{section}{Souhrnné srovnání a~typické zkouškové otázky}
\label{sec:zaver}
%=====================================================================
\subsection*{Jeden pohled na všechny algoritmy}

Když jsou všechny metody probrané, vyplatí se podívat se na ně jedním společným pohledem. Slidy
EVA~II tuto myšlenku formulují až u~CMA-ES, platí ale pro celý okruh: probírané algoritmy dělají
v~jádru jednu a~tutéž věc. U~zkoušky je to nejlepší možné zahájení kterékoli odpovědi, protože
rovnou ukáže, že jednotlivé metody nejsou nesouvislá sbírka triků.

\begin{defbox}[Stochastické black-box hledání]
\textbf{Black-box optimalizace} znamená, že minimalizujeme funkci
$f : \mathbb{R}^n \to \mathbb{R}$, přičemž jedinou dostupnou informací o~$f$ jsou \emph{její
hodnoty ve zvolených bodech}. Gradient k~dispozici není a~analytický popis také ne. Cenou
hledání je proto \emph{počet vyhodnocení}~$f$.

Všechny probírané metody mají pak tentýž skelet. Algoritmus si udržuje \textbf{rozdělení
pravděpodobnosti} $p_\theta(x)$ nad prostorem řešení a~pořád dokola opakuje tři kroky:
\[
\text{navzorkuj } \{x_i\}_{i=1}^n \sim p_\theta
\;\longrightarrow\;
\text{vyhodnoť } \{(x_i, f(x_i))\}
\;\longrightarrow\;
\theta \leftarrow h(\theta, D).
\]
Parametr $\theta$ je \emph{jediná} informace, která se přenáší mezi iteracemi. Kóduje všechno,
co se algoritmus dosud naučil, a~určuje, kde bude hledat dál. Jednotlivé metody se pak liší jen
ve dvou věcech: \emph{co je} $\theta$ a~jak zní aktualizační heuristika~$h$.
\end{defbox}

\begin{center}
\small
\begin{tabular}{@{}p{3.2cm}p{11.6cm}@{}}
\toprule
\textbf{algoritmus} & \textbf{co je $\theta$} \\
\midrule
EA (GA, GP, \dots) & rodičovská populace --- rozdělení je empirické, dané vzorkem;
  aktualizace $=$ selekce $+$ variace \\
ES & gaussián: střed a~šířka $\vec\sigma$ (u~samoadaptace nese $\sigma$ každý jedinec sám) \\
CMA-ES & $(\vec m, \sigma, \mathbf{C}, \vec p_\sigma, \vec p_c)$ --- plná kovariance
  $+$ evoluční cesty \\
EDA & parametry explicitního modelu: marginály (UMDA, PBIL), párové závislosti (MIMIC),
  bayesovská síť (BOA) \\
PSO & pozice a~rychlosti částic $+$ paměť $\vec p_i$, $\vec g$ \\
ACO & feromonová matice $\tau$ --- rozdělení nad \emph{konstrukcemi}, ne nad body \\
SA & aktuální bod $+$ teplota $T$ \\
tabu search & aktuální bod $+$ seznam zakázaných tahů \\
\bottomrule
\end{tabular}
\end{center}

Z~téhle perspektivy přestávají být některé metody překvapením. EDA (\S\ref{ssec:eda}) je jen
důsledné domyšlení celé myšlenky: když je populace stejně jen implicitní model rozdělení, není
důvod ho nemodelovat rovnou explicitně. A~naopak, poznání, že genetický algoritmus není nic
jiného než transformace rozdělení, je přesně to, co formalizuje Voseho model
(\S\ref{ssec:vose}).

\subsection*{Mapa oboru}

Celý okruh se dá rozdělit do čtyř skupin podle toho, odkud metoda bere svou inspiraci a~co je
jejím předmětem.

\begin{itemize}
\item \textbf{Evoluční algoritmy} stojí na populaci, fitness, křížení, mutaci a~selekci. Patří
  sem genetické algoritmy ve všech svých kódováních (binární, celočíselné, permutační,
  reálné), evoluční strategie, diferenciální evoluce, evoluční programování, genetické
  programování se svými variantami (stromové, lineární, kartézské, gramatická evoluce),
  learning classifier systems a~algoritmy odhadu rozdělení (EDA).
\item \textbf{Přírodou inspirované metody bez evoluce} čerpají předlohu z~přírody, evoluční
  schéma ale nepoužívají. Jsou to rojové algoritmy PSO, ACO, ABC a~firefly.
\item \textbf{Nebiologicky inspirované metaheuristiky} berou inspiraci odjinud než z~živé
  přírody. Řadí se mezi ně tabu search, scatter search, novelty search, simulované žíhání
  a~memetické algoritmy.
\item \textbf{Aplikační oblasti} si vynucují vlastní reprezentace a~operátory, a~proto se
  probírají samostatně. Patří mezi ně kombinatorická optimalizace, neuroevoluce, zpětnovazební
  učení, vícekriteriální optimalizace a~AutoML/NAS.
\end{itemize}

\subsection*{Souhrnná srovnávací tabulka algoritmů}

Následující tabulka staví algoritmy vedle sebe podle čtyř hledisek: čím reprezentují řešení,
jak vytvářejí nové kandidáty, jak vybírají mezi nimi a~co se u~každého z~nich nejčastěji chce
slyšet.

\begin{center}
\small
\begin{tabular}{@{}p{2.5cm}p{2.9cm}p{3.1cm}p{3.1cm}p{3.1cm}@{}}
\toprule
\textbf{algoritmus} & \textbf{reprezentace} & \textbf{variace} &
\textbf{selekce} & \textbf{co si zapamatovat} \\
\midrule
SGA & binární řetězec & 1-bod.\ křížení, bit-flip & ruleta, generační &
  teorie schémat, BBH \\
ES & $\mathbb{R}^n$ + $\sigma$ & gaussovská mutace, rekombinace &
  náhodní rodiče, $(\mu,\lambda)$/$(\mu{+}\lambda)$ & pravidlo 1/5,
  samoadaptace \\
CMA-ES & $\mathbb{R}^n$, model $\mathcal N(m,\sigma^2 C)$ & vzorkování
  z~rozdělení & pořadová, $\mu$ nejlepších & evoluční cesty, rank-1/rank-$\mu$ \\
DE & $\mathbb{R}^n$ & $\vec x_a + F(\vec x_b - \vec x_c)$, binom.\ křížení &
  souboj rodič--potomek & měřítko z~populace \\
EP & doménová (automat) & jen mutace & turnaj rodiče+potomci &
  genotyp = fenotyp \\
GP & syntaktický strom & výměna podstromů, mutace & turnaj &
  bloat, ADF \\
PSO & $\mathbb{R}^n$ + rychlost & pohybové rovnice & žádná (jen paměť) &
  $w$, $c_1$, $c_2$, gbest/lbest \\
ACO & konstrukce cesty & feromon + heuristika & implicitní &
  $\tau^\alpha\eta^\beta$, odpařování \\
SA & jedno řešení & soused & Metropolis & rozvrh chlazení \\
MA & libovolná & EA operátory + lokální hledání & libovolná &
  Lamarck vs.\ Baldwin \\
NSGA-II & libovolná & SBX + polynom.\ mutace & fronty + crowding &
  elitismus, crowding distance \\
NEAT & seznam hran & mutace struktury, křížení dle ID & sdílená fitness
  v~druhu & inovační čísla, speciace \\
\bottomrule
\end{tabular}
\end{center}

\subsection*{Průřezová témata --- co se ptá napříč otázkami}

Některá témata nepatří k~jedné metodě, ale prostupují celým okruhem. Právě na nich se pozná,
jestli student vidí souvislosti, a~proto se na ně zkoušející ptají nejraději.

\begin{itemize}
\item \textbf{Explorace vs.\ exploatace}: u~každého algoritmu se dá ukázat, čím se rovnováha
  mezi nimi řídí. U~GA je to napětí mezi křížením a~selekčním tlakem, u~ES velikost~$\sigma$
  a~pravidlo 1/5, u~simulovaného žíhání teplota, u~PSO váha~$w$ a~topologie, u~ACO parametr
  $q_0$ a~odpařování, u~tabu search samotný tabu list. Novelty search je krajní případ, protože
  stojí jen na exploraci.
\item \textbf{Udržení diverzity} se řeší celou sadou nástrojů: fitness sharing, crowding,
  speciace, ostrovní model, nenáhodné páření, restart a~hypermutace.
\item \textbf{Adaptace parametrů} se probírá jako dvojice ladění versus řízení a~jako trojice
  pevná, adaptivní a~samoadaptivní strategie. U~zkoušky se pak chce vědět, kde se která z~nich
  v~kterém algoritmu vyskytuje.
\item \textbf{Práce s~omezeními} má několik standardních řešení: penalizaci, opravu, dekodér,
  uzavřené operátory a~vícekriteriální formulaci.
\item \textbf{Lamarckismus vs.\ baldwinismus} se objevuje u~memetických algoritmů, při opravě
  jedinců i~v~neuroevoluci s~douč\-ováním vah.
\item \textbf{Teoretické zázemí} okruhu tvoří věta o~schématech i~její kritika, Voseho model,
  Markovovy řetězce a~konvergence elitistického GA, No Free Lunch a~konvergence SA.
\end{itemize}

\subsection*{Typické zkouškové otázky}

Následující otázky se u~zkoušky objevují opakovaně; jsou to doslovná zadání, ne parafráze.

\begin{enumerate}
\item Popište obecné schéma evolučního algoritmu a~vyjmenujte typy selekce
  včetně jejich vlastností. Spočtěte ruletovou selekci pro danou populaci.
\item Formulujte a~dokažte větu o~schématech. Co je řád a~definující délka?
  Jaké jsou slabiny věty?
\item Vysvětlete hypotézu stavebních bloků a~implicitní paralelismus.
  Co je deceptivní úloha?
\item Popište Voseho model SGA: co jsou vektory $p$ a~$s$, co dělá operátor
  $\mathcal{F}$, co $\mathcal{M}$, jaké mají pevné body? Jak vypadá model
  konečné populace jako Markovův řetězec?
\item Co jsou algoritmy odhadu rozdělení (EDA)? Jaký je jejich vztah
  k~Voseho pohledu na GA a~jak se liší UMDA/PBIL od BOA?
\item Jaký je rozdíl mezi $(\mu,\lambda)$ a~$(\mu+\lambda)$-ES?
  Vysvětlete pravidlo 1/5 a~samoadaptaci $\sigma$. Jak funguje CMA-ES?
\item Popište jeden krok diferenciální evoluce včetně vzorců
  a~číselného příkladu. Jaké jsou varianty mutace?
\item Jak funguje stromové genetické programování? Co je bloat a~jak se proti
  němu bojuje? K~čemu slouží ADF?
\item Vysvětlete gramatickou evoluci a~kartézské GP --- jaký je genotyp
  a~fenotyp?
\item Co je koevoluce, jaké má typy a~jaké patologie? Jak se hodnotí fitness?
\item Co je otevřená evoluce, jaké má znaky a~proč se umělé systémy
  \uv{zasekávají}? Jak funguje novelty search a~proč se u~deceptivních úloh
  vyplatí zahodit cíl?
\item Vysvětlete iterované vězňovo dilema, strategii TFT a~vlastnosti
  úspěšných strategií.
\item Napište rovnice PSO a~vysvětlete význam jednotlivých členů.
  Jaký je rozdíl mezi gbest a~lbest topologií?
\item Popište ACO: pravidlo přechodu, aktualizaci feromonu, rozdíl mezi
  AS a~ACS.
\item Vysvětlete simulované žíhání, Metropolisovo kritérium a~rozvrh chlazení.
  Kdy konverguje?
\item Co je memetický algoritmus? Vysvětlete rozdíl mezi lamarckovským
  a~baldwinovským učením.
\item Porovnejte michiganský a~pittsburghský přístup. Jak funguje XCS a~proč
  je fitness založena na přesnosti?
\item Popište NEAT: k~čemu slouží inovační čísla a~speciace?
  Co přidává HyperNEAT?
\item Jak se řeší TSP evolučním algoritmem? Popište PMX, OX, CX a~ER.
\item Jak se v~EA pracuje s~omezeními (na příkladu problému batohu)?
\item Definujte dominanci a~Pareto frontu. Popište NSGA-II a~spočtěte
  crowding distance na příkladu. Co je hypervolume?
\end{enumerate}

\subsection*{Závěrečné shrnutí}

\begin{itemize}
\item Evoluční a~rojové algoritmy jsou \textbf{robustní metaheuristiky} pro úlohy, na kterých
  selhávají exaktní a~gradientní metody. Typicky jde o~úlohy black-box,
  nedi\-fe\-ren\-co\-va\-tel\-né, diskrétní, struk\-tu\-ro\-va\-né,
  ví\-ce\-kri\-te\-ri\-ál\-ní a~dynamické.
\item Jejich síla neleží v~jednom \uv{nejlepším} algoritmu, protože takový podle No Free Lunch
  neexistuje. Leží ve \textbf{schopnosti přizpůsobit reprezentaci a~operátory dané doméně}
  a~v~možnosti \textbf{hybridizace} s~lokálními heuristikami a~doménovými znalostmi.
\item Teorie, tedy schémata, Voseho model, konvergenční věty a~analýza doby běhu, poskytuje
  intuici a~a\-sym\-pto\-tic\-ké záruky. Praktický návrh přesto zůstává z~velké části
  ex\-pe\-ri\-men\-tál\-ní dis\-cip\-lí\-nou.
\end{itemize}

\end{document}
